% Generated mechanically from the frozen artin.tex target.
% Do not edit this generated file; edit the bound locale source instead.

% OK, start here.
%

\setcounter{chapter}{97}
\chapter{Artin 公理}


\phantomsection
\label{artin-section-phantom}





\section{引言}
\label{artin-section-introduction}

\noindent
本章讨论函子可由代数空间表示的 Artin 公理。推荐参考论文
\cite{ArtinI}, \cite{ArtinII}, \cite{ArtinVersal}.

\medskip\noindent
% P11-E0200: frozen source has plural coordinated subjects with singular "is awkward"; translated naturally.
本章中的某些记号、约定和术语颇为别扭，经验更丰富的读者甚至可能觉得
它们前后颠倒。这是有意为之。解释见《Quot 空间与 Hilbert 空间》第
\ref{quot-section-conventions} 节。

\medskip\noindent
设 $S$ 是局部 Noether 基概形，并设
$$
p : \mathcal{X} \longrightarrow (\Sch/S)_{fppf}
$$
是群胚纤维化范畴。设 $x_0$ 是 $\mathcal{X}$ 在域 $k$ 上的一个对象，
其中该域在 $S$ 上有限型。贯穿本章，与 $k$ 上的 $x_0$ 相伴的预形变范畴
（见《形式形变理论》定义
\ref{formal-defos-definition-predeformation-category}）
$$
\mathcal{F}_{\mathcal{X}, k, x_0}
\longrightarrow
\{\text{Artinian local }S\text{-algebras with residue field }k\}
$$
起着重要作用。我们为 $\mathcal{X}$ 引入 Rim--Schlessinger 条件 (RS)，
并证明该条件保证 $\mathcal{F}_{\mathcal{X}, k, x_0}$ 是形变范畴；也就是说，
$\mathcal{F}_{\mathcal{X}, k, x_0}$ 本身满足 (RS)。我们讨论用有限扩张
替换 $k$ 时 $\mathcal{F}_{\mathcal{X}, k, x_0}$ 如何变化，并讨论切空间。

\medskip\noindent
接着，我们讨论 $\mathcal{X}$ 的形式对象 $\xi = (\xi_n)$；它们是位于商
$R/\mathfrak m^n$ 上的对象所成的逆系统，其中 $R$ 是 Noether 完备局部
$S$-代数，其剩余域在 $S$ 上有限型。这等价于对某个 $x_0$ 和 $k$，
在 $\mathcal{F}_{\mathcal{X}, k, x_0}$ 中有一个形式对象。若 $\mathcal{X}$
在 $R$ 上有一个对象可产生该逆系统，就称形式对象为有效的。若
$\mathcal{X}$ 的形式对象产生 $\mathcal{F}_{\mathcal{X}, k, x_0}$ 的通用
形式对象，就称它是通用的。最后，给定有限型 $S$-概形 $U$、
$\mathcal{X}$ 在 $U$ 上的对象 $x$ 和闭点 $u_0 \in U$，若完备局部环
$\mathcal{O}_{U, u_0}^\wedge$ 上诱导的形式对象是通用的，就称 $x$
在 $u_0$ 处通用。

\medskip\noindent
完成这些准备后，可以陈述 Artin 的著名定理：若下列条件成立，则
$\mathcal{X}$ 是代数叠：
\begin{enumerate}
\item 对所有 $s \in S$，$\mathcal{O}_{S, s}$ 都是 G-环\footnote{这意味着
$S$ 是 G-概形，参见《性质》第 \ref{properties-section-G-scheme} 节。}；
\item $\Delta : \mathcal{X} \to \mathcal{X} \times \mathcal{X}$
可由代数空间表示；
\item $\mathcal{X}$ 是 \'etale 拓扑的叠；
\item $\mathcal{X}$ 保极限；
\item $\mathcal{X}$ 满足 (RS)；
\item 形变范畴 $\mathcal{F}_{\mathcal{X}, k, x_0}$ 的切空间和无穷小
自同构空间是有限维的；
\item 形式对象有效；
\item $\mathcal{X}$ 满足通用性的开性。
\end{enumerate}
这就是引理 \ref{artin-lemma-diagonal-representable}；一个略有改进的版本另见
命题 \ref{artin-proposition-second-diagonal-representable}。还有一个类似命题刻画
哪些函子 $F : (\Sch/S)_{fppf}^{opp} \to \textit{Sets}$ 是代数空间，
参见第 \ref{artin-section-algebraic-spaces} 节。

\medskip\noindent
下面粗略概述 Artin 定理的证明。首先，利用 (RS) 以及切空间和自同构空间
的有限维性，证明存在充分多的通用形式对象；例如参见《形式形变理论》引理
\ref{formal-defos-lemma-minimal-groupoid-in-functors-construction}。按假设，
这些形式对象有效。有效形式对象可由有限型 $S$-概形 $U$ 上的对象 $x$
“逼近”；参见引理 \ref{artin-lemma-approximate}。
% P11-E0201: frozen source omits "that" before "the local rings of S are G-rings"; translated as the intended coordination.
这一逼近利用 $S$ 的局部环都是 G-环以及 $\mathcal{X}$ 保极限这两个事实；
这也许是证明中最困难的部分，因为它依赖一般 N\'eron 去奇异化，用
Hensel 化中的解去逼近 Noether 局部 G-环上代数方程的形式解。接着，
通用性的开性蕴含：缩小 $U$ 后，可以假设 $x$ 在 $U$ 的每个闭点处通用。
完成这一切后，我们证明 $U \to \mathcal{X}$ 是光滑态射。取充分多的
$U \to \mathcal{X}$，即可得到 $\mathcal{X}$ 的一个“光滑图册”，从而证明
$\mathcal{X}$ 是代数叠。

\medskip\noindent
对给定的群胚纤维化范畴 $\mathcal{X}$ 检验 Artin 公理时，最困难的一步
通常是验证通用性的开性。以下讨论假设 $\mathcal{X}/S$ 已满足上列其他条件。
本章给出两种方法，用以证明 $\mathcal{X}$ 满足通用性的开性：
\begin{enumerate}
\item 第一种方法是假设一个称为 (RS*) 的更强 Rim--Schlessinger 条件，
并假设更强版本的形式有效性；本质上，它要求加厚逆系统上的对象有效。
事实证明，在这些假设下，通用性的开性自动成立；参见引理
\ref{artin-lemma-SGE-implies-openness-versality}。
% P11-E0202: frozen source says "that category of all schemes" where the definite article "the" is intended.
请注意，这里至关重要地用到了 $\mathcal{X}$ 定义在 $S$ 上所有概形的范畴
之上，而不只定义在 Noether 概形的范畴之上！
\item 第二种方法依循 Artin，要求 $\mathcal{X}$ 配备障碍理论。若该障碍理论
在适当意义下“与积可交换”，则 $\mathcal{X}$ 满足通用性的开性；参见引理
\ref{artin-lemma-get-openness-obstruction-theory}。
\end{enumerate}
障碍理论可以用许多不同方式公理化；文献中确实可见许多变体，往往针对特定
模叠作了调整。引理 \ref{artin-lemma-dual-openness} 说明一个使用导出范畴的变体；
它往往自然产生于文献中的形变理论计算。

\medskip\noindent
在第 \ref{artin-section-algebraic-spaces-noetherian} 节中，我们讨论需要作何修改，
才能把结果用于定义在 $S$ 上局部 Noether 概形范畴
$(\textit{Noetherian}/S)_\etale$ 上的函子。

\medskip\noindent
在本章最后一节，作为 Artin 公理的应用，我们证明 Artin 关于收缩存在性的
定理；参见第 \ref{artin-section-contractions} 节。该定理粗略地说：给定
% P11-E0203: frozen source says "X' separated of finite type"; translated with the intended conjunction.
在 $S$ 上分离且有限型的代数空间 $X'$、闭子集 $T' \subset |X'|$，以及形式修改
$$
\mathfrak{f} : X'_{/T'} \longrightarrow \mathfrak{X}
$$
其中 $\mathfrak{X}$ 是 $S$ 上的 Noether 形式代数空间，则存在一个真态射
$f : X' \to X$，“实现该收缩”。这指的是：存在等同
$\mathfrak{X} = X_{/T}$，使得
$\mathfrak{f} = f_{/T'} : X'_{/T'} \to X_{/T}$，其中 $T = f(T')$；
而且 $f$ 在 $X \setminus T$ 上是同构。证明的做法是在 $S$ 上局部 Noether
概形的范畴上定义函子 $F$，并证明 $F$ 满足 Artin 公理。有趣的是，在
Artin 公理的这个应用中，通用性的开性并非最难证明之处；
% P11-E0204: frozen source joins two independent clauses with a comma before "instead"; translated as two coordinated clauses.
反而是证明 $F$ 保极限需要大量工作和预备结果。




\section{约定}
\label{artin-section-conventions}

\noindent
本章采用的约定与《代数叠》一章中的约定相同；参见《代数叠》第
\ref{algebraic-section-conventions} 节。本章中的基概形 $S$ 往往是局部
Noether 的（不过在陈述结果时，我们总会重申这一条件）。





\section{预形变范畴}
\label{artin-section-predeformation-categories}

\noindent
设 $S$ 是局部 Noether 基概形，并设
$$
p : \mathcal{X} \longrightarrow (\Sch/S)_{fppf}
$$
是群胚纤维化范畴。设 $k$ 是域，而 $\Spec(k) \to S$ 是有限型态射
（参见《态射》引理 \ref{morphisms-lemma-point-finite-type}）。有时，
我们简称 {\it $k$ 是 $S$ 上有限型的域}。设 $x_0$ 是 $\mathcal{X}$
中位于 $\Spec(k)$ 上的对象。给定 $S$、$\mathcal{X}$、$k$ 和 $x_0$，
我们将构造《形式形变理论》定义
\ref{formal-defos-definition-predeformation-category} 意义下的预形变范畴。
该构造类似于《形式形变理论》注记
\ref{formal-defos-remark-localize-cofibered-groupoid} 的构造。

\medskip\noindent
首先，由《态射》引理 \ref{morphisms-lemma-point-finite-type}，可以选取仿射开集
$\Spec(\Lambda) \subset S$，使 $\Spec(k) \to S$ 经 $\Spec(\Lambda)$ 分解，
且相伴的环同态 $\Lambda \to k$ 有限。这给出范畴 $\mathcal{C}_\Lambda$；
参见《形式形变理论》定义 \ref{formal-defos-definition-CLambda}。在典范等价
意义下，范畴 $\mathcal{C}_\Lambda$ 不依赖于 $S$ 的仿射开集
$\Spec(\Lambda)$ 的选择。确切地说，$\mathcal{C}_\Lambda$ 等价于下述分解
所成范畴的反范畴：
\begin{equation}
\label{artin-equation-factor}
\Spec(k) \to \Spec(A) \to S
\end{equation}
这里要求 $A$ 是 Artin 局部环，并且 $\Spec(k) \to \Spec(A)$ 对应于
环同态 $A \to k$，它把 $k$ 等同于 $A$ 的剩余域。

\medskip\noindent
令 $\mathcal{F} = \mathcal{F}_{\mathcal{X}, k, x_0}$ 为如下范畴：
\begin{enumerate}
\item 对象是 $\mathcal{X}$ 的态射 $x_0 \to x$，其中
$p(x) = \Spec(A)$，$A$ 是 Artin 局部环，而
$p(x_0) \to p(x) \to S$ 是形如 (\ref{artin-equation-factor}) 的分解；并且
\item 态射 $(x_0 \to x) \to (x_0 \to x')$ 是 $\mathcal{X}$ 中的交换图
$$
\xymatrix{
x & & x' \ar[ll] \\
& x_0 \ar[lu] \ar[ru]
}
$$
。（注意箭头方向的反转。）
\end{enumerate}
若 $x_0 \to x$ 是 $\mathcal{F}$ 的对象，写 $p(x) = \Spec(A)$，便得到
$\mathcal{C}_\Lambda$ 的对象 $A$。我们常说 $x_0 \to x$ 或 $x$ 位于
$A$ 上。$\mathcal{F}$ 中从位于 $A$ 上的对象 $x_0 \to x$ 到位于 $A'$
上的对象 $x_0 \to x'$ 的态射，对应于 $\mathcal{X}$ 的态射
$x' \to x$，从而对应于态射
$p(x' \to x) : \Spec(A') \to \Spec(A)$；后者又对应于环同态
$A \to A'$。由于 $\mathcal{X}$ 是 $S$ 上概形范畴之上的范畴，
% P11-E0205: frozen source omits the article "a" before Lambda-algebra homomorphism; translated naturally.
可知 $A \to A'$ 是 $\Lambda$-代数同态。于是得到函子
\begin{equation}
\label{artin-equation-predeformation-category}
p : \mathcal{F} = \mathcal{F}_{\mathcal{X}, k, x_0}
\longrightarrow
\mathcal{C}_\Lambda.
\end{equation}
我们用记号 $\mathcal{F}(A)$ 表示 $\mathcal{C}_\Lambda$ 的对象 $A$ 上的
纤维范畴。$\mathcal{F}(A)$ 的对象就是 $\mathcal{X}$ 的态射
$x_0 \to x$，使得 $x$ 位于 $\Spec(A)$ 上，而 $x_0 \to x$ 位于
$\Spec(k) \to \Spec(A)$ 上。

\begin{lemma}
\label{artin-lemma-predeformation-category}
上面定义的函子 $p : \mathcal{F} \to \mathcal{C}_\Lambda$ 是预形变范畴。
\end{lemma}

\begin{proof}
须证明 (a) $\mathcal{F}$ 在 $\mathcal{C}_\Lambda$ 上以群胚为纤维余纤维化；
并且 (b) $\mathcal{F}(k)$ 等价于仅有一个对象和一个态射的范畴。

\medskip\noindent
(a) 的证明。由于 $\mathcal{X}$ 的纤维范畴是群胚，$\mathcal{F}$ 在
$\mathcal{C}_\Lambda$ 上的纤维范畴也是群胚。设 $A \to A'$ 是
$\mathcal{C}_\Lambda$ 的态射，而 $x_0 \to x$ 是 $\mathcal{F}(A)$
的对象。由于 $\mathcal{X}$ 以群胚为纤维纤维化，可以找到位于
$\Spec(A') \to \Spec(A)$ 上的态射 $x' \to x$。
% P11-E0206: frozen source says "is equal the given map"; "to" is missing. Translated naturally.
由于复合 $A \to A' \to k$ 等于给定映射 $A \to k$，由拉回在同构意义下
的唯一性可知，$x'$ 经 $\Spec(k) \to \Spec(A')$ 的拉回是 $x_0$；
也就是说，存在位于 $\Spec(k) \to \Spec(A')$ 上的态射 $x_0 \to x'$，
它与 $x_0 \to x$ 和 $x' \to x$ 相容。这证明 $\mathcal{F}$ 具有推前。
结论由《范畴》引理 \ref{categories-lemma-fibred-groupoids} 的对偶形式给出。

\medskip\noindent
(b) 的证明。若 $A = k$，则 $\Spec(k) = \Spec(A)$。由于 $\mathcal{X}$
在 $(\Sch/S)_{fppf}$ 上以群胚为纤维纤维化，对 $\mathcal{F}(k)$
中的任意对象 $x_0 \to x$，态射 $x_0 \to x$ 都是同构。因此，
$\mathcal{F}(k)$ 的每个对象都同构于 $x_0 \to x_0$。显然，
$\mathcal{F}$ 中 $x_0 \to x_0$ 的唯一自态射是恒等态射。
\end{proof}

\noindent
设 $S$ 是局部 Noether 基概形，$F : \mathcal{X} \to \mathcal{Y}$ 是
$(\Sch/S)_{fppf}$ 上群胚纤维化范畴之间的 $1$-态射。
% P11-E0207: frozen source says "Let k is a field" instead of "Let k be a field"; translated naturally.
设 $k$ 是 $S$ 上有限型的域。设 $x_0$ 是 $\mathcal{X}$ 中位于
$\Spec(k)$ 上的对象。令 $y_0 = F(x_0)$；这是 $\mathcal{Y}$ 中位于
$\Spec(k)$ 上的对象。则 $F$ 诱导函子
\begin{equation}
\label{artin-equation-functoriality}
F :
\mathcal{F}_{\mathcal{X}, k, x_0}
\longrightarrow
\mathcal{F}_{\mathcal{Y}, k, y_0}
\end{equation}
，其源和靶都是 $\mathcal{C}_\Lambda$ 上的余纤维化范畴。确切地说，
它把 $\mathcal{F}_{\mathcal{X}, k, x_0}(A)$ 的对象 $x_0 \to x$
映为 $\mathcal{F}_{\mathcal{Y}, k, y_0}(A)$ 的对象 $F(x_0) \to F(x)$。

\begin{lemma}
\label{artin-lemma-formally-smooth-on-deformation-categories}
设 $S$ 是局部 Noether 概形，$F : \mathcal{X} \to \mathcal{Y}$ 是
$(\Sch/S)_{fppf}$ 上群胚纤维化范畴的一个 $1$-态射。假设下列条件之一成立：
\begin{enumerate}
\item $F$ 在对象上形式光滑（《可表示性判据》第
\ref{criteria-section-formally-smooth} 节）；
\item $F$ 可由代数空间表示且形式光滑；或
\item $F$ 可由代数空间表示且光滑。
\end{enumerate}
则对每个 $S$ 上有限型的域 $k$ 以及 $\mathcal{X}$ 在 $k$ 上的对象
$x_0$，函子 (\ref{artin-equation-functoriality}) 在《形式形变理论》定义
\ref{formal-defos-definition-smooth-morphism} 的意义下光滑。
\end{lemma}

\begin{proof}
情形 (1) 只需展开定义。由《可表示性判据》引理
\ref{criteria-lemma-representable-by-spaces-formally-smooth}，假设 (2)
蕴含 (1)。由《空间的态射进阶》引理
\ref{spaces-more-morphisms-lemma-smooth-formally-smooth} 以及《代数叠》引理
\ref{algebraic-lemma-representable-transformations-property-implication}
的一般原理，假设 (3) 蕴含 (2)。
\end{proof}

\begin{lemma}
\label{artin-lemma-fibre-product-deformation-categories}
设 $S$ 是局部 Noether 概形。设
$$
\xymatrix{
\mathcal{W} \ar[d] \ar[r] & \mathcal{Z} \ar[d] \\
\mathcal{X} \ar[r] & \mathcal{Y}
}
$$
是 $(\Sch/S)_{fppf}$ 上群胚纤维化范畴的 $2$-纤维积。设 $k$ 是
$S$ 上有限型的域，而 $w_0$ 是 $\mathcal{W}$ 在 $k$ 上的对象。
设 $x_0, z_0, y_0$ 是 $w_0$ 在图中各态射下的像。则
$$
\xymatrix{
\mathcal{F}_{\mathcal{W}, k, w_0} \ar[d] \ar[r] &
\mathcal{F}_{\mathcal{Z}, k, z_0} \ar[d] \\
\mathcal{F}_{\mathcal{X}, k, x_0} \ar[r] & \mathcal{F}_{\mathcal{Y}, k, y_0}
}
$$
是预形变范畴的纤维积。
\end{lemma}

\begin{proof}
只需展开定义。细节从略。
\end{proof}






\section{推出与叠}
\label{artin-section-pushouts}

\noindent
本节证明代数叠对于某些推出具有良好行为。本节结果对任意基概形都成立。

\medskip\noindent
% P11-E0208: the frozen source contains an unresolved editorial-reference placeholder.
当 $Y$、$X'$、$X$、$Y'$ 是代数空间时，下述引理也成立；
参见（此处补入将来的参考文献）。

\begin{lemma}
\label{artin-lemma-pushout}
\begin{slogan}
代数叠满足（强）Rim--Schlessinger 条件
\end{slogan}
设 $S$ 是概形。设
$$
\xymatrix{
X \ar[r] \ar[d] & X' \ar[d] \\
Y \ar[r] & Y'
}
$$
是 $S$ 上概形范畴中的推出，其中 $X \to X'$ 是加厚，且 $X \to Y$
是仿射态射；参见《更多态射》引理
\ref{more-morphisms-lemma-pushout-along-thickening}。
设 $\mathcal{Z}$ 是 $S$ 上的代数叠。则纤维范畴的函子
$$
\mathcal{Z}_{Y'}
\longrightarrow
\mathcal{Z}_Y \times_{\mathcal{Z}_X} \mathcal{Z}_{X'}
$$
是范畴等价。
\end{lemma}

\begin{proof}
% P11-E0209: the frozen full-faithfulness argument names one object and only its self-Isom sheaf.
设 $y'$ 是左端的一个对象。定义在 $Y'$ 上概形范畴上的层
$\mathit{Isom}(y', y')$ 可由 $Y'$ 上的代数空间 $I$ 表示；参见
《代数叠》引理 \ref{algebraic-lemma-representable-diagonal}。
由于 $Y'$ 不仅在概形范畴中，而且在代数空间范畴中也是该推出，
我们断定引理中的函子是全忠实的；参见《空间的推出》引理
\ref{spaces-pushouts-lemma-pushout-along-thickening-schemes}。

\medskip\noindent
设 $(y, x', f)$ 是右端的一个对象，其中
$f : y|_X \to x'|_X$ 是同构。为完成证明，我们必须构造
$\mathcal{Z}_{Y'}$ 的一个对象 $y'$，使其在 $Y$ 和 $X'$ 上的限制
分别与 $y$ 和 $x'$ 相符，并且这种相符方式与 $f$ 相容。事实上，
只需在 $Y'$ 上 fppf 局部地构造 $y'$；参见《叠》引理
\ref{stacks-lemma-characterize-essentially-surjective-when-ff}。
选取一个可表代数叠 $\mathcal{W}$ 和一个满光滑态射
$\mathcal{W} \to \mathcal{Z}$。于是
$$
(\Sch/Y)_{fppf} \times_{y, \mathcal{Z}} \mathcal{W}
\quad\text{且}\quad
(\Sch/X')_{fppf} \times_{x', \mathcal{Z}} \mathcal{W}
$$
是分别由代数空间 $V$ 和 $U'$ 表示的代数叠，且相应表示空间
分别在 $Y$ 和 $X'$ 上光滑。同构 $f$ 诱导 $X$ 上的同构
$\varphi : V \times_Y X \to U' \times_{X'} X$。由《空间的推出》引理
\ref{spaces-pushouts-lemma-pushout-along-thickening} 和
\ref{spaces-pushouts-lemma-equivalence-categories-spaces-pushout-flat}
可知，推出 $V' = V \amalg_{V \times_Y X} U'$ 是 $Y'$ 上光滑的
代数空间；它到 $Y$ 和 $X'$ 的基变换分别恢复 $V$ 和 $U'$，
并且恢复的方式与 $\varphi$ 相容。

\medskip\noindent
令 $W$ 为表示 $\mathcal{W}$ 的代数空间。投影 $V \to W$ 和
$U' \to W$ 在 $V \times_Y X \cong U' \times_{X'} X$ 上作为态射相同，
故推出的泛性质确定代数空间的态射 $V' \to W$。选取概形 $Y_1'$
和满 \'etale 态射 $Y_1' \to V'$。令
$Y_1 = Y \times_{Y'} Y_1'$、$X_1' = X' \times_{Y'} Y_1'$、
$X_1 = X \times_{Y'} Y_1'$。复合
$$
(\Sch/Y_1') \to (\Sch/V') \to (\Sch/W) = \mathcal{W} \to \mathcal{Z}
$$
由 $2$-Yoneda 引理对应于 $\mathcal{Z}$ 在 $Y_1'$ 上的一个对象 $y_1'$；
它在 $Y_1$ 和 $X_1'$ 上的限制分别与 $y|_{Y_1}$ 和 $x'|_{X_1'}$
相符，且相符方式与 $f|_{X_1}$ 相容。因此，我们已经在 $Y'$
上光滑局部地构造出所需对象，证明完成。
\end{proof}








\section{Rim--Schlessinger 条件}
\label{artin-section-RS}

\noindent
下述定义的动机来自引理 \ref{artin-lemma-pushout}，以及《形式形变理论》定义
\ref{formal-defos-definition-RS} 和引理
\ref{formal-defos-lemma-RS-2-categorical}。

\begin{definition}
\label{artin-definition-RS}
设 $S$ 是局部 Noether 概形，$\mathcal{Z}$ 是
$(\Sch/S)_{fppf}$ 上的群胚纤维化范畴。若对每个如下推出
$$
\xymatrix{
X \ar[r] \ar[d] & X' \ar[d] \\
Y \ar[r] & Y' = Y \amalg_X X'
}
$$
（它位于 $S$ 上概形的范畴中），其中
\begin{enumerate}
\item $X$、$X'$、$Y$、$Y'$ 都是 Artin 局部环的谱；
\item $X$、$X'$、$Y$、$Y'$ 在 $S$ 上都是有限型的；并且
\item $X \to X'$（因而 $Y \to Y'$）是闭浸入，
\end{enumerate}
纤维范畴的函子
$$
\mathcal{Z}_{Y'}
\longrightarrow
\mathcal{Z}_Y \times_{\mathcal{Z}_X} \mathcal{Z}_{X'}
$$
都是范畴等价，则称 $\mathcal{Z}$ 满足{\it 条件 (RS)}。
\end{definition}

\noindent
若 $A$ 是剩余域为 $k$ 的 Artin 局部环，则任意态射
$\Spec(A) \to S$ 都是仿射且有限型的，当且仅当诱导态射
$\Spec(k) \to S$ 是有限型的；参见《态射》引理
\ref{morphisms-lemma-Artinian-affine} 和
\ref{morphisms-lemma-artinian-finite-type}。

\begin{lemma}
\label{artin-lemma-algebraic-stack-RS}
设 $\mathcal{X}$ 是局部 Noether 基概形 $S$ 上的代数叠。
则 $\mathcal{X}$ 满足 (RS)。
\end{lemma}

\begin{proof}
由定义和引理 \ref{artin-lemma-pushout} 立即得到。
\end{proof}

\begin{lemma}
\label{artin-lemma-fibre-product-RS}
设 $S$ 是概形。设 $p : \mathcal{X} \to \mathcal{Y}$ 和
$q : \mathcal{Z} \to \mathcal{Y}$ 是 $(\Sch/S)_{fppf}$ 上
群胚纤维化范畴的 $1$-态射。若 $\mathcal{X}$、$\mathcal{Y}$、
$\mathcal{Z}$ 都满足 (RS)，则
$\mathcal{X} \times_\mathcal{Y} \mathcal{Z}$ 也满足 (RS)。
\end{lemma}

\begin{proof}
这是形式论证。设
$$
\xymatrix{
X \ar[r] \ar[d] & X' \ar[d] \\
Y \ar[r] & Y' = Y \amalg_X X'
}
$$
是定义 \ref{artin-definition-RS} 中的图。我们必须证明
$$
(\mathcal{X} \times_{\mathcal{Y}} \mathcal{Z})_{Y'}
\longrightarrow
(\mathcal{X} \times_{\mathcal{Y}} \mathcal{Z})_Y
\times_{(\mathcal{X} \times_{\mathcal{Y}} \mathcal{Z})_X}
(\mathcal{X} \times_{\mathcal{Y}} \mathcal{Z})_{X'}
$$
是等价。利用 $2$-纤维积的定义，这化为
\begin{equation}
\label{artin-equation-RS-fibre-product}
\mathcal{X}_{Y'} \times_{\mathcal{Y}_{Y'}} \mathcal{Z}_{Y'}
\longrightarrow
(\mathcal{X}_Y \times_{\mathcal{Y}_Y} \mathcal{Z}_Y)
\times_{(\mathcal{X}_X \times_{\mathcal{Y}_X} \mathcal{Z}_X)}
(\mathcal{X}_{X'} \times_{\mathcal{Y}_{X'}} \mathcal{Z}_{X'}).
\end{equation}
% P11-E0210: the three frozen-source functors below mix the X, Y, and Z fibres instead of stating each category's own RS equivalence.
已知下列各函子
$$
\mathcal{X}_{Y'} \to \mathcal{X}_Y \times_{\mathcal{Y}_Y} \mathcal{Z}_Y,
\quad
\mathcal{Y}_{Y'} \to \mathcal{X}_X \times_{\mathcal{Y}_X} \mathcal{Z}_X,
\quad
\mathcal{Z}_{Y'} \to
\mathcal{X}_{X'} \times_{\mathcal{Y}_{X'}} \mathcal{Z}_{X'}
$$
都是等价。式 (\ref{artin-equation-RS-fibre-product}) 右端的一个对象是系统
$$
((x_Y, z_Y, \phi_Y), (x_{X'}, z_{X'}, \phi_{X'}), (\alpha, \beta)).
$$
% P11-E0211: the frozen source next uses x_{Y'} and z_{Y'} inside the input triples where the displayed right-hand object supplies x_{X'} and z_{X'}.
于是 $(x_Y, x_{Y'}, \alpha)$ 同构于 $\mathcal{X}_{Y'}$ 中某个对象
$x_{Y'}$ 的像，而 $(z_Y, z_{Y'}, \beta)$ 同构于
$\mathcal{Z}_{Y'}$ 中某个对象 $z_{Y'}$ 的像。态射对
$(\phi_Y, \phi_{X'})$ 对应于 $x_{Y'}$ 和 $z_{Y'}$ 在
$\mathcal{Y}_{Y'}$ 中的像之间的一个态射 $\psi$。于是
$(x_{Y'}, z_{Y'}, \psi)$ 是式 (\ref{artin-equation-RS-fibre-product})
左端的一个对象，它映到给定的右端对象。这证明
(\ref{artin-equation-RS-fibre-product}) 是本质满的。其全忠实性的证明从略。
\end{proof}





\section{形变范畴}
\label{artin-section-deformation-categories}

\noindent
我们把上文引入的记号与《形式形变理论》一章中的记号对应起来。

\begin{lemma}
\label{artin-lemma-deformation-category}
设 $S$ 是局部 Noether 概形，$\mathcal{X}$ 是
$(\Sch/S)_{fppf}$ 上满足 (RS) 的群胚纤维化范畴。对任意在 $S$
上有限型的域 $k$，以及 $\mathcal{X}$ 在 $k$ 上的任意对象 $x_0$，
预形变范畴
$p : \mathcal{F}_{\mathcal{X}, k, x_0} \to \mathcal{C}_\Lambda$
（式 (\ref{artin-equation-predeformation-category})）是形变范畴；参见
《形式形变理论》定义
\ref{formal-defos-definition-deformation-category}。
\end{lemma}

\begin{proof}
令 $\mathcal{F} = \mathcal{F}_{\mathcal{X}, k, x_0}$。设
$f_1 : A_1 \to A$ 和 $f_2 : A_2 \to A$ 是
$\mathcal{C}_\Lambda$ 中的环同态，且 $f_2$ 为满射。须证明函子
$$
\mathcal{F}(A_1 \times_A A_2)
\longrightarrow
\mathcal{F}(A_1) \times_{\mathcal{F}(A)} \mathcal{F}(A_2)
$$
是等价；参见《形式形变理论》引理
\ref{formal-defos-lemma-RS-2-categorical}。令
$X = \Spec(A)$、$X' = \Spec(A_2)$、$Y = \Spec(A_1)$ 以及
$Y' = \Spec(A_1 \times_A A_2)$。注意，在概形范畴中
$Y' = Y \amalg_X X'$；参见《更多态射》引理
\ref{more-morphisms-lemma-pushout-along-thickening}。在纤维范畴函子的图
$$
\xymatrix{
\mathcal{X}_{Y'} \ar[r] \ar[d] &
\mathcal{X}_Y \times_{\mathcal{X}_X} \mathcal{X}_{X'} \ar[d] \\
\mathcal{X}_{\Spec(k)} \ar@{=}[r] & \mathcal{X}_{\Spec(k)}
}
$$
中，由定义 \ref{artin-definition-RS}，上方横箭头是等价。由于
$\mathcal{F}(B)$ 是 $\mathcal{X}_{\Spec(B)}$ 中带有在 $k$ 上与
$x_0$ 的一个认同的对象所成的范畴，结论成立。
\end{proof}

\begin{remark}
\label{artin-remark-deformation-category-implies}
设 $S$ 是局部 Noether 概形，$\mathcal{X}$ 是
$(\Sch/S)_{fppf}$ 上的群胚纤维化范畴。设 $k$ 是 $S$ 上有限型的域，
$x_0$ 是 $\mathcal{X}$ 在 $k$ 上的对象。设
$p : \mathcal{F} \to \mathcal{C}_\Lambda$ 如式
(\ref{artin-equation-predeformation-category})。若 $\mathcal{F}$ 是形变范畴，
即 $\mathcal{F}$ 满足 Rim--Schlessinger 条件 (RS)，则由
《形式形变理论》引理 \ref{formal-defos-lemma-RS-implies-S1-S2}，
$\mathcal{F}$ 满足 Schlessinger 条件 (S1) 和 (S2)。令
$\overline{\mathcal{F}}$ 为同构类函子；参见《形式形变理论》评注
\ref{formal-defos-remarks-cofibered-groupoids}
中的 (\ref{formal-defos-item-associated-functor-isomorphism-classes})。
于是 $\overline{\mathcal{F}}$ 也满足 (S1) 和 (S2)；参见
《形式形变理论》引理
\ref{formal-defos-lemma-S1-S2-associated-functor}。特别地，
这适用于引理 \ref{artin-lemma-deformation-category} 的情形。
\end{remark}




\section{域的变换}
\label{artin-section-change-of-field}

\noindent
本节对应于《形式形变理论》一节
\ref{formal-defos-section-change-of-field}。正如该处所指出的，为讨论
域变换下发生的情形，需要写成 $\mathcal{C}_{\Lambda, k}$，而不是
$\mathcal{C}_\Lambda$。在下述引理中，我们采用《形式形变理论》情形
\ref{formal-defos-situation-change-of-fields} 中引入的记号
$\mathcal{F}_{l/k}$。

\begin{lemma}
\label{artin-lemma-change-of-field}
设 $S$ 是局部 Noether 概形，$\mathcal{X}$ 是
$(\Sch/S)_{fppf}$ 上的群胚纤维化范畴。设 $k$ 是 $S$ 上有限型的域，
$l/k$ 是有限扩张。
% P11-E0212: the frozen source says x_0 is an object of undefined \mathcal{F}, while the surrounding construction requires \mathcal{X}; the mathematical token is preserved.
设 $x_0$ 是 $\mathcal{F}$ 位于 $\Spec(k)$ 上的一个对象。
% P11-E0213: the frozen source omits "by" in the naming construction "Denote x_{l,0} the restriction".
记 $x_{l, 0}$ 为 $x_0$ 在 $\Spec(l)$ 上的限制。则存在
在 $\mathcal{C}_{\Lambda, l}$ 上以群胚为纤维的余纤维化范畴的典范函子
$$
(\mathcal{F}_{\mathcal{X}, k , x_0})_{l/k}
\longrightarrow
\mathcal{F}_{\mathcal{X}, l, x_{l, 0}}
$$
。若 $\mathcal{X}$ 满足 (RS)，则此函子是等价。
\end{lemma}

\begin{proof}
考虑式 (\ref{artin-equation-factor}) 中的分解
$$
\Spec(l) \to \Spec(B) \to S
$$
。由定义，
$$
(\mathcal{F}_{\mathcal{X}, k , x_0})_{l/k}(B) =
\mathcal{F}_{\mathcal{X}, k, x_0}(B \times_l k)
$$
；参见《形式形变理论》情形
\ref{formal-defos-situation-change-of-fields}。因此，其中一个对象是
$\mathcal{X}$ 的态射 $x_0 \to x$，它位于态射
$\Spec(k) \to \Spec(B \times_l k)$ 之上。
% P11-E0214: the frozen source omits the article in "Choosing pullback functor".
为 $\mathcal{X}$ 选取拉回函子后，可将 $x_0 \to x$ 对应到态射
$x_{l, 0} \to x_B$，其中 $x_B$ 是 $x$ 在 $\Spec(B)$ 上的限制
（经由来自 $B \times_l k \subset B$ 的态射
$\Spec(B) \to \Spec(B \times_l k)$）。此构造关于 $B$ 函子性地变化，
并与态射相容。

\medskip\noindent
下面假设 $\mathcal{X}$ 满足 (RS)。考虑诸图
$$
\vcenter{
\xymatrix{
l & B \ar[l] \\
k \ar[u] & B \times_l k \ar[l] \ar[u]
}
}
\quad\text{且}\quad
\vcenter{
\xymatrix{
\Spec(l) \ar[d] \ar[r] & \Spec(B) \ar[d] \\
\Spec(k) \ar[r] & \Spec(B \times_l k)
}
}
$$
左图是环的纤维积，右图是概形范畴中的推出；参见《更多态射》引理
\ref{more-morphisms-lemma-pushout-along-thickening}。这些概形在 $S$
上都是有限型的（参见定义 \ref{artin-definition-RS} 后的评注）。因此，
(RS) 给出纤维范畴的等价
$$
\mathcal{X}_{\Spec(B \times_l k)}
\longrightarrow
\mathcal{X}_{\Spec(k)}
\times_{\mathcal{X}_{\Spec(l)}}
\mathcal{X}_{\Spec(B)}
$$
。这说明上面定义的函子给出纤维范畴的等价。因此，由《范畴》引理
\ref{categories-lemma-equivalence-fibred-categories} 的对偶，
该函子是以群胚为纤维的余纤维化范畴之间的等价。
\end{proof}








\section{切空间}
\label{artin-section-tangent-spaces}

\noindent
设 $S$ 是局部 Noether 概形，$\mathcal{X}$ 是
$(\Sch/S)_{fppf}$ 上的群胚纤维化范畴。设 $k$ 是 $S$ 上有限型的域，
$x_0$ 是 $\mathcal{X}$ 在 $k$ 上的对象。在《形式形变理论》一节
\ref{formal-defos-section-tangent-spaces} 中，我们定义了预形变范畴
$\mathcal{F}_{\mathcal{X}, k, x_0}$ 的{\it 切空间}
\begin{equation}
\label{artin-equation-tangent-space}
T\mathcal{F}_{\mathcal{X}, k, x_0} =
\left\{
\begin{matrix}
\text{isomorphism classes of morphisms}\\
x_0 \to x\text{ 相对于 }\Spec(k) \to \Spec(k[\epsilon])
\end{matrix}
\right\}
\end{equation}
。在《形式形变理论》一节
\ref{formal-defos-section-infinitesimal-automorphisms} 中，我们定义了
\begin{equation}
\label{artin-equation-infinitesimal-automorphisms}
\text{Inf}(\mathcal{F}_{\mathcal{X}, k, x_0}) =
\Ker\left(
\text{Aut}_{\Spec(k[\epsilon])}(x'_0) \to \text{Aut}_{\Spec(k)}(x_0)
\right)
\end{equation}
，其中 $x_0'$ 是 $x_0$ 到 $\Spec(k[\epsilon])$ 的拉回。
若 $\mathcal{X}$ 满足 Rim--Schlessinger 条件 (RS)，则由
《形式形变理论》引理
\ref{formal-defos-lemma-tangent-space-vector-space}
（由引理 \ref{artin-lemma-deformation-category} 和评注
\ref{artin-remark-deformation-category-implies}，其假设成立），
$T\mathcal{F}_{\mathcal{X}, k, x_0}$ 带有自然的 $k$-向量空间结构。
此外，《形式形变理论》引理
\ref{formal-defos-lemma-infaut-vector-space} 表明，
$\text{Inf}(\mathcal{F}_{\mathcal{X}, k, x_0})$ 带有自然的
$k$-向量空间结构，使加法与自同构的复合一致。一个自然条件是要求
这些向量空间为有限维。

\medskip\noindent
下述引理说明，若 $\mathcal{X}$ 在 $S$ 上局部有限型，则此条件成立
（参见《叠的态射》一节 \ref{stacks-morphisms-section-finite-type}）。

\begin{lemma}
\label{artin-lemma-finite-dimension}
设 $S$ 是局部 Noether 概形。假设
\begin{enumerate}
\item $\mathcal{X}$ 是代数叠；
\item $U$ 是 $S$ 上局部有限型的概形；并且
\item $(\Sch/U)_{fppf} \to \mathcal{X}$ 是满光滑态射。
\end{enumerate}
% P11-E0215: the frozen source omits punctuation between the section-reference modifier and the main clause.
则对第 \ref{artin-section-predeformation-categories} 节中的任意
$\mathcal{F} = \mathcal{F}_{\mathcal{X}, k, x_0}$，切空间
$T\mathcal{F}$ 和无穷小自同构空间 $\text{Inf}(\mathcal{F})$
在 $k$ 上都是有限维的。
\end{lemma}

\begin{proof}
记 $\mathcal{U} = (\Sch/U)_{fppf}$。由代数叠的定义，$1$-态射
$\mathcal{U} \to \mathcal{X}$ 可由代数空间表示。因此，特别地，
2-纤维积
$$
\mathcal{U}_{x_0} = (\Sch/\Spec(k))_{fppf} \times_\mathcal{X} \mathcal{U}
$$
可由 $\Spec(k)$ 上的代数空间 $U_{x_0}$ 表示。于是
$U_{x_0} \to \Spec(k)$ 光滑且满（特别地，$U_{x_0}$ 非空）。由
《域上的空间》引理
\ref{spaces-over-fields-lemma-smooth-separable-closed-points-dense}
，可找到有限扩张 $l/k$ 和 $k$ 上的点
$\Spec(l) \to U_{x_0}$。由引理 \ref{artin-lemma-change-of-field} 以及
$\mathcal{X}$ 满足 (RS) 这一事实，有
$$
(\mathcal{F}_{\mathcal{X}, k , x_0})_{l/k} =
\mathcal{F}_{\mathcal{X}, l, x_{l, 0}}
$$
。因而
$$
T\mathcal{F} \otimes_k l \cong T\mathcal{F}_{\mathcal{X}, l, x_{l, 0}}
\quad\text{且}\quad
\text{Inf}(\mathcal{F}) \otimes_k l \cong
\text{Inf}(\mathcal{F}_{\mathcal{X}, l, x_{l, 0}})
$$
；这是由《形式形变理论》引理
\ref{formal-defos-lemma-tangent-space-change-of-field} 和
\ref{formal-defos-lemma-inf-aut-change-of-field}
得到的（由引理 \ref{artin-lemma-algebraic-stack-RS}、
\ref{artin-lemma-deformation-category} 和评注
\ref{artin-remark-deformation-category-implies}，这些引理可以应用）。
因此，只需证明 $T\mathcal{F}_{\mathcal{X}, l, x_{l, 0}}$ 和
$\text{Inf}(\mathcal{F}_{\mathcal{X}, l, x_{l, 0}})$ 在 $l$ 上有限维。
注意，$x_{l, 0}$ 来自 $\mathcal{U}$ 在 $l$ 上的一个点 $u_0$。

\medskip\noindent
我们暂时中断论证，说明无穷小自同构的结论可由切空间的结论推出。
令 $\mathcal{R} = \mathcal{U} \times_\mathcal{X} \mathcal{U}$，
并令 $r_0$ 为 $\mathcal{R}$ 的 $l$-值点
$(u_0, u_0, \text{id}_{x_0})$。结合引理
\ref{artin-lemma-fibre-product-deformation-categories} 和
《形式形变理论》引理
\ref{formal-defos-lemma-deformation-functor-diagonal}
可知
$$
\text{Inf}(\mathcal{F}_{\mathcal{X}, l, x_{l, 0}})
\subset
T\mathcal{F}_{\mathcal{R}, l, r_0}
$$
。注意，$\mathcal{R}$ 是代数叠；参见《代数叠》引理
\ref{algebraic-lemma-2-fibre-product-general}。此外，$\mathcal{R}$
可由代数空间 $R$ 表示，且该空间经任一投影在 $U$ 上光滑；参见
《代数叠》引理 \ref{algebraic-lemma-stack-presentation}。
% P11-E0216: the frozen source has "choose an scheme" and joins the choice directly to "we see" without a grammatical connector.
因此，选取概形 $U'$ 和满 \'etale 态射 $U' \to R$，可知 $U'$
在 $U$ 上光滑，因而在 $S$ 上局部有限型。由于
$(\Sch/U')_{fppf} \to \mathcal{R}$ 满且光滑，问题已经归结为切空间情形。

\medskip\noindent
函子 (\ref{artin-equation-functoriality})
$$
\mathcal{F}_{\mathcal{U}, l, u_0}
\longrightarrow
\mathcal{F}_{\mathcal{X}, l, x_{l, 0}}
$$
由引理 \ref{artin-lemma-formally-smooth-on-deformation-categories} 是光滑的。
所诱导的切空间映射
$$
T\mathcal{F}_{\mathcal{U}, l, u_0}
\longrightarrow
T\mathcal{F}_{\mathcal{X}, l, x_{l, 0}}
$$
是 $l$-线性的（由《形式形变理论》引理
\ref{formal-defos-lemma-k-linear-differential}），也是满的（因为由
《形式形变理论》引理
\ref{formal-defos-lemma-smooth-morphism-essentially-surjective}，
预形变范畴的光滑映射诱导切空间的满映射）。因此，只需证明与可表代数叠
$\mathcal{U}$ 相联系的形变空间在点 $u_0$ 处的切空间为有限维。
% P11-E0217: the frozen source says "deformation space" although the surrounding construction is the predeformation/deformation category associated to \mathcal{U}.
取仿射开集 $\Spec(R) \subset U$，使 $u_0 : \Spec(l) \to U$ 分解通过
$\Spec(R)$，且 $\Spec(R) \to S$ 分解通过
$\Spec(\Lambda) \subset S$。令 $\mathfrak m_R \subset R$ 为与 $u_0$
对应的 $\Lambda$-代数同态 $\varphi_0 : R \to l$ 的核。注意，$R$
在 Noether 环 $\Lambda$ 上有限型，因而是 Noether 环。因此，
$\mathfrak m_R = (f_1, \ldots, f_n)$ 是有限生成理想。我们有
$$
T\mathcal{F}_{\mathcal{U}, l, u_0}
=
\{\varphi : R \to l[\epsilon] \mid
\varphi \text{ is a } \Lambda\text{-algebra map and }
\varphi \bmod \epsilon = \varphi_0\}
$$
% P11-E0218: the frozen source joins the independent clauses around "hence" without punctuation.
右端的一个元素由它在 $f_1, \ldots, f_n$ 上的值确定；因此其维数至多为
$n$，结论成立。这里略去若干细节。
\end{proof}

\begin{lemma}
\label{artin-lemma-fibre-product-tangent-spaces}
设 $S$ 是局部 Noether 概形。设 $p : \mathcal{X} \to \mathcal{Y}$ 和
$q : \mathcal{Z} \to \mathcal{Y}$ 是 $(\Sch/S)_{fppf}$ 上
群胚纤维化范畴的 $1$-态射。假设 $\mathcal{X}, \mathcal{Y}, \mathcal{Z}$
满足 (RS)。设 $k$ 是 $S$ 上有限型的域，$w_0$ 是
$\mathcal{W} = \mathcal{X} \times_\mathcal{Y} \mathcal{Z}$ 在 $k$
上的一个对象。
% P11-E0219: the frozen source omits "by" in "Denote x_0, y_0, z_0 the objects".
将由 $w_0$ 得到的 $\mathcal{X}$、$\mathcal{Y}$、$\mathcal{Z}$ 的对象
分别记为 $x_0, y_0, z_0$。则存在 $k$-向量空间的 $6$-项正合列
$$
\xymatrix{
0 \ar[r] &
\text{Inf}(\mathcal{F}_{\mathcal{W}, k, w_0}) \ar[r] &
\text{Inf}(\mathcal{F}_{\mathcal{X}, k, x_0}) \oplus
\text{Inf}(\mathcal{F}_{\mathcal{Z}, k, z_0}) \ar[r] &
\text{Inf}(\mathcal{F}_{\mathcal{Y}, k, y_0}) \ar[lld] \\
 &
T\mathcal{F}_{\mathcal{W}, k, w_0} \ar[r] &
T\mathcal{F}_{\mathcal{X}, k, x_0} \oplus
T\mathcal{F}_{\mathcal{Z}, k, z_0} \ar[r] &
T\mathcal{F}_{\mathcal{Y}, k, y_0}
}
$$
。
\end{lemma}

\begin{proof}
由引理 \ref{artin-lemma-fibre-product-RS}，$\mathcal{W}$ 满足 (RS)，
因而本引理的陈述有意义。为证明本引理，应用引理
\ref{artin-lemma-fibre-product-deformation-categories}、
\ref{artin-lemma-deformation-category} 以及《形式形变理论》引理
\ref{formal-defos-lemma-deformation-categories-fiber-product-morphisms}。
\end{proof}






\section{形式对象}
\label{artin-section-formal-objects}

\noindent
本节把《形式形变理论》一章中已经定义的若干概念移到当前情形。
下文说“$R$ 是 $S$-代数”，是指 $R$ 是一个环，并配备概形态射
$\Spec(R) \to S$。

\begin{definition}
\label{artin-definition-formal-objects}
设 $S$ 是局部 Noether 概形，并设
$p : \mathcal{X} \to (\Sch/S)_{fppf}$ 是群胚纤维化范畴。
\begin{enumerate}
\item $\mathcal{X}$ 的一个{\it 形式对象} $\xi = (R, \xi_n, f_n)$
由如下数据组成：Noether 完备局部 $S$-代数 $R$；$\mathcal{X}$ 中位于
$\Spec(R/\mathfrak m_R^n)$ 上的对象 $\xi_n$；以及 $\mathcal{X}$ 中位于
% P11-E0220: the frozen source switches from \mathfrak m_R to undefined \mathfrak m in the transition morphisms and residue-field condition.
$\Spec(R/\mathfrak m^n) \to \Spec(R/\mathfrak m^{n + 1})$ 上的态射
$f_n : \xi_n \to \xi_{n + 1}$。这里要求 $R/\mathfrak m$ 是 $S$
上有限型的域。
\item 一个{\it 形式对象的态射}
$a : \xi = (R, \xi_n, f_n) \to \eta = (T, \eta_n, g_n)$
由态射 $a_n : \xi_n \to \eta_n$ 给出，并且对每个 $n$，图
$$
\xymatrix{
\xi_n \ar[r]_{f_n} \ar[d]_{a_n} & \xi_{n + 1} \ar[d]^{a_{n + 1}} \\
\eta_n \ar[r]^{g_n} & \eta_{n + 1}
}
$$
交换。应用函子 $p$，得到相容的态射族
$\Spec(R/\mathfrak m_R^n) \to \Spec(T/\mathfrak m_T^n)$，从而得到
$S$ 上的态射 $a_0 : \Spec(R) \to \Spec(T)$。我们称 $a$
{\it 位于} $a_0$ 之上。
\end{enumerate}
\end{definition}

\noindent
这样就得到 $\mathcal{X}$ 的形式对象所成的范畴。

\begin{remark}
\label{artin-remark-formal-objects-match}
设 $S$ 是局部 Noether 概形，并设
$p : \mathcal{X} \to (\Sch/S)_{fppf}$ 是群胚纤维化范畴。
设 $\xi = (R, \xi_n, f_n)$ 是形式对象。令 $k = R/\mathfrak m$ 且
$x_0 = \xi_1$。形式对象 $\xi$ 定义了预形变范畴
$\mathcal{F}_{\mathcal{X}, k, x_0}$ 的一个形式对象 $\xi$。
这立即来自上面的定义 \ref{artin-definition-formal-objects}、
《形式形变理论》定义 \ref{formal-defos-definition-formal-objects}，
以及第 \ref{artin-section-predeformation-categories} 节中对预形变范畴
$\mathcal{F}_{\mathcal{X}, k, x_0}$ 的构造。
\end{remark}

\noindent
若 $F : \mathcal{X} \to \mathcal{Y}$ 是 $(\Sch/S)_{fppf}$ 上
群胚纤维化范畴的 $1$-态射，则 $F$ 也诱导形式对象范畴之间的函子。

\begin{lemma}
\label{artin-lemma-smooth-lift-formal}
设 $S$ 是局部 Noether 概形。设 $F : \mathcal{X} \to \mathcal{Y}$ 是
$(\Sch/S)_{fppf}$ 上群胚纤维化范畴的 $1$-态射。设
$\eta = (R, \eta_n, g_n)$ 是 $\mathcal{Y}$ 的形式对象，$\xi_1$ 是
$\mathcal{X}$ 的对象，且 $F(\xi_1) \cong \eta_1$。若 $F$ 在对象上
形式光滑（参见《可表示性判据》一节
\ref{criteria-section-formally-smooth}），则存在 $\mathcal{X}$ 的形式对象
$\xi = (R, \xi_n, f_n)$，使 $F(\xi) \cong \eta$。
\end{lemma}

\begin{proof}
注意，每个态射
$\Spec(R/\mathfrak m^n) \to \Spec(R/\mathfrak m^{n + 1})$ 都是 $S$
上仿射概形的一阶加厚。因此，对 $F$ 的假设意味着可逐次把 $\xi_1$
提升为 $\mathcal{X}$ 的对象 $\xi_2, \xi_3, \ldots$，并配备相容同构
% P11-E0221: the frozen source switches xi and eta in both displayed compatibility isomorphisms.
$\eta_n|_{\Spec(R/\mathfrak m^{n - 1})} \cong \eta_{n - 1}$
和 $F(\eta_n) \cong \xi_n$。
\end{proof}

\noindent
设 $S$ 是局部 Noether 概形，并设
$p : \mathcal{X} \to (\Sch/S)_{fppf}$ 是群胚纤维化范畴。假设 $x$
是 $\mathcal{X}$ 在 $R$ 上的对象，其中 $R$ 是 Noether 完备局部
$S$-代数，其剩余域在 $S$ 上有限型。于是，可以考虑限制系统
$\xi_n = x|_{\Spec(R/\mathfrak m^n)}$，并赋予它由限制的传递性所得的
自然态射 $\xi_1 \to \xi_2 \to \ldots$。因此，
$\xi = (R, \xi_n, \xi_n \to \xi_{n + 1})$ 是 $\mathcal{X}$ 的形式对象。
此构造关于对象 $x$ 函子性地变化。这样得到函子
\begin{equation}
\label{artin-equation-approximation}
\left\{
\begin{matrix}
\text{objects }x\text{ 属于 }\mathcal{X} \text{ 使得 }p(x) = \Spec(R) \\
\text{其中 }R\text{ is Noetherian complete local}\\
\text{其中 }R/\mathfrak m\text{ of finite type over }S
\end{matrix}
\right\}
\longrightarrow
\left\{
\begin{matrix}
\text{formal objects of }\mathcal{X}
\end{matrix}
\right\}
\end{equation}
确切地说，左端是 $\mathcal{X}$ 中由所述对象组成的全子范畴，
右端是定义 \ref{artin-definition-formal-objects} 中 $\mathcal{X}$ 的形式对象范畴。

\begin{definition}
\label{artin-definition-effective}
设 $S$ 是局部 Noether 概形，$\mathcal{X}$ 是
$(\Sch/S)_{fppf}$ 上的群胚纤维化范畴。若 $\mathcal{X}$ 的形式对象
$\xi = (R, \xi_n, f_n)$ 位于函子 (\ref{artin-equation-approximation})
的本质像中，则称它是{\it 有效的}。
\end{definition}

\noindent
若这个群胚纤维化范畴是代数叠，则由下述引理，每个形式对象都是有效的。

\begin{lemma}
\label{artin-lemma-effective}
设 $S$ 是局部 Noether 概形，$\mathcal{X}$ 是 $S$ 上的代数叠。
函子 (\ref{artin-equation-approximation}) 是等价。
\end{lemma}

\begin{proof}
情形 I：$\mathcal{X}$ 可由概形表示。设
$\mathcal{X} = (\Sch/X)_{fppf}$，其中 $X$ 是 $S$ 上的某个概形。
展开定义后，需要证明如下事实：给定 Noether 完备局部 $S$-代数 $R$，
且 $R/\mathfrak m$ 在 $S$ 上有限型，则
$$
\Mor_S(\Spec(R), X) \longrightarrow \lim \Mor_S(\Spec(R/\mathfrak m^n), X)
$$
为双射。这来自《形式空间》引理
\ref{formal-spaces-lemma-map-into-scheme}。

\medskip\noindent
情形 II：$\mathcal{X}$ 可由代数空间表示。设 $\mathcal{X}$ 由 $X$ 表示。
仍须证明
$$
\Mor_S(\Spec(R), X) \longrightarrow \lim \Mor_S(\Spec(R/\mathfrak m^n), X)
$$
对上述 $R$ 为双射。这就是《形式空间》引理
\ref{formal-spaces-lemma-map-into-algebraic-space}。

\medskip\noindent
情形 III：一般的代数叠情形。一般地，式
(\ref{artin-equation-approximation}) 的左右两端，都是下述范畴上的
群胚纤维化范畴：对象为 $S$ 上的仿射概形，它们是剩余域在 $S$
上有限型的 Noether 完备局部环之谱。在下面的证明中还会看到，
对于这个范畴上的某个拓扑，它们构成叠。

\medskip\noindent
% P11-E0222: the frozen source says "prove fully faithfulness" instead of "prove full faithfulness".
先证明全忠实性。设 $R$ 是 Noether 完备局部 $S$-代数，且
$k = R/\mathfrak m$ 在 $S$ 上有限型。设 $x, x'$ 是 $\mathcal{X}$
在 $R$ 上的对象。由于 $\mathcal{X}$ 是代数叠，
$\mathit{Isom}(x, x')$ 可由 $\Spec(R)$ 上的代数空间 $I$ 表示；
参见《代数叠》引理 \ref{algebraic-lemma-representable-diagonal}。
把情形 II 应用于 $\Spec(R)$ 上的 $I$，立即得到
(\ref{artin-equation-approximation}) 在 $\Spec(R)$ 上的纤维范畴之间全忠实。
因此，由《范畴》引理 \ref{categories-lemma-equivalence-fibred-categories}，
该函子全忠实。

\medskip\noindent
本质满性。设 $\xi = (R, \xi_n, f_n)$ 是 $\mathcal{X}$ 的形式对象。
选取 $S$ 上的概形 $U$ 和满光滑态射
$f : (\Sch/U)_{fppf} \to \mathcal{X}$。对每个 $n$，考虑纤维积
$$
(\Sch/\Spec(R/\mathfrak m^n))_{fppf}
\times_{\xi_n, \mathcal{X}, f}
(\Sch/U)_{fppf}
$$
。由假设，它可由代数空间 $V_n$ 表示，且该空间在
$\Spec(R/\mathfrak m^n)$ 上满且光滑。态射
$f_n : \xi_n \to \xi_{n + 1}$ 诱导代数空间的 Cartesian 方块
$$
\xymatrix{
V_{n + 1} \ar[d] & V_n \ar[d] \ar[l] \\
\Spec(R/\mathfrak m^{n + 1}) & \Spec(R/\mathfrak m^n) \ar[l]
}
$$
。由《域上的空间》引理
\ref{spaces-over-fields-lemma-smooth-separable-closed-points-dense}
，可找到有限可分扩张 $k'/k$ 和 $k$ 上的点
$v'_1 : \Spec(k') \to V_1$。令 $R \subset R'$ 为剩余域扩张是
$k'/k$ 的有限 \'etale 扩张；由《代数》引理
\ref{algebra-lemma-henselian-cat-finite-etale} 和
\ref{algebra-lemma-complete-henselian}，它存在且唯一。对
$n = 2, 3, 4, \ldots$，把光滑性的无穷小提升判据（参见
《代数空间态射进阶》引理
\ref{spaces-more-morphisms-lemma-smooth-formally-smooth})
应用于 $V_n \to \Spec(R/\mathfrak m^n)$，可逐次找到
$\Spec(R/\mathfrak m^n)$ 上的态射
$v'_n : \Spec(R'/(\mathfrak m')^n) \to V_n$，使下图交换：
$$
\xymatrix{
\Spec(R'/(\mathfrak m')^{n + 1}) \ar[d]_{v'_{n + 1}} &
\Spec(R'/(\mathfrak m')^n) \ar[d]^{v'_n} \ar[l] \\
V_{n + 1} & V_n \ar[l]
}
$$
与投影态射 $V_n \to U$ 复合，得到相容态射系
$u'_n : \Spec(R'/(\mathfrak m')^n) \to U$。由情形 I，族 $(u'_n)$
来自唯一态射 $u' : \Spec(R') \to U$。
% P11-E0223: the frozen source omits "by" in "Denote x' the object".
把 $1$-态射 $f$ 应用于 $u'$ 所得的 $\mathcal{X}$ 在
$\Spec(R')$ 上的对象记为 $x'$。由构造，存在形式对象的态射
$$
(\ref{artin-equation-approximation})(x') =
(R', x'|_{\Spec(R'/(\mathfrak m')^n)}, \ldots)
\longrightarrow
(R, \xi_n, f_n)
$$
，它位于 $\Spec(R') \to \Spec(R)$ 之上。
% P11-E0224: the frozen source says the ring R' tensor_R R' is a finite product of spectra, mixing ring and scheme types.
注意，$R' \otimes_R R'$ 是若干 Noether 完备局部环的谱的有限乘积，
而当前讨论适用于这些谱。
% P11-E0225: the frozen source omits "by" in the naming construction for p_0 and p_1.
把两个投影记为
$p_0, p_1 : \Spec(R' \otimes_R R') \to \Spec(R')$。
% P11-E0226: the frozen source says "By the fully faithfulness" instead of "By the full faithfulness".
由上面证明的全忠实性，存在典范同构
$\varphi : p_0^*x' \to p_1^*x'$，因为在
$\Spec((R' \otimes_R R')/\mathfrak m^n(R' \otimes_R R'))$ 上已有这样的同构。
同构 $\varphi$ 满足上链条件的证明从略（参见《叠》定义
\ref{stacks-definition-descent-data}）。由于
$\{\Spec(R') \to \Spec(R)\}$ 是 fppf 覆盖，断定 $x'$ 下降为
$\mathcal{X}$ 在 $\Spec(R)$ 上的对象 $x$。
% P11-E0227: the frozen source refers to x_n, whereas the formal object's components are xi_n.
$x_n$ 是 $x$ 在 $\Spec(R/\mathfrak m^n)$ 上的限制这一点的证明从略。
\end{proof}

\begin{lemma}
\label{artin-lemma-fibre-product-effective}
设 $S$ 是概形。设 $p : \mathcal{X} \to \mathcal{Y}$ 和
$q : \mathcal{Z} \to \mathcal{Y}$ 是 $(\Sch/S)_{fppf}$ 上
群胚纤维化范畴的 $1$-态射。若函子
(\ref{artin-equation-approximation}) 对 $\mathcal{X}$、$\mathcal{Y}$ 和
$\mathcal{Z}$ 都是等价，则它对
$\mathcal{X} \times_\mathcal{Y} \mathcal{Z}$ 也是等价。
\end{lemma}

\begin{proof}
对 $\mathcal{X} \times_\mathcal{Y} \mathcal{Z}$ 而言，式
(\ref{artin-equation-approximation}) 的左端和右端，分别就是关于
$\mathcal{X}$、$\mathcal{Z}$ 在 $\mathcal{Y}$ 上的式
(\ref{artin-equation-approximation}) 左端和右端的 $2$-纤维积。由于取
$2$-纤维积与范畴等价相容，结论随即得到；参见《范畴》引理
\ref{categories-lemma-equivalence-2-fibre-product}。
\end{proof}






\section{逼近}
\label{artin-section-approximation}

\noindent
Michael Artin 的一个根本洞见是：可以逼近保极限叠的对象。具体地，
给定这个叠在某个 Noether 完备局部环上的对象 $x$，可以找到代数环上的
对象 $x_A$，使其“接近”$x$。这里，代数环是指有限型 $S$-代数，
而接近是指在进拓扑意义下接近。本节以简单而一般的形式陈述这一结果。

\medskip\noindent
为表述这个结果，需要汇集 Stacks Project 不同部分中的若干定义。首先，
在《可表示性判据》一节 \ref{criteria-section-limit-preserving} 中，
我们对概形范畴上群胚纤维化范畴的 $1$-态射引入了
{\it 在对象上保极限}的概念。在《更多代数》定义
\ref{more-algebra-definition-G-ring} 中，我们定义了 {\it G-环}。
设 $S$ 是局部 Noether 概形，并设 $A$ 是 $S$-代数。若
$\Spec(A) \to S$ 为有限型态射，则称 $A$ {\it 在 $S$ 上有限型}，
或称它是{\it 有限型 $S$-代数}。在这种情形下，$A$ 是 Noether 环。
最后，给定环 $A$ 和理想 $I$，记
$\text{Gr}_I(A) = \bigoplus I^n/I^{n + 1}$。

\begin{lemma}
\label{artin-lemma-approximate}
设 $S$ 是局部 Noether 概形，并设
$p : \mathcal{X} \to (\Sch/S)_{fppf}$ 是群胚纤维化范畴。设 $x$
是 $\mathcal{X}$ 位于 $\Spec(R)$ 上的对象，其中 $R$ 是 Noether
完备局部环，其剩余域 $k$ 在 $S$ 上有限型。令 $s \in S$ 为
$\Spec(k) \to S$ 的像。假设 (a) $\mathcal{O}_{S, s}$ 是 G-环，
并且 (b) $p$ 在对象上保极限。则对每个整数 $N \geq 1$，存在
\begin{enumerate}
\item 有限型 $S$-代数 $A$；
\item 极大理想 $\mathfrak m_A \subset A$；
\item $\mathcal{X}$ 在 $\Spec(A)$ 上的对象 $x_A$；
\item $S$-同构 $R/\mathfrak m_R^N \cong A/\mathfrak m_A^N$；
\item 同构
$x|_{\Spec(R/\mathfrak m_R^N)} \cong x_A|_{\Spec(A/\mathfrak m_A^N)}$
，它与 (4) 相容；以及
\item 分次 $k$-代数的同构
$\text{Gr}_{\mathfrak m_R}(R) \cong \text{Gr}_{\mathfrak m_A}(A)$
。
\end{enumerate}
\end{lemma}

\begin{proof}
选取仿射开集 $\Spec(\Lambda) \subset S$，使 $k$ 是有限
$\Lambda$-代数；参见《态射》引理
\ref{morphisms-lemma-point-finite-type}。我们可以并且确实用
$\Spec(\Lambda)$ 代替 $S$。

\medskip\noindent
可将 $R$ 写成有向余极限 $R = \colim C_j$，其中每个 $C_j$ 都是
有限型 $\Lambda$-代数（参见《代数》引理
\ref{algebra-lemma-ring-colimit-fp}）。由假设 (b)，对象 $x$ 同构于
某个 $C_j$ 上对象的限制。因此，可以选取有限型 $\Lambda$-代数 $C$、
$\Lambda$-代数同态 $C \to R$，以及 $\mathcal{X}$ 在
$\Spec(C)$ 上的对象 $x_C$，使得 $x = x_C|_{\Spec(R)}$。对 $C$ 的选取
只是记账手段，可以避免。为后文使用，写成
$C = \Lambda[y_1, \ldots, y_u]/(f_1, \ldots, f_v)$，并把 $y_i$
在映射 $C \to R$ 下的像记为 $\overline{a}_i \in R$。令
$\mathfrak m_C = C \cap \mathfrak m_R$。

\medskip\noindent
选取 $\Lambda$-代数满射 $\Lambda[x_1, \ldots, x_s] \to k$。
% P11-E0228: the frozen source omits "by" in "denote m' the kernel".
把它的核记为 $\mathfrak m'$。由多项式环的泛性质，可将其提升为
$\Lambda$-代数同态 $\Lambda[x_1, \ldots, x_s] \to R$。再添加若干变量
（即略微增大 $s$），使它们映到 $\mathfrak m_R$ 的一组生成元。
完成这一步后，$\Lambda[x_1, \ldots, x_s] \to R/\mathfrak m_R^2$
是满射。因此，
\begin{equation}
\label{artin-equation-surjection}
P = \Lambda[x_1, \ldots, x_s]_{\mathfrak m'}^\wedge \longrightarrow R
\end{equation}
是 Noether 完备局部环的满同态；例如参见《形式形变理论》引理
\ref{formal-defos-lemma-surjective-cotangent-space}。

\medskip\noindent
选取上面得到的 $\overline{a}_i$ 的提升 $a_i \in P$。
选取式 (\ref{artin-equation-surjection}) 之核的一组生成元
$b_1, \ldots, b_r \in P$。选取 $c_{ji} \in P$，使
$$
f_j(a_1, \ldots, a_u) = \sum c_{ji} b_i
$$
在 $P$ 中成立；根据此前的选取，这是可能的。选取如下核的一组生成元
$$
k_1, \ldots, k_t \in
\Ker(P^{\oplus r} \xrightarrow{(b_1, \ldots, b_r)} P)
$$
，并写 $k_i = (k_{i1}, \ldots, k_{ir})$ 和 $K = (k_{ij})$，使
$$
P^{\oplus t} \xrightarrow{K}
P^{\oplus r} \xrightarrow{(b_1, \ldots, b_r)}
P \to R \to 0
$$
成为 $P$-模的正合列。特别地，$\sum k_{ij} b_j = 0$。必要时增大
$N$ 后，可假设 $N - 1$ 对此正合列的前两个映射适用 Artin--Rees
引理（术语参见《更多代数》一节
\ref{more-algebra-section-artin-rees}）。

\medskip\noindent
由假设，$\mathcal{O}_{S, s} = \Lambda_{\Lambda \cap \mathfrak m'}$ 是
G-环。因此，由《更多代数》命题
\ref{more-algebra-proposition-finite-type-over-G-ring}
，环 $\Lambda[x_1, \ldots, x_s]_{\mathfrak m'}$ 是 $G$-环。
因而，由《环同态的光滑化》定理
\ref{smoothing-theorem-approximation-property-variant}
，存在 \'etale 环同态
$$
\Lambda[x_1, \ldots, x_s]_{\mathfrak m'} \to B,
$$
、$B$ 中位于 $\mathfrak m'$ 之上的极大理想 $\mathfrak m_B$，以及
% P11-E0229: the frozen source places the approximating elements in undefined B' although only B has been introduced.
元素 $a'_i, b'_i, c'_{ij}, k'_{ij} \in B'$，使得
\begin{enumerate}
\item $\kappa(\mathfrak m') = \kappa(\mathfrak m_B)$；这蕴含
$\Lambda[x_1, \ldots, x_s]_{\mathfrak m'} \subset B_{\mathfrak m_B}
\subset P$，并且 $P$ 被认同为 $B$ 在 $\mathfrak m_B$ 处的完备化；
参见《环同态的光滑化》定理
\ref{smoothing-theorem-approximation-property-variant} 之前的评注；
\item $a_i - a'_i, b_i - b'_i, c_{ij} - c'_{ij}, k_{ij} - k'_{ij} \in
(\mathfrak m')^N P$；并且
\item $f_j(a'_1, \ldots, a'_u) = \sum c'_{ji} b'_i$ 且
$\sum k'_{ij}b'_j = 0$。
\end{enumerate}
令 $A = B/(b'_1, \ldots, b'_r)$。
% P11-E0230: the frozen source omits "by" in "denote m_A the image".
把 $\mathfrak m_B$ 在 $A$ 中的像记为 $\mathfrak m_A$。（注意，$A$
在 $\Lambda$ 上本质有限型；证明末尾将说明如何得到在 $\Lambda$
上有限型的 $A$。）存在环同态 $C \to A$，把 $y_i \mapsto a'_i$，
因为 $a'_i$ 模 $(b'_1, \ldots, b'_r)$ 满足所需方程。注意，由上面的
性质 (2)，作为 $P = B^\wedge$ 的商，有
$A/\mathfrak m_A^N = R/\mathfrak m_R^N$。令
$x_A = x_C|_{\Spec(A)}$。由于映射
$$
C \to A \to A/\mathfrak m_A^N \cong R/\mathfrak m_R^N
\quad\text{且}\quad
C \to R \to R/\mathfrak m_R^N
$$
相等，经同构 $A/\mathfrak m_A^N = R/\mathfrak m_R^N$ 可知 $x_A$
和 $x$ 模 $\mathfrak m_R^N$ 相符。至此，已经证明引理陈述中的性质
(1)--(5)。为证明 (6)，注意
$$
P^{\oplus t} \xrightarrow{K}
P^{\oplus r} \xrightarrow{(b_1, \ldots, b_r)}
P
\quad\text{且}\quad
P^{\oplus t} \xrightarrow{K'}
P^{\oplus r} \xrightarrow{(b'_1, \ldots, b'_r)}
P
$$
是两个 $P$-模复形，它们模 $(\mathfrak m')^N$ 同余，且第一个正合。
根据上面对 $N$ 的选取，由《更多代数》引理
\ref{more-algebra-lemma-approximate-complex-graded} 可知，
$R = P/(b_1, \ldots, b_r)$ 和
$P/(b'_1, \ldots, b'_r) = B^\wedge/(b'_1, \ldots, b'_r) = A^\wedge$
具有同构的相伴分次代数。这正是所要证明的。

\medskip\noindent
证明的最后一段用来处理 $A$ 在 $S$ 上本质有限型、但尚非有限型的问题。
上面的构造给出 $A = B/(b'_1, \ldots, b'_r)$ 和
$\mathfrak m_A \subset A$，其中 $B$ 在
$\Lambda[x_1, \ldots, x_s]_{\mathfrak m'}$ 上为 \'etale。因此，$A$
在 Noether 环 $\Lambda[x_1, \ldots, x_s]_{\mathfrak m'}$ 上有限型。
% P11-E0231: the frozen source changes the polynomial-variable upper index from s to n without introducing n.
于是，对某个有限型 $\Lambda[x_1, \ldots, x_n]$-代数 $A_0$，可写成
$A = (A_0)_{\mathfrak m'}$。然后
$A = \colim (A_0)_f$，其中
$f \in \Lambda[x_1, \ldots, x_n] \setminus \mathfrak m'$；参见
《代数》引理 \ref{algebra-lemma-localization-colimit}。由于
$p : \mathcal{X} \to (\Sch/S)_{fppf}$ 在对象上保极限，可知对某个
上述 $f$，$x_A$ 来自 $\Spec((A_0)_f)$ 上的某个对象
$x_{(A_0)_f}$。用 $(A_0)_f$ 代替 $A$、用 $x_{(A_0)_f}$ 代替
$x_A$，并用 $(A_0)_f \cap \mathfrak m_A$ 代替 $\mathfrak m_A$ 后，
证明完成。
\end{proof}





\section{保极限性}
\label{artin-section-limits}

\noindent
若下述定义中的函子本质满，则态射
$p : \mathcal{X} \to (\Sch/S)_{fppf}$ 在《可表示性判据》一节
\ref{criteria-section-limit-preserving} 所定义的意义下在对象上保极限。
不过，《例》一节 \ref{examples-section-limit-preserving} 中的例子表明，
这并不等价于保极限。

\begin{definition}
\label{artin-definition-limit-preserving}
设 $S$ 是概形，$\mathcal{X}$ 是 $(\Sch/S)_{fppf}$ 上的
群胚纤维化范畴。若对每个 $S$ 上的仿射概形 $T$，只要它是
$S$ 上仿射概形 $T_i$ 的某个有向逆系统的极限 $T = \lim T_i$，
就有纤维范畴的等价
$$
\colim \mathcal{X}_{T_i} \longrightarrow \mathcal{X}_T
$$
，则称 $\mathcal{X}$ {\it 保极限}。
\end{definition}

\noindent
详细说明其含义。首先，给定 $\mathcal{X}$ 在 $T_i$ 上的对象 $x, y$，
应有
$$
\Mor_{\mathcal{X}_T}(x|_T, y|_T) =
\colim_{i' \geq i} \Mor_{\mathcal{X}_{T_{i'}}}(x|_{T_{i'}}, y|_{T_{i'}})
$$
；其次，$\mathcal{X}_T$ 的每个对象都同构于某个 $T_i$ 上对象的限制，
其中 $i$ 适当选取。注意，第一个条件意味着预层
$\mathit{Isom}_\mathcal{X}(x, y)$ 保极限（参见《叠》定义
\ref{stacks-definition-mor-presheaf}）。

\begin{lemma}
\label{artin-lemma-fibre-product-limit-preserving}
设 $S$ 是概形。设 $p : \mathcal{X} \to \mathcal{Y}$ 和
$q : \mathcal{Z} \to \mathcal{Y}$ 是 $(\Sch/S)_{fppf}$ 上
群胚纤维化范畴的 $1$-态射。
\begin{enumerate}
\item 若 $\mathcal{X} \to (\Sch/S)_{fppf}$ 和
$\mathcal{Z} \to (\Sch/S)_{fppf}$ 在对象上保极限，且
$\mathcal{Y}$ 保极限，则
$\mathcal{X} \times_\mathcal{Y} \mathcal{Z} \to (\Sch/S)_{fppf}$
在对象上保极限。
\item 若 $\mathcal{X}$、$\mathcal{Y}$ 和 $\mathcal{Z}$ 保极限，则
$\mathcal{X} \times_\mathcal{Y} \mathcal{Z}$ 也保极限。
\end{enumerate}
\end{lemma}

\begin{proof}
这是形式论证。证明 (1)。设 $T = \lim_{i \in I} T_i$ 是 $S$ 上
仿射概形 $T_i$ 的有向极限。我们将证明函子
$\colim \mathcal{X}_{T_i} \to \mathcal{X}_T$ 本质满。回忆，纤维积在
$T$ 上的一个对象是四元组 $(T, x, z, \alpha)$，其中 $x$ 是
$\mathcal{X}$ 位于 $T$ 上的对象，$z$ 是 $\mathcal{Z}$ 位于 $T$
上的对象，而 $\alpha : p(x) \to q(z)$ 是 $\mathcal{Y}$ 在 $T$
上的纤维范畴中的态射。由对 $\mathcal{X}$ 和 $\mathcal{Z}$ 的假设，
可以找到 $i$ 以及 $T_i$ 上的对象 $x_i$ 和 $z_i$，使得
% P11-E0232: the frozen source writes x_i|_T congruent T where the intended object is x.
$x_i|_T \cong T$ 且 $z_i|_T \cong z$。于是，$\alpha$ 对应于同构
$p(x_i)|_T \to q(z_i)|_T$。由对 $\mathcal{Y}$ 的假设，此同构来自同构
$\alpha_{i'} : p(x_i)|_{T_{i'}} \to q(z_i)|_{T_{i'}}$。用 $i'$ 代替
$i$、用 $x_i|_{T_{i'}}$ 代替 $x_i$，并用 $z_i|_{T_{i'}}$ 代替
$z_i$ 后，可知 $(T_i, x_i, z_i, \alpha_i)$ 是纤维积在 $T_i$
上的一个对象，而它限制为 $T$ 上同构于 $(T, x, z, \alpha)$ 的对象，
如所要求。

\medskip\noindent
% P11-E0233: the frozen proof of part (2) names the X-fibre functor rather than the fibre-product functor whose full faithfulness is at issue.
在 (2) 中，证明 $\colim \mathcal{X}_{T_i} \to \mathcal{X}_T$
全忠实的论证从略。
\end{proof}

\begin{lemma}
\label{artin-lemma-limit-preserving-algebraic-space}
设 $S$ 是概形，$\mathcal{X}$ 是 $S$ 上的代数叠。则下列条件等价：
\begin{enumerate}
\item $\mathcal{X}$ 是广群叠，且
$\mathcal{X} \to (\Sch/S)_{fppf}$ 在对象上保极限；
\item $\mathcal{X}$ 是广群叠且保极限；
\item $\mathcal{X}$ 可由一个局部有限表示的代数空间表示。
\end{enumerate}
\end{lemma}

\begin{proof}
在三个条件中的每一个之下，$\mathcal{X}$ 都可由 $S$ 上的代数空间
$X$ 表示；参见《代数叠》命题
\ref{algebraic-proposition-algebraic-stack-no-automorphisms}。
(1) 与 (2) 显然等价，因为广群之间的函子是等价，当且仅当它在同构类上
为满射。最后，由《空间的极限》命题
\ref{spaces-limits-proposition-characterize-locally-finite-presentation}，
(1) 与 (3) 等价。
\end{proof}

\begin{lemma}
\label{artin-lemma-diagonal}
设 $S$ 是概形，$\mathcal{X}$ 是 $(\Sch/S)_{fppf}$ 上的
群胚纤维化范畴。假设
$\Delta : \mathcal{X} \to \mathcal{X} \times \mathcal{X}$ 可由代数空间
表示，且 $\mathcal{X}$ 保极限。则 $\Delta$ 局部有限型。
\end{lemma}

\begin{proof}
应用《可表示性判据》引理
\ref{criteria-lemma-check-property-limit-preserving}。
% P11-E0234: the frozen source redundantly says "Let V be an affine scheme V".
设 $V$ 是在 $S$ 上局部有限表示的仿射概形 $V$，并设 $\theta$ 是
$\mathcal{X} \times \mathcal{X}$ 在 $V$ 上的对象。设 $F_\theta$ 为表示
$\mathcal{X} \times_{\Delta, \mathcal{X} \times \mathcal{X}, \theta}
(\Sch/V)_{fppf}$ 的代数空间，并设 $f_\theta : F_\theta \to V$ 为
典范态射（参见《代数叠》一节
\ref{algebraic-section-morphisms-representable-by-algebraic-spaces}）。
只需证明 $F_\theta \to V$ 具有相应性质。由引理
\ref{artin-lemma-fibre-product-limit-preserving} 和
\ref{artin-lemma-limit-preserving-algebraic-space}，$F_\theta \to S$
局部有限表示。因此，由《空间的态射》引理
\ref{spaces-morphisms-lemma-permanence-finite-type}，
$F_\theta \to V$ 局部有限型。
\end{proof}






\section{通用性}
\label{artin-section-versality}

\noindent
上一节说明了如何用代数对象逼近完备局部环上的对象。不过，为证明叠
$\mathcal{X}$ 是代数叠，需要找到从概形到 $\mathcal{X}$ 的光滑
$1$-态射。由于不预先假设 $\mathcal{X}$ 具有可表对角态射，甚至无法谈论
到 $\mathcal{X}$ 的光滑态射。因此，我们借用形变理论的术语，引入通用对象。

\begin{definition}
\label{artin-definition-versal-formal-object}
设 $S$ 是局部 Noether 概形，并设
$p : \mathcal{X} \to (\Sch/S)_{fppf}$ 是群胚纤维化范畴。设
$\xi = (R, \xi_n, f_n)$ 是形式对象。令 $k = R/\mathfrak m$ 且
$x_0 = \xi_1$。若把 $\xi$ 视为 $\mathcal{F}_{\mathcal{X}, k, x_0}$
的形式对象（评注 \ref{artin-remark-formal-objects-match}）时，它在
《形式形变理论》定义 \ref{formal-defos-definition-versal} 的意义下通用，
则称 $\xi$ 是{\it 通用的}。
\end{definition}

\noindent
简要说明其含义。采用定义中的记号，假设给定 $\mathcal{X}$ 的态射
$\xi_1 = x_0 \to y \to z$，它们位于闭浸入
$\Spec(k) \to \Spec(A) \to \Spec(B)$
之上，其中 $A, B$ 是剩余域为 $k$ 的 Artin 局部环。再假设给定
$n \geq 1$ 和交换图
$$
\vcenter{
\xymatrix{
& y \ar[ld] \\
\xi_n & \xi_1 \ar[u] \ar[l]
}
}
\quad\text{位于下列对象之上：}\quad
\vcenter{
\xymatrix{
& \Spec(A) \ar[ld] \\
\Spec(R/\mathfrak m^n) & \Spec(k) \ar[u] \ar[l]
}
}
$$
通用性意味着，对任意上述数据，都存在 $m \geq n$ 和交换图
$$
\vcenter{
\xymatrix{
& & z \ar[lldd] \\
& & y \ar[ld] \ar[u] \\
\xi_m & \xi_n \ar[l] & \xi_1 \ar[u] \ar[l]
}
}
\quad\text{位于下列对象之上：}
\vcenter{
\xymatrix{
& & \Spec(B) \ar[lldd] \\
& & \Spec(A) \ar[ld] \ar[u] \\
\Spec(R/\mathfrak m^m) &
\Spec(R/\mathfrak m^n) \ar[l] &
\Spec(k) \ar[u] \ar[l]
}
}
$$
请与《形式形变理论》评注 \ref{formal-defos-remark-versal-object} 比较。

\medskip\noindent
设 $S$ 是局部 Noether 概形，$U$ 是 $S$ 上的概形，其结构态射
$U \to S$ 局部有限型。设 $u_0 \in U$ 是 $U$ 的有限型点；参见
《态射》定义 \ref{morphisms-definition-finite-type-point}。令
$k = \kappa(u_0)$。注意，复合 $\Spec(k) \to S$ 也是有限型的；
参见《态射》引理 \ref{morphisms-lemma-composition-finite-type}。设
$p : \mathcal{X} \to (\Sch/S)_{fppf}$ 是群胚纤维化范畴。设 $x$
是 $\mathcal{X}$ 位于 $U$ 上的对象。
% P11-E0235: the frozen source omits "by" in "Denote x_0 the pullback".
把 $x$ 由 $u_0$ 的拉回记为 $x_0$。由 $2$-Yoneda 引理，$x$ 对应于
$1$-态射
$$
x : (\Sch/U)_{fppf} \longrightarrow \mathcal{X},
$$
；参见《代数叠》一节 \ref{algebraic-section-2-yoneda}。由此得到
预形变范畴的态射
\begin{equation}
\label{artin-equation-hat-x}
\hat x :
\mathcal{F}_{(\Sch/U)_{fppf}, k, u_0}
\longrightarrow
\mathcal{F}_{\mathcal{X}, k, x_0},
\end{equation}
，它位于 $\mathcal{C}_\Lambda$ 之上；参见
(\ref{artin-equation-functoriality})。

\begin{definition}
\label{artin-definition-versal}
设 $S$ 是局部 Noether 概形，$\mathcal{X}$ 是
$(\Sch/S)_{fppf}$ 上的群胚纤维化范畴。设 $U$ 是 $S$ 上局部有限型的
概形，$x$ 是 $\mathcal{X}$ 位于 $U$ 上的对象，而 $u_0$ 是 $U$
的有限型点。若态射 $\hat x$（式 (\ref{artin-equation-hat-x})）光滑，
则称 $x$ 在 $u_0$ 处{\it 通用}；参见《形式形变理论》定义
\ref{formal-defos-definition-smooth-morphism}。
\end{definition}

\noindent
这个定义与 $\mathcal{X}$ 的形式对象的通用性概念相吻合。

\begin{lemma}
\label{artin-lemma-versality-matches}
% P11-E0236: the frozen source ends the introductory "With notation" phrase with a period, leaving a sentence fragment.
采用定义 \ref{artin-definition-versal} 中的记号。令
$R = \mathcal{O}_{U, u_0}^\wedge$。令 $\xi$ 为与
$x|_{\Spec(R)}$ 相联系的 $\mathcal{X}$ 在 $R$ 上的形式对象；
参见 (\ref{artin-equation-approximation})。则
$$
x\text{ is versal at }u_0
\Leftrightarrow
\xi\text{ is versal}
$$
\end{lemma}

\begin{proof}
注意，$\mathcal{O}_{U, u_0}$ 是剩余域为 $k$ 的 Noether 局部
$S$-代数。因此，$R = \mathcal{O}_{U, u_0}^\wedge$ 是
$\mathcal{C}_\Lambda^\wedge$ 的对象；参见《形式形变理论》定义
\ref{formal-defos-definition-completion-CLambda}。回忆，$\xi$ 通用是指
$\underline{\xi} : \underline{R}|_{\mathcal{C}_\Lambda} \to
\mathcal{F}_{\mathcal{X}, k, x_0}$
光滑，而 $x$ 在 $u_0$ 处通用是指
$\hat x : \mathcal{F}_{(\Sch/U)_{fppf}, k, u_0}
\to \mathcal{F}_{\mathcal{X}, k, x_0}$ 光滑。存在预形变范畴的认同
$$
\underline{R}|_{\mathcal{C}_\Lambda}
=
\mathcal{F}_{(\Sch/U)_{fppf}, k, u_0},
$$
；记号参见《形式形变理论》评注
\ref{formal-defos-remark-formal-objects-yoneda}。具体地，给定 Artin
局部 $S$-代数 $A$，并把它的剩余域与 $k$ 认同，就有
$$
\Mor_{\mathcal{C}_\Lambda^\wedge}(R, A) =
\{\varphi \in \Mor_S(\Spec(A), U) \mid \varphi|_{\Spec(k)} = u_0\}
$$
。展开定义即可验证，所得映射
$$
\underline{R}|_{\mathcal{C}_\Lambda} =
\mathcal{F}_{(\Sch/U)_{fppf}, k, u_0}
\xrightarrow{\hat x}
\mathcal{F}_{\mathcal{X}, k, x_0},
$$
等于 $\underline{\xi}$，故引理成立。
\end{proof}

\noindent
下面给出一个合理性检验。

\begin{lemma}
\label{artin-lemma-versal-implies-smooth}
设 $S$ 是局部 Noether 概形。设 $f : U \to V$ 是在 $S$ 上局部有限型的
概形态射，并设 $u_0 \in U$ 是有限型点。下列条件等价：
\begin{enumerate}
\item $f$ 在 $u_0$ 处光滑；
\item 把 $f$ 视为 $(\Sch/V)_{fppf}$ 位于 $U$ 上的对象时，
它在 $u_0$ 处通用。
\end{enumerate}
\end{lemma}

\begin{proof}
这只是《更多态射》引理
\ref{more-morphisms-lemma-lifting-along-artinian-at-point} 的重述。
\end{proof}

\noindent
这个概念对于域扩张具有良好行为。

\begin{lemma}
\label{artin-lemma-versal-change-of-field}
设 $S$、$\mathcal{X}$、$U$、$x$、$u_0$ 如定义
\ref{artin-definition-versal}。设 $l$ 是域，并设
$u_{l, 0} : \Spec(l) \to U$ 是像为 $u_0$ 的态射，且
% P11-E0237: the frozen source equates the extension l/k with the residue field kappa(u_0), rather than defining k as that residue field and requiring l/k finite.
$l/k = \kappa(u_0)$ 是有限的。令 $x_{l, 0} = x_0|_{\Spec(l)}$。
若 $\mathcal{X}$ 满足 (RS)，且 $x$ 在 $u_0$ 处通用，则
$$
\mathcal{F}_{(\Sch/U)_{fppf}, l, u_{l, 0}}
\longrightarrow
\mathcal{F}_{\mathcal{X}, l, x_{l, 0}}
$$
光滑。
\end{lemma}

\begin{proof}
注意，由引理 \ref{artin-lemma-algebraic-stack-RS}，
$(\Sch/U)_{fppf}$ 满足 (RS)。因此，引理中的函子就是与 $\hat x$
相联系的函子
$$
(\mathcal{F}_{(\Sch/U)_{fppf}, k , u_0})_{l/k}
\longrightarrow
(\mathcal{F}_{\mathcal{X}, k , x_0})_{l/k}
$$
；参见引理 \ref{artin-lemma-change-of-field}。因而，引理由《形式形变理论》引理
\ref{formal-defos-lemma-change-of-fields-smooth} 得到。
\end{proof}

\noindent
下述引理是另一个合理性检验。其含义大致是：若 $x$ 如定义
\ref{artin-definition-versal} 那样在 $u_0$ 处通用，则把 $x$ 视为从 $U$
到 $\mathcal{X}$ 的态射后，每当用概形作基变换时，它都是光滑的。

\begin{lemma}
\label{artin-lemma-base-change-versal}
设 $S$、$\mathcal{X}$、$U$、$x$、$u_0$ 如定义
\ref{artin-definition-versal}。假设
\begin{enumerate}
\item $\Delta : \mathcal{X} \to \mathcal{X} \times \mathcal{X}$
可由代数空间表示；
\item $\Delta$ 局部有限型（例如，当 $\mathcal{X}$ 保极限时）；并且
\item $\mathcal{X}$ 具有 (RS)。
\end{enumerate}
设 $V$ 是 $S$ 上局部有限型的概形，而 $y$ 是 $\mathcal{X}$ 在 $V$
上的对象。作 $2$-纤维积
$$
\xymatrix{
\mathcal{Z} \ar[r] \ar[d] & (\Sch/U)_{fppf} \ar[d]^x \\
(\Sch/V)_{fppf} \ar[r]^y & \mathcal{X}
}
$$
设 $Z$ 是表示 $\mathcal{Z}$ 的代数空间，并设 $z_0 \in |Z|$ 是位于
$u_0$ 之上的有限型点。若 $x$ 在 $u_0$ 处通用，则态射
$Z \to V$ 在 $z_0$ 处光滑。
\end{lemma}

\begin{proof}
（陈述中的括号评注由引理 \ref{artin-lemma-diagonal} 得到。）由假设 (1)
和《代数叠》引理 \ref{algebraic-lemma-representable-diagonal}，$Z$ 存在。
由假设 (2)，$Z \to V \times_S U$ 局部有限型。选取概形 $W$、闭点
$w_0 \in W$，以及把 $w_0$ 映到 $z_0$ 的 \'etale 态射 $W \to Z$；
参见《空间的态射》定义
\ref{spaces-morphisms-definition-finite-type-point}。于是 $W$ 在 $S$
上局部有限型，且 $w_0$ 是 $W$ 的有限型点。令 $l = \kappa(z_0)$。
% P11-E0238: the frozen source omits "by" in the four-object naming construction.
把通过拉回到 $\Spec(l) = w_0$ 所得的 $\mathcal{Z}$、
$(\Sch/V)_{fppf}$、$(\Sch/U)_{fppf}$ 和 $\mathcal{X}$ 在 $\Spec(l)$
上的对象分别记为 $z_{l, 0}$、$v_{l, 0}$、$u_{l, 0}$ 和
$x_{l, 0}$。考虑
$$
\xymatrix{
\mathcal{F}_{(\Sch/W)_{fppf}, l, w_0} \ar[r] &
\mathcal{F}_{\mathcal{Z}, l, z_{l, 0}} \ar[d] \ar[r] &
\mathcal{F}_{(\Sch/U)_{fppf}, l, u_{l, 0}} \ar[d] \\
& \mathcal{F}_{(\Sch/V)_{fppf}, l, v_{l, 0}} \ar[r] &
\mathcal{F}_{\mathcal{X}, l, x_{l, 0}}
}
$$
由引理 \ref{artin-lemma-fibre-product-deformation-categories}，此方块是
预形变范畴的纤维积。由引理 \ref{artin-lemma-versal-change-of-field}，
右边的竖箭头光滑。由《形式形变理论》引理
\ref{formal-defos-lemma-smooth-properties}，左边的竖箭头光滑。
由引理 \ref{artin-lemma-formally-smooth-on-deformation-categories}，
左边的横箭头光滑。由此断定映射
$$
\mathcal{F}_{(\Sch/W)_{fppf}, l, w_0} \to
\mathcal{F}_{(\Sch/V)_{fppf}, l, v_{l, 0}}
$$
由《形式形变理论》引理 \ref{formal-defos-lemma-smooth-properties}
是光滑的。因此，由《更多态射》引理
\ref{more-morphisms-lemma-lifting-along-artinian-at-point}，
$W \to V$ 在 $w_0$ 处光滑。这恰好意味着 $Z \to V$ 在 $z_0$
处光滑，证明完成。
\end{proof}

\noindent
用通用对象的语言重述逼近结果。

\begin{lemma}
\label{artin-lemma-approximate-versal}
设 $S$ 是局部 Noether 概形，并设
$p : \mathcal{X} \to (\Sch/S)_{fppf}$ 是群胚纤维化范畴。设
$\xi = (R, \xi_n, f_n)$ 是 $\mathcal{X}$ 的形式对象，其中 $\xi_1$
位于 $\Spec(k) \to S$ 之上，而该态射的像为 $s \in S$。假设
\begin{enumerate}
\item $\xi$ 通用；
\item $\xi$ 有效；
\item $\mathcal{O}_{S, s}$ 是 G-环；并且
\item $p : \mathcal{X} \to (\Sch/S)_{fppf}$ 在对象上保极限。
\end{enumerate}
则存在有限型态射 $U \to S$、剩余域为 $k$ 的有限型点
$u_0 \in U$，以及 $\mathcal{X}$ 在 $U$ 上的对象 $x$，使得 $x$
在 $u_0$ 处通用，并且
$x|_{\Spec(\mathcal{O}_{U, u_0}/\mathfrak m_{u_0}^n)} \cong \xi_n$。
\end{lemma}

\begin{proof}
选取 $\mathcal{X}$ 位于 $\Spec(R)$ 上的对象 $x_R$，使其相联系的形式对象
为 $\xi$。令 $N = 2$，并应用引理 \ref{artin-lemma-approximate}，得到
$A, \mathfrak m_A, x_A, \ldots$。令
$\eta = (A^\wedge, \eta_n, g_n)$ 为与
$x_A|_{\Spec(A^\wedge)}$ 相联系的形式对象。存在图
$$
\vcenter{
\xymatrix{
& \eta \ar[d] \\
\xi \ar[r] \ar@{..>}[ru] & \xi_2 = \eta_2
}
}
\quad\text{位于下列对象之上：}\quad
\vcenter{
\xymatrix{
& A^\wedge \ar[d] \\
R \ar[r] \ar@{..>}[ru] & R/\mathfrak m_R^2 = A/\mathfrak m_A^2
}
}
$$
由 $\xi$ 的通用性，恰好可以找到诸图中的虚线箭头，因为由
《形式形变理论》评注 \ref{formal-defos-remark-versal-object}，
可以逐次找到态射 $\xi \to \eta_3$、$\xi \to \eta_4$ 等等。
由《形式形变理论》引理
\ref{formal-defos-lemma-surjective-cotangent-space}，相应环同态
$R \to A^\wedge$ 是满射。另一方面，由构造，对所有 $n$ 都有
$\dim_k \mathfrak m_R^n/\mathfrak m_R^{n + 1} =
\dim_k \mathfrak m_A^n/\mathfrak m_A^{n + 1}$。因此，由长度的可加性和
《形式形变理论》引理 \ref{formal-defos-lemma-length}，
$R/\mathfrak m_R^n$ 和 $A/\mathfrak m_A^n$ 作为 $\Lambda$-模具有相同的
有限长度。于是，对所有 $n$，
$R/\mathfrak m_R^n \to A/\mathfrak m_A^n$ 都是同构，从而
$R \to A^\wedge$ 是同构。因此，$\eta$ 同构于通用对象，故其自身也通用。
由引理 \ref{artin-lemma-versality-matches}，$x_A$ 在
$U = \Spec(A)$ 中与 $\mathfrak m_A$ 对应的点 $u_0$ 处通用。
\end{proof}

\begin{example}
\label{artin-example-approximate-versal-implies}
本例说明，要使引理 \ref{artin-lemma-approximate-versal} 的结论成立，局部环
$\mathcal{O}_{S, s}$ 必须是 G-环。具体地，设 $\Lambda$ 是 Noether 环，
$\mathfrak m$ 是 $\Lambda$ 的极大理想。令
$R = \Lambda_\mathfrak m^\wedge$。设 $\Lambda \to C \to R$ 是一个分解，
其中 $C$ 在 $\Lambda$ 上有限型。令 $S = \Spec(\Lambda)$、
$U = S \setminus \{\mathfrak m\}$，以及 $S' = U \amalg \Spec(C)$。
考虑由下述规则定义的函子
$F : (\Sch/S)_{fppf}^{opp} \to \textit{Sets}$：
$$
F(T) = 
\left\{
\begin{matrix}
* & \text{若 }T \to S\text{ factors through }S' \\
\emptyset & \text{否则}
\end{matrix}
\right.
$$
% P11-E0239: the frozen source says "Let X = S_F is the category", combining incompatible let- and copular constructions.
令 $\mathcal{X} = \mathcal{S}_F$ 是与 $F$ 相联系的集合纤维化范畴；
参见《代数叠》一节 \ref{algebraic-section-split}。于是
$\mathcal{X} \to (\Sch/S)_{fppf}$ 在对象上保极限，并且在 $R$ 上存在
一个有效的通用形式对象 $\xi$。因此，若引理
\ref{artin-lemma-approximate-versal} 的结论对 $\mathcal{X}$ 成立，则存在
有限型环同态 $\Lambda \to A$ 和位于 $\mathfrak m$ 之上的极大理想
$\mathfrak m_A$，使得
\begin{enumerate}
\item $\kappa(\mathfrak m) = \kappa(\mathfrak m_A)$；
\item $\Lambda \to A$ 和 $\mathfrak m_A$ 满足《代数》引理
\ref{algebra-lemma-smooth-test-artinian} 的条件 (4)；并且
\item 存在 $\Lambda$-代数同态 $C \to A$。
\end{enumerate}
由所引引理，$\Lambda \to A$ 在 $\mathfrak m_A$ 处光滑。对 $A$ 作切片后，
可假设 $\Lambda \to A$ 在 $\mathfrak m_A$ 处为 \'etale；例如参见
《更多态射》引理 \ref{more-morphisms-lemma-slice-smooth}，也可以直接论证。
写 $C = \Lambda[y_1, \ldots, y_n]/(f_1, \ldots, f_m)$。于是，
$C \to R$ 对应于方程组 $f_1 = \ldots = f_m = 0$ 在 $R$ 中的一个解；
参见《环同态的光滑化》一节
\ref{smoothing-section-approximation-G-rings}。因此，若引理
\ref{artin-lemma-approximate-versal} 的结论对每个上述 $\mathcal{X}$ 成立，
则在 $R$ 中有解的方程组在 $\Lambda_{\mathfrak m}$ 的 Hensel 化中也有解。
换言之，逼近性质对 $\Lambda_{\mathfrak m}^h$ 成立。这蕴含
$\Lambda_{\mathfrak m}^h$ 是 G-环
% P11-E0240: the frozen source contains a second unresolved future-reference placeholder.
（此处补入将来的参考文献；另见《环同态的光滑化》一节
\ref{smoothing-section-introduction} 中的讨论），进而蕴含
$\Lambda_{\mathfrak m}$ 是 G-环。
\end{example}






\section{通用性的开性}
\label{artin-section-openness-versality}

\noindent
接下来讨论通用性的开性。

\begin{definition}
\label{artin-definition-openness-versality}
设 $S$ 是局部 Noether 概形。
\begin{enumerate}
\item 设 $\mathcal{X}$ 是 $(\Sch/S)_{fppf}$ 上的群胚纤维化范畴。
若给定 $S$ 上局部有限型的概形 $U$、$\mathcal{X}$ 位于 $U$ 上的对象
$x$ 和有限型点 $u_0 \in U$，且 $x$ 在 $u_0$ 处通用时，总存在开邻域
$u_0 \in U' \subset U$，使得 $x$ 在 $U'$ 的每个有限型点处都通用，
则称 $\mathcal{X}$ 满足{\it 通用性的开性}。
\item 设 $f : \mathcal{X} \to \mathcal{Y}$ 是 $(\Sch/S)_{fppf}$ 上
群胚纤维化范畴的 $1$-态射。若给定 $S$ 上局部有限型的概形 $U$ 和
$\mathcal{Y}$ 位于 $U$ 上的对象 $y$ 时，
$(\Sch/U)_{fppf} \times_\mathcal{Y} \mathcal{X}$ 满足通用性的开性，
则称 $f$ 满足{\it 通用性的开性}。
\end{enumerate}
\end{definition}

\noindent
通用性的开性往往是最难核验的条件。不过，下面的例子表明这项要求确属必要。

\begin{example}
\label{artin-example-versality}
设 $k$ 是域，并令 $\Lambda = k[s, t]$。考虑依下述规则定义的函子
$F : \Lambda\text{-algebras} \longrightarrow \textit{Sets}$
$$
F(A) =
\left\{
\begin{matrix}
* & \text{if there exist }f_1, \ldots, f_n \in A\text{ 使得 } \\
  & A = (s, t, f_1, \ldots, f_n)\text{ 且 } f_i s = 0\ \forall i \\
\emptyset & \text{否则}
\end{matrix}
\right.
$$
从几何上说，$F(A) = *$ 意味着 $V(s, t) \subset \Spec(A)$ 有一个
拟紧开邻域 $W$，使得 $s|_W = 0$。令
$\mathcal{X} \subset (\Sch/\Spec(\Lambda))_{fppf}$ 为如下概形 $T$
所成的全子范畴：它有仿射开覆盖 $T = \bigcup \Spec(A_j)$，且对每个
$j$ 都有 $F(A_j) = *$。于是 $\mathcal{X}$ 满足 [0]、[1]、[2]、[3]
和 [4]，但不满足 [5]。具体地，在
$U = \Spec(k[s, t]/(s))$ 上存在对象 $x$，它在 $u_0 = (s, t)$ 处通用，
但不在任何其他点处通用。细节从略。
\end{example}

\noindent
设 $S$ 是局部 Noether 概形。设
$f : \mathcal{X} \to \mathcal{Y}$ 是 $(\Sch/S)_{fppf}$ 上群胚纤维化
范畴的 $1$-态射。考虑如下性质
\begin{equation}
\label{artin-equation-smooth}
\begin{matrix}
\text{for all fields }k\text{ of finite type over }S
\text{ and all }x_0 \in \Ob(\mathcal{X}_{\Spec(k)})\text{ the}\\
\text{map }
\mathcal{F}_{\mathcal{X}, k, x_0} \to \mathcal{F}_{\mathcal{Y}, k, f(x_0)}
\text{ of predeformation categories is smooth}
\end{matrix}
\end{equation}
下面围绕这一概念表述若干引理。首先把它与通用性（及其开性）联系起来。

\begin{lemma}
\label{artin-lemma-versal-smooth}
设 $S$ 是局部 Noether 概形，$\mathcal{X}$ 是 $(\Sch/S)_{fppf}$ 上的
群胚纤维化范畴。设 $U$ 是 $S$ 上局部有限型的概形，$x$ 是
$\mathcal{X}$ 位于 $U$ 上的对象。假设 $x$ 在 $U$ 的每个有限型点处
都通用，且 $\mathcal{X}$ 满足 (RS)。则
$x : (\Sch/U)_{fppf} \to \mathcal{X}$ 满足 (\ref{artin-equation-smooth})。
\end{lemma}

\begin{proof}
设 $\Spec(l) \to U$ 是态射，且 $l$ 在 $S$ 上有限型。则其像
$u_0 \in U$ 是 $U$ 的有限型点，而 $l/\kappa(u_0)$ 是有限扩张；参见
《态射》一节 \ref{morphisms-section-points-finite-type} 中的讨论。因此，
$\mathcal{F}_{(\Sch/U)_{fppf}, l, u_{l, 0}} \to
\mathcal{F}_{\mathcal{X}, l, x_{l, 0}}$
由引理 \ref{artin-lemma-versal-change-of-field} 可知是光滑的。
\end{proof}

\begin{lemma}
\label{artin-lemma-composition-smooth}
设 $S$ 是局部 Noether 概形。设 $f : \mathcal{X} \to \mathcal{Y}$ 和
$g : \mathcal{Y} \to \mathcal{Z}$ 是 $(\Sch/S)_{fppf}$ 上群胚纤维化
范畴的可复合 $1$-态射。若 $f$ 和 $g$ 满足
(\ref{artin-equation-smooth})，则 $g \circ f$ 也满足它。
\end{lemma}

\begin{proof}
这由《形式形变理论》引理
\ref{formal-defos-lemma-smooth-properties} 形式地推出。
\end{proof}

\begin{lemma}
\label{artin-lemma-base-change-smooth}
设 $S$ 是局部 Noether 概形。设 $f : \mathcal{X} \to \mathcal{Y}$ 和
$\mathcal{Z} \to \mathcal{Y}$ 是 $(\Sch/S)_{fppf}$ 上群胚纤维化范畴的
$1$-态射。若 $f$ 满足 (\ref{artin-equation-smooth})，则投影
$\mathcal{X} \times_\mathcal{Y} \mathcal{Z} \to \mathcal{Z}$ 也满足它。
\end{lemma}

\begin{proof}
% P11-E0241: the frozen source proof begins with the sentence fragment "Follows immediately from".
这立即由引理 \ref{artin-lemma-fibre-product-deformation-categories} 和
《形式形变理论》引理
\ref{formal-defos-lemma-smooth-properties} 得到。
\end{proof}

\begin{lemma}
\label{artin-lemma-smooth-smooth}
% P11-E0242: the frozen source uses plural "morphisms" after the singular article and subject.
设 $S$ 是局部 Noether 概形。设 $f : \mathcal{X} \to \mathcal{Y}$ 是
$(\Sch/S)_{fppf}$ 上群胚纤维化范畴的 $1$-态射。若 $f$ 在对象上形式
光滑，则 $f$ 满足 (\ref{artin-equation-smooth})。若 $f$ 可由代数空间表示且
光滑，则 $f$ 满足 (\ref{artin-equation-smooth})。
\end{lemma}

\begin{proof}
% P11-E0243: the frozen source proof is a noun phrase rather than a complete sentence.
这就是引理 \ref{artin-lemma-formally-smooth-on-deformation-categories} 的改述。
\end{proof}

\begin{lemma}
\label{artin-lemma-implies-smooth}
设 $S$ 是局部 Noether 概形。设 $f : \mathcal{X} \to \mathcal{Y}$ 是
$(\Sch/S)_{fppf}$ 上群胚纤维化范畴的 $1$-态射。假设
\begin{enumerate}
\item $f$ 可由代数空间表示；
\item $f$ 满足 (\ref{artin-equation-smooth})；
\item $\mathcal{X} \to (\Sch/S)_{fppf}$ 在对象上保极限；并且
\item $\mathcal{Y}$ 保极限。
\end{enumerate}
则 $f$ 光滑。
\end{lemma}

\begin{proof}
证明的关键是《更多态射》引理
\ref{more-morphisms-lemma-lifting-along-artinian-at-point}；它（几乎）断言，
$S$ 上有限型概形之间满足 (\ref{artin-equation-smooth}) 的态射是光滑态射。
证明中的其他论证本质上只是组织记号。

\medskip\noindent
设 $V$ 是 $S$ 上的概形，而 $y$ 是 $\mathcal{Y}$ 位于 $V$ 上的对象。
% P11-E0244: the frozen 2-fibre-product subscript names X, although f and y both map to Y.
设 $Z$ 是表示 $2$-纤维积
$\mathcal{Z} = \mathcal{X} \times_{f, \mathcal{X}, y} (\Sch/V)_{fppf}$.
的代数空间。需要证明投影态射 $Z \to V$ 光滑；参见《代数叠》定义
\ref{algebraic-definition-relative-representable-property}。事实上，只须在
$V$ 是 $S$ 上局部有限表示的仿射概形时证明这一点；参见《可表性判据》引理
\ref{criteria-lemma-check-property-limit-preserving}。此时，由引理
\ref{artin-lemma-limit-preserving-algebraic-space}，$(\Sch/V)_{fppf}$ 保极限。
再由引理 \ref{artin-lemma-fibre-product-limit-preserving} 和
\ref{artin-lemma-limit-preserving-algebraic-space}，$Z \to S$ 局部有限表示。
选取概形 $W$ 和满 \'etale 态射 $W \to Z$。于是 $W$ 在 $S$ 上局部有限表示。

\medskip\noindent
由于 $f$ 满足 (\ref{artin-equation-smooth})，由引理
\ref{artin-lemma-base-change-smooth}，$\mathcal{Z} \to (\Sch/V)_{fppf}$ 也满足它。
其次，$(\Sch/W)_{fppf} \to \mathcal{Z}$ 满足
(\ref{artin-equation-smooth})；这是引理 \ref{artin-lemma-smooth-smooth} 的结论。因此，复合
$$
(\Sch/W)_{fppf} \to \mathcal{Z} \to (\Sch/V)_{fppf}
$$
满足 (\ref{artin-equation-smooth})；这由引理 \ref{artin-lemma-composition-smooth} 得到。
《更多态射》引理
\ref{more-morphisms-lemma-lifting-along-artinian-at-point} 表明，复合
$W \to Z \to V$ 在 $W$ 的每个有限型点 $w_0$ 处光滑。由于光滑点集是开的，
由《态射》引理 \ref{morphisms-lemma-enough-finite-type-points}，可知
$W \to V$ 是概形的光滑态射。因此按定义，$Z \to V$ 是代数空间的光滑态射。
\end{proof}

\noindent
下面的引理说明我们将如何使用通用性的开性。

\begin{lemma}
\label{artin-lemma-get-smooth}
设 $S$ 是局部 Noether 概形。设
$p : \mathcal{X} \to (\Sch/S)_{fppf}$ 是群胚纤维化范畴。设 $k$ 是
$S$ 上的有限型域，而 $x_0$ 是 $\mathcal{X}$ 位于 $\Spec(k)$ 上的对象，
其像为 $s \in S$。假设
\begin{enumerate}
\item $\Delta : \mathcal{X} \to \mathcal{X} \times \mathcal{X}$
可由代数空间表示；
\item $\mathcal{X}$ 满足公理 [1]、[2]、[3]（参见一节
\ref{artin-section-axioms}）；
\item $\mathcal{X}$ 的每个形式对象都有效；
\item $\mathcal{X}$ 满足通用性的开性；并且
\item $\mathcal{O}_{S, s}$ 是 G-环。
\end{enumerate}
则存在有限型态射 $U \to S$ 和 $\mathcal{X}$ 位于 $U$ 上的对象 $x$，使得
$$
x : (\Sch/U)_{fppf} \longrightarrow \mathcal{X}
$$
光滑，并且存在有限型点 $u_0 \in U$，其剩余域为 $k$，且
$x|_{u_0} \cong x_0$。
\end{lemma}

\begin{proof}
由公理 [2]、引理 \ref{artin-lemma-deformation-category} 和评注
\ref{artin-remark-deformation-category-implies}，可知
$\mathcal{F}_{\mathcal{X}, k, x_0}$ 满足 (S1) 和 (S2)。再由公理 [3]，
其切空间是有限维的，故《形式形变理论》引理
\ref{formal-defos-lemma-versal-object-existence} 表明
$\mathcal{F}_{\mathcal{X}, k, x_0}$ 有通用形式对象 $\xi$。假设 (3)
断言 $\xi$ 有效。由公理 [1] 和引理 \ref{artin-lemma-approximate-versal}，
存在有限型态射 $U \to S$、$\mathcal{X}$ 位于 $U$ 上的对象 $x$，以及
$U$ 中剩余域为 $k$ 的有限型点 $u_0$，使得 $x$ 在 $u_0$ 处通用，且
$x|_{\Spec(k)} \cong x_0$。由通用性的开性，可以缩小 $U$，并假设 $x$
在 $U$ 的每个有限型点处都通用。我们断言
$$
x : (\Sch/U)_{fppf} \longrightarrow \mathcal{X}
$$
光滑，这就证明了引理。事实上，由引理 \ref{artin-lemma-versal-smooth}，$x$
满足 (\ref{artin-equation-smooth})；再由引理 \ref{artin-lemma-implies-smooth} 即得结论。
\end{proof}






\section{公理}
\label{artin-section-axioms}

\noindent
设 $S$ 是局部 Noether 概形。设
$p : \mathcal{X} \to (\Sch/S)_{fppf}$ 是群胚纤维化范畴。下面列出我们将对
$\mathcal{X}$ 考察的公理。
\begin{enumerate}
% P11-E0245: the frozen source says "is <= than", redundantly combining the inequality with "than".
\item[{[-1]}] 一个集合论条件\footnote{该条件如下：在 $k$ 遍历 $S$ 上
有限型域时，所有基数
$|\Ob(\mathcal{X}_{\Spec(k)})/\cong|$ 和
$|\text{Arrows}(\mathcal{X}_{\Spec(k)})|$ 的上确界 $\leq$
$(\Sch/S)_{fppf}$ 某个对象的大小。}；不关心集合论问题的读者可以忽略它；
\item[{[0]}] $\mathcal{X}$ 是 \'etale 拓扑下的广群叠；
\item[{[1]}] $\mathcal{X}$ 保极限；
\item[{[2]}] $\mathcal{X}$ 满足 Rim-Schlessinger 条件 (RS)；
\item[{[3]}] 对每个 $k$ 和 $x_0$，空间
$T\mathcal{F}_{\mathcal{X}, k, x_0}$ 和
$\text{Inf}(\mathcal{F}_{\mathcal{X}, k, x_0})$
都是有限维的；参见 (\ref{artin-equation-tangent-space}) 和
(\ref{artin-equation-infinitesimal-automorphisms})；
\item[{[4]}] 函子 (\ref{artin-equation-approximation}) 是等价；
\item[{[5]}] $\mathcal{X}$ 和
$\Delta : \mathcal{X} \to \mathcal{X} \times \mathcal{X}$ 都满足
通用性的开性。
\end{enumerate}










\section{函子的公理}
\label{artin-section-axioms-functors}

\noindent
设 $S$ 是概形，并设 $F : (\Sch/S)_{fppf}^{opp} \to \textit{Sets}$ 是函子。
% P11-E0246: the frozen source omits "by" in the naming construction for X.
把与 $F$ 相联系的集合纤维化范畴记为
$\mathcal{X} = \mathcal{S}_F$；参见《代数叠》一节
\ref{algebraic-section-split}。本节说明前述内容在 $\mathcal{X}$ 上的表述
与关于 $F$ 的陈述如何互译。

\medskip\noindent
设 $S$ 是局部 Noether 概形。设
$F : (\Sch/S)_{fppf}^{opp} \to \textit{Sets}$ 是函子。设 $k$ 是 $S$
上的有限型域，并设 $x_0 \in F(\Spec(k))$。相联系的预形变范畴
(\ref{artin-equation-predeformation-category}) 对应于函子
$$
F_{k, x_0} : \mathcal{C}_\Lambda \longrightarrow \textit{Sets},
\quad
A \longmapsto \{x \in F(\Spec(A)) \mid x|_{\Spec(k)} = x_0 \}.
$$
% P11-E0247: the frozen source uses the bare singular "functor" after "between ... and".
回忆，我们不区分 $\mathcal{C}_\Lambda$ 上的集合余纤维化范畴与函子
$\mathcal{C}_\Lambda \to \textit{Sets}$；参见《形式形变理论》评注
\ref{formal-defos-remarks-cofibered-groupoids}
(\ref{formal-defos-item-convention-cofibered-sets}).
给定函子变换 $a : F \to G$，令 $y_0 = a(x_0)$，便得到态射
$$
F_{k, x_0} \longrightarrow G_{k, y_0}
$$
；参见 (\ref{artin-equation-functoriality})。引理
\ref{artin-lemma-formally-smooth-on-deformation-categories} 表明，若
$a : F \to G$ 形式光滑（在《空间的更多态射》定义
\ref{spaces-more-morphisms-definition-formally-smooth-etale-unramified}
的意义下），则 $F_{k, x_0} \longrightarrow G_{k, y_0}$ 在《形式形变理论》
评注 \ref{formal-defos-remark-compare-smooth-schlessinger} 的意义下光滑。

\medskip\noindent
引理 \ref{artin-lemma-pushout} 断言，若在 $S$ 上概形的范畴中
$Y' = Y \amalg_X X'$，其中 $X \to X'$ 是加厚而 $X \to Y$ 是仿射的，则映射
$$
F(Y \amalg_X X') \to F(Y) \times_{F(X)} F(X')
$$
在 $F$ 是代数空间时为双射。对一般函子 $F$，若每当给定推出
$Y' = Y \amalg_X X'$，其中 $Y, X, X'$ 都是 $S$ 上有限型 Artin 局部环
之谱时，映射
$$
F(Y \amalg_X X') \to F(Y) \times_{F(X)} F(X')
$$
都是双射，则称其满足{\it Rim-Schlessinger 条件}，或称 $F$
{\it 满足 (RS)}。因此，每个代数空间都满足 (RS)。

\medskip\noindent
引理 \ref{artin-lemma-deformation-category} 断言，若函子 $F$ 满足 (RS)，则所有
$F_{k, x_0}$ 都是《形式形变理论》定义
\ref{formal-defos-definition-deformation-category} 意义下的形变函子；换言之，
它们满足《形式形变理论》评注
\ref{formal-defos-remark-compare-schlessinger-H4} 意义下的 (RS)。特别地，切空间
$$
TF_{k, x_0} = \{x \in F(\Spec(k[\epsilon])) \mid x|_{\Spec(k)} = x_0\}
$$
由《形式形变理论》引理
\ref{formal-defos-lemma-tangent-space-vector-space} 具有 $k$-向量空间结构。

\medskip\noindent
引理 \ref{artin-lemma-finite-dimension} 断言，$S$ 上局部有限型的代数空间 $F$
给出形变函子 $F_{k, x_0}$，其切空间 $TF_{k, x_0}$ 是有限维的。

\medskip\noindent
$F$ 的一个{\it 形式对象}\footnote{Artin 称之为形式形变。}
$\xi = (R, \xi_n)$ 由如下数据组成：Noether 完备局部 $S$-代数 $R$，
其剩余域在 $S$ 上有限型；以及元素
$\xi_n \in F(\Spec(R/\mathfrak m^n))$，满足
$\xi_{n + 1}|_{\Spec(R/\mathfrak m^n)} = \xi_n$。形式对象 $\xi$ 定义了
$F_{R/\mathfrak m, \xi_1}$ 的形式对象 $\xi$。当且仅当 $\xi$ 在
《形式形变理论》定义 \ref{formal-defos-definition-versal} 的意义下通用时，
称它{\it 通用}。若存在 $x \in F(\Spec(R))$，使得对所有 $n \geq 1$ 都有
$\xi_n = x|_{\Spec(R/\mathfrak m^n)}$，则称形式对象
$\xi = (R, \xi_n)$ {\it 有效}。引理 \ref{artin-lemma-effective} 断言，若 $F$
是代数空间，则每个形式对象都有效。

\medskip\noindent
设 $U$ 是 $S$ 上局部有限型的概形，并设 $x \in F(U)$。设
$u_0 \in U$ 是有限型点。当且仅当
$\xi = (\mathcal{O}_{U, u_0}^\wedge,
x|_{\Spec(\mathcal{O}_{U, u_0}/\mathfrak m_{u_0}^n)})$
是上述意义下的通用形式对象时，称 $x$ 在 $u_0$ 处通用。

\medskip\noindent
设 $S$ 是局部 Noether 概形。设
$F : (\Sch/S)_{fppf}^{opp} \to \textit{Sets}$ 是函子。下面列出我们将对
$F$ 考察的公理。
\begin{enumerate}
% P11-E0248: the frozen source repeats the malformed comparison "is <= than".
\item[{[-1]}] 一个集合论条件\footnote{该条件如下：在 $k$ 遍历 $S$ 上
有限型域时，所有基数 $|F(\Spec(k))|$ 的上确界 $\leq$
$(\Sch/S)_{fppf}$ 某个对象的大小。}；不关心集合论问题的读者可以忽略它；
\item[{[0]}] $F$ 是 \'etale 拓扑下的层；
\item[{[1]}] $F$ 保极限；
\item[{[2]}] $F$ 满足 Rim-Schlessinger 条件 (RS)；
\item[{[3]}] 每个切空间 $TF_{k, x_0}$ 都是有限维的；
\item[{[4]}] 每个形式对象都有效；
\item[{[5]}] $F$ 满足通用性的开性。
\end{enumerate}
这里，{\it 保极限}是《空间的极限》定义
\ref{spaces-limits-definition-locally-finite-presentation} 中的概念，而
{\it 通用性的开性}是指：给定 $S$ 上局部有限型的概形 $U$、
$x \in F(U)$，以及有限型点 $u_0 \in U$，若 $x$ 在 $u_0$ 处通用，
则存在开邻域 $u_0 \in U' \subset U$，使得 $x$ 在 $U'$ 的每个有限型点处
都通用。






\section{代数空间}
\label{artin-section-algebraic-spaces}

\noindent
下面给出关于代数空间的第一个主要结果。

\begin{proposition}
\label{artin-proposition-spaces-diagonal-representable}
设 $S$ 是局部 Noether 概形。设
$F : (\Sch/S)_{fppf}^{opp} \to \textit{Sets}$ 是函子。假设
\begin{enumerate}
\item $\Delta : F \to F \times F$ 可由代数空间表示；
\item $F$ 满足公理 [-1]、[0]、[1]、[2]、[3]、[4]、[5]
（参见一节 \ref{artin-section-axioms-functors}）；并且
\item 对 $S$ 的每个有限型点 $s$，$\mathcal{O}_{S, s}$ 都是 G-环。
\end{enumerate}
则 $F$ 是代数空间。
\end{proposition}

\begin{proof}
引理 \ref{artin-lemma-get-smooth} 适用于 $F$。利用这一点，对 $S$ 上每个有限型域
$k$ 和每个 $x_0 \in F(\Spec(k))$，选取 $S$ 上有限型的仿射概形
$U_{k, x_0}$ 及光滑态射 $U_{k, x_0} \to F$，使得存在剩余域为 $k$
的有限型点 $u_{k, x_0} \in U_{k, x_0}$，而 $u_{k, x_0}$ 的像为 $x_0$。于是
$$
U = \coprod\nolimits_{k, x_0} U_{k, x_0} \longrightarrow F
$$
光滑\footnote{集合论评注：由公理 [-1]，指标集的大小有界，所以这个余积
（同构于）$(\Sch/S)_{fppf}$ 的一个对象；参见《集合》引理
\ref{sets-lemma-what-is-in-it}。}。为完成证明，只须证明此映射是满的；
参见《自举》引理 \ref{bootstrap-lemma-spaces-etale-smooth-cover}
（此处用到公理 [0]）。由《可表性判据》引理
\ref{criteria-lemma-check-property-limit-preserving}，只须对那些 $V \to F$
证明 $U \times_F V \to V$ 是满的，其中 $V$ 是 $S$ 上局部有限表示的仿射
概形。由于 $U \times_F V \to V$ 光滑，其像是开的。因此，只须证明
$U \times_F V \to V$ 的像包含 $V$ 的所有有限型点；参见《态射》引理
\ref{morphisms-lemma-enough-finite-type-points}。设 $v_0 \in V$ 是有限型点。
则 $k = \kappa(v_0)$ 是 $S$ 上的有限型域。
% P11-E0249: the frozen source omits "by" in "Denote x_0 the composition".
把复合 $\Spec(k) \xrightarrow{v_0} V \to F$ 记为 $x_0$。于是
$(u_{k, x_0}, v_0) : \Spec(k) \to U \times_F V$ 是映到 $v_0$ 的点，
即得结论。
\end{proof}

\begin{lemma}
\label{artin-lemma-monomorphism}
设 $S$ 是局部 Noether 概形。设 $a : F \to G$ 是函子
$(\Sch/S)_{fppf}^{opp} \to \textit{Sets}$ 的变换。假设
\begin{enumerate}
\item $a$ 是单射；
\item $F$ 满足公理 [0]、[1]、[2]、[4] 和 [5]；
\item 对 $S$ 的每个有限型点 $s$，$\mathcal{O}_{S, s}$ 都是 G-环；
\item $G$ 是 $S$ 上局部有限型的代数空间。
\end{enumerate}
则 $F$ 是代数空间。
\end{lemma}

\begin{proof}
由引理 \ref{artin-lemma-finite-dimension}，函子 $G$ 满足 [3]。由于
$F \to G$ 是单射，可知 $F$ 也满足 [3]。此外，由于 $F \to G$ 是单射，
给定概形 $U$、$V$ 以及态射 $U \to F$ 和 $V \to F$ 时，有
$U \times_F V = U \times_G V$。因此，$\Delta : F \to F \times F$
可（由概形）表示，因为按假设这对 $G$ 成立。于是命题
\ref{artin-proposition-spaces-diagonal-representable}
适用\footnote{由于集合论条件 [-1] 对 $G$ 成立，它对 $F$ 也成立。细节从略。}。
\end{proof}












\section{代数叠}
\label{artin-section-algebraic-stacks}

\noindent
命题 \ref{artin-proposition-second-diagonal-representable} 是关于代数叠的第一个
主要结果。

\begin{lemma}
\label{artin-lemma-diagonal-representable}
设 $S$ 是局部 Noether 概形。设
$p : \mathcal{X} \to (\Sch/S)_{fppf}$ 是群胚纤维化范畴。假设
\begin{enumerate}
\item $\Delta : \mathcal{X} \to \mathcal{X} \times \mathcal{X}$
可由代数空间表示；
\item $\mathcal{X}$ 满足公理 [-1]、[0]、[1]、[2]、[3]（参见一节
\ref{artin-section-axioms}）；
\item $\mathcal{X}$ 的每个形式对象都有效；
\item $\mathcal{X}$ 满足通用性的开性；并且
\item 对 $S$ 的每个有限型点 $s$，$\mathcal{O}_{S, s}$ 都是 G-环。
\end{enumerate}
则 $\mathcal{X}$ 是代数叠。
\end{lemma}

\begin{proof}
引理 \ref{artin-lemma-get-smooth} 适用于 $\mathcal{X}$。利用这一点，对 $S$ 上
每个有限型域 $k$ 和每个对象同构类
$x_0 \in \Ob(\mathcal{X}_{\Spec(k)})$，选取 $S$ 上有限型的仿射概形
$U_{k, x_0}$ 及光滑态射 $(\Sch/U_{k, x_0})_{fppf} \to \mathcal{X}$，
使得存在剩余域为 $k$ 的有限型点
$u_{k, x_0} \in U_{k, x_0}$，而 $x_0$ 是 $u_{k, x_0}$ 的像。于是
$$
(\Sch/U)_{fppf} \to \mathcal{X},
\quad\text{其中}\quad
U = \coprod\nolimits_{k, x_0} U_{k, x_0}
$$
光滑\footnote{集合论评注：由公理 [-1]，指标集的大小有界，所以这个余积
（同构于）$(\Sch/S)_{fppf}$ 的一个对象；参见《集合》引理
\ref{sets-lemma-what-is-in-it}。}。为完成证明，只须证明此映射是满的；
参见《可表性判据》引理 \ref{criteria-lemma-stacks-etale}
（此处用到公理 [0]）。由《可表性判据》引理
\ref{criteria-lemma-check-property-limit-preserving}，只须证明
$(\Sch/U)_{fppf} \times_\mathcal{X} (\Sch/V)_{fppf} \to (\Sch/V)_{fppf}$
对那些 $y : (\Sch/V)_{fppf} \to \mathcal{X}$ 是满的，其中 $V$ 是
$S$ 上局部有限表示的仿射概形。由假设 (1)，纤维积
$(\Sch/U)_{fppf} \times_\mathcal{X} (\Sch/V)_{fppf}$ 可由代数空间 $W$
表示。于是 $W \to V$ 光滑，故其像是开的。因此，只须
证明 $W \to V$ 的像包含 $V$ 的所有有限型点；参见《态射》引理
\ref{morphisms-lemma-enough-finite-type-points}。设 $v_0 \in V$ 是有限型点。
则 $k = \kappa(v_0)$ 是 $S$ 上的有限型域。
% P11-E0250: the frozen source uses an ungrammatical "Denote x_0 = ... the pullback" construction.
把 $y$ 沿 $v_0$ 的拉回记为 $x_0 = y|_{\Spec(k)}$。于是
$(u_{k, x_0}, v_0)$ 给出态射 $\Spec(k) \to W$，它与 $W \to V$ 的复合
为 $v_0$，即得结论。
\end{proof}

\begin{proposition}
\label{artin-proposition-second-diagonal-representable}
设 $S$ 是局部 Noether 概形。设
$p : \mathcal{X} \to (\Sch/S)_{fppf}$ 是群胚纤维化范畴。假设
\begin{enumerate}
\item $\Delta_\Delta : \mathcal{X} \to
\mathcal{X} \times_{\mathcal{X} \times \mathcal{X}} \mathcal{X}$
可由代数空间表示；
\item $\mathcal{X}$ 满足公理 [-1]、[0]、[1]、[2]、[3]、[4] 和 [5]
（参见一节 \ref{artin-section-axioms}）；
\item 对 $S$ 的每个有限型点 $s$，$\mathcal{O}_{S, s}$ 都是 G-环。
\end{enumerate}
则 $\mathcal{X}$ 是代数叠。
\end{proposition}

\begin{proof}
先证明 $\Delta : \mathcal{X} \to \mathcal{X} \times \mathcal{X}$ 可由
代数空间表示。为此，只须证明
$$
\mathcal{Y} =
\mathcal{X} \times_{\Delta, \mathcal{X} \times \mathcal{X}, y} (\Sch/V)_{fppf}
$$
对任意 $S$ 上局部有限表示的仿射概形 $V$ 和 $\mathcal{X} \times \mathcal{X}$
位于 $V$ 上的对象 $y$，都可由代数空间表示；参见《可表性判据》引理
\ref{criteria-lemma-check-representable-limit-preserving}\footnote{《可表性判据》
引理 \ref{criteria-lemma-check-representable-limit-preserving} 中的集合论条件
成立：表示 $\mathcal{Y}$ 的代数空间 $Y$ 的大小适当地有界。具体地，
$Y \to S$ 将局部有限型，并且 $Y$ 将满足公理 [-1]。细节从略。}。
注意，$\mathcal{Y}$ 以集合胚为纤维（《叠》引理
\ref{stacks-lemma-isom-as-2-fibre-product}）。令
$Y : (\Sch/S)_{fppf}^{opp} \to \textit{Sets}$，
$T \mapsto \Ob(\mathcal{Y}_T)/\cong$ 为同构类函子。我们将应用命题
\ref{artin-proposition-spaces-diagonal-representable}，证明 $Y$ 是代数空间。

\medskip\noindent
注意，
$\Delta_\mathcal{Y} : \mathcal{Y} \to \mathcal{Y} \times \mathcal{Y}$
（因而 $Y \to Y \times Y$ 也一样）由条件 (1) 和《可表性判据》引理
\ref{criteria-lemma-second-diagonal} 可由代数空间表示。由《叠》引理
\ref{stacks-lemma-stack-in-setoids-characterize} 和
\ref{stacks-lemma-2-fibre-product-gives-stack-in-setoids}，$Y$ 是 \'etale
拓扑下的层；换言之，公理 [0] 成立。又由引理
\ref{artin-lemma-fibre-product-limit-preserving}，$Y$ 保极限，即 [1] 成立。
由引理 \ref{artin-lemma-algebraic-stack-RS} 和
\ref{artin-lemma-fibre-product-RS}，$Y$ 具有 (RS)，即公理 [2] 成立。$Y$ 的
公理 [3] 由引理 \ref{artin-lemma-finite-dimension} 和
\ref{artin-lemma-fibre-product-tangent-spaces} 得到。公理 [4] 由引理
\ref{artin-lemma-effective} 和 \ref{artin-lemma-fibre-product-effective} 得到。
$Y$ 的公理 [5] 直接来自 $\Delta_\mathcal{X}$ 的通用性的开性；后者是
$\mathcal{X}$ 的公理 [5] 的一部分。因此，命题
\ref{artin-proposition-spaces-diagonal-representable} 的所有假设都满足，且 $Y$
是代数空间。

\medskip\noindent
此时，由引理 \ref{artin-lemma-diagonal-representable}，$\mathcal{X}$ 是代数叠。
\end{proof}





\section{强 Rim--Schlessinger 条件}
\label{artin-section-RS-star}

\noindent
在本章余下部分，下面这个严格更强的 Rim--Schlessinger 条件将起重要作用。

\begin{definition}
\label{artin-definition-RS-star}
设 $S$ 是概形，$\mathcal{X}$ 是 $(\Sch/S)_{fppf}$ 上的群胚纤维化范畴。
若给定 $S$-代数的纤维积图
$$
\xymatrix{
B' \ar[r] & B \\
A' = A \times_B B' \ar[u] \ar[r] & A \ar[u]
}
$$
，其中 $B' \to B$ 为具有平方零核的满射时，纤维范畴的函子
$$
\mathcal{X}_{\Spec(A')}
\longrightarrow
\mathcal{X}_{\Spec(A)} \times_{\mathcal{X}_{\Spec(B)}} \mathcal{X}_{\Spec(B')}
$$
是范畴等价，则称 $\mathcal{X}$ 满足{\it 条件 (RS*)}。
\end{definition}

\noindent
作如下观察。设 $A \to B \leftarrow B'$ 如定义 \ref{artin-definition-RS-star} 中所述。
\begin{enumerate}
\item 在概形范畴中有
$\Spec(A') = \Spec(A) \amalg_{\Spec(B)} \Spec(B')$；参见《更多态射》引理
\ref{more-morphisms-lemma-pushout-along-thickening}；并且
\item 若 $\mathcal{X}$ 是代数叠，则由引理
\ref{artin-lemma-algebraic-stack-RS-star}，$\mathcal{X}$ 满足 (RS*)。
\end{enumerate}
若 $S$ 局部 Noether，则
\begin{enumerate}
\item[(3)] 若 $A$、$B$、$B'$ 在 $S$ 上有限型，且 $B$ 在 $A$ 上有限，
则 $A'$ 在 $S$ 上有限型\footnote{若 $\Spec(A)$ 映入 $S$ 的某个仿射开集，
这由《更多代数》引理 \ref{more-algebra-lemma-fibre-product-finite-type} 得到。
一般情形使用《更多代数》引理
\ref{more-algebra-lemma-diagram-localize} 得到。}；并且
\item[(4)] 若 $\mathcal{X}$ 满足 (RS*)，则 $\mathcal{X}$ 满足 (RS)，
因为 (RS) 恰好涵盖 (RS*) 中 $A$、$B$、$B'$ 为 Artin 局部环的情形。
\end{enumerate}

\begin{lemma}
\label{artin-lemma-algebraic-stack-RS-star}
设 $\mathcal{X}$ 是基概形 $S$ 上的代数叠。则 $\mathcal{X}$ 满足 (RS*)。
\end{lemma}

\begin{proof}
这由引理 \ref{artin-lemma-pushout} 推出；参见定义
\ref{artin-definition-RS-star} 后的评注。
\end{proof}

\begin{lemma}
\label{artin-lemma-fibre-product-RS-star}
设 $S$ 是概形。设 $p : \mathcal{X} \to \mathcal{Y}$ 和
$q : \mathcal{Z} \to \mathcal{Y}$ 是 $(\Sch/S)_{fppf}$ 上群胚纤维化
范畴的 $1$-态射。若 $\mathcal{X}$、$\mathcal{Y}$ 和 $\mathcal{Z}$
满足 (RS*)，则 $\mathcal{X} \times_\mathcal{Y} \mathcal{Z}$ 也满足 (RS*)。
\end{lemma}

\begin{proof}
证明与引理 \ref{artin-lemma-fibre-product-RS} 的证明完全相同。
\end{proof}








\section{通用性与一般化}
\label{artin-section-generalize-versality}

\noindent
我们证明，对具有 (RS*) 且保极限的叠，通用性在一般化下保持。
初次阅读时，建议跳过本节。

\begin{lemma}
\label{artin-lemma-single-point}
设 $S$ 是局部 Noether 概形，$\mathcal{X}$ 是 $(\Sch/S)_{fppf}$ 上具有
(RS*) 的群胚纤维化范畴。设 $x$ 是 $\mathcal{X}$ 位于 $S$ 上有限型仿射
概形 $U$ 上的对象。设 $u \in U$ 是有限型点，且 $x$ 不在 $u$ 处通用。
则存在 $\mathcal{X}$ 的态射 $x \to y$，它位于 $U \to T$ 之上并满足
\begin{enumerate}
\item 态射 $U \to T$ 是一阶加厚；
\item 有短正合列
$$
0 \to \kappa(u) \to \mathcal{O}_T \to \mathcal{O}_U \to 0
$$
% P11-E0251: condition (3) names the pair with alpha but defines and later uses beta.
\item {\bf 不}存在二元组 $(W, \alpha)$，它由 $u$ 的开邻域
$W \subset T$ 和态射 $\beta : y|_W \to x$ 组成，并使得复合
$$
x|_{U \cap W} \xrightarrow{\text{下列对象的限制： }x \to y}
y|_W \xrightarrow{\beta} x
$$
是典范态射 $x|_{U \cap W} \to x$。
\end{enumerate}
\end{lemma}

\begin{proof}
令 $R = \mathcal{O}_{U, u}^\wedge$。令 $k = \kappa(u)$ 为 $R$ 的剩余域。
令 $\xi$ 为与 $x$ 相联系的 $\mathcal{X}$ 在 $R$ 上的形式对象。由于 $x$
不在 $u$ 处通用，由引理 \ref{artin-lemma-versality-matches}，$\xi$ 不通用。
根据定义 \ref{artin-definition-versal-formal-object} 后的讨论，这意味着可以找到
$\mathcal{X}$ 的态射 $\xi_1 \to x_A \to x_B$，它位于闭浸入
$\Spec(k) \to \Spec(A) \to \Spec(B)$ 之上，其中 $A,B$ 是剩余域为 $k$
的 Artin 局部环；还可以找到 $n \geq 1$ 和交换图
$$
\vcenter{
\xymatrix{
& x_A \ar[ld] \\
\xi_n & \xi_1 \ar[u] \ar[l]
}
}
\quad\text{位于下列对象之上：}\quad
\vcenter{
\xymatrix{
& \Spec(A) \ar[ld] \\
\Spec(R/\mathfrak m^n) & \Spec(k) \ar[u] \ar[l]
}
}
$$
，使得{\bf 不}存在 $m \geq n$ 和交换图
$$
\vcenter{
\xymatrix{
& & x_B \ar[lldd] \\
& & x_A \ar[ld] \ar[u] \\
\xi_m & \xi_n \ar[l] & \xi_1 \ar[u] \ar[l]
}
}
\quad\text{位于下列对象之上：}
\vcenter{
\xymatrix{
& & \Spec(B) \ar[lldd] \\
& & \Spec(A) \ar[ld] \ar[u] \\
\Spec(R/\mathfrak m^m) &
\Spec(R/\mathfrak m^n) \ar[l] &
\Spec(k) \ar[u] \ar[l]
}
}
$$
还可假设 $B \to A$ 是小扩张；换言之，满射 $B \to A$ 的核 $I$ 作为
$A$-模同构于 $k$。这由《形式形变理论》评注
\ref{formal-defos-remark-versal-object} 得到。然后直接定义
$$
T = U \amalg_{\Spec(A)} \Spec(B)
$$
由性质 (RS*)，可找到 $T$ 上的 $y$，其在 $\Spec(B)$ 上的限制为 $x_B$，
在 $U$ 上的限制为 $x$（这给出位于 $U \to T$ 之上的箭头 $x \to y$）。
为完成证明，核验条件 (1)、(2) 和 (3)。

\medskip\noindent
由推出的构造，有交换图
$$
\xymatrix{
0 \ar[r] &
I \ar[r] &
B \ar[r] &
A \ar[r] &
0 \\
0 \ar[r] &
I \ar[r] \ar[u] &
\Gamma(T, \mathcal{O}_T) \ar[r] \ar[u] &
\Gamma(U, \mathcal{O}_U) \ar[r] \ar[u] &
0
}
$$
，其两行正合。这立即证明了 (1) 和 (2)。为完成证明，采用反证法。
假设有 (3) 中的二元组 $(W, \beta)$。由于 $\Spec(B) \to T$ 分解经过
$W$，得到态射
$$
x_B \to y|_W \xrightarrow{\beta} x
$$
由于 $B$ 是剩余域为 $k = \kappa(u)$ 的 Artin 局部环，$x_B \to x$
位于某个态射 $\Spec(B) \to U$ 之上，而该态射对某个 $m \geq n$ 分解经过
$\Spec(\mathcal{O}_{U, u}/\mathfrak m_u^m)$。换言之，$x_B \to x$ 分解经过
$\xi_m$，从而给出映射 $x_B \to \xi_m$。
% P11-E0252: the proof invokes alpha although condition (3) defines the relevant morphism as beta.
条件 (3) 中态射 $\alpha$ 的相容条件转化为下图交换：
$$
\xymatrix{
x_B \ar[d] & x_A \ar[d] \ar[l] \\
\xi_m & \xi_n \ar[l]
}
$$
这就给出所需的矛盾。
\end{proof}

\begin{lemma}
\label{artin-lemma-generalization-versality}
设 $S$ 是局部 Noether 概形，$\mathcal{X}$ 是 $(\Sch/S)_{fppf}$ 上的
群胚纤维化范畴。假设
\begin{enumerate}
\item $\Delta : \mathcal{X} \to \mathcal{X} \times \mathcal{X}$
可由代数空间表示；
\item $\mathcal{X}$ 具有 (RS*)；
\item $\mathcal{X}$ 保极限。
\end{enumerate}
设 $x$ 是 $\mathcal{X}$ 位于 $S$ 上有限型概形 $U$ 上的对象。设
$u \leadsto u_0$ 是 $U$ 的有限型点之间的特殊化，且 $x$ 在 $u_0$ 处通用。
则 $x$ 在 $u$ 处通用。
\end{lemma}

\begin{proof}
缩小 $U$ 后，可假设 $U$ 仿射，且 $U$ 映入 $S$ 的一个仿射开集
$\Spec(\Lambda)$。若 $x$ 不在 $u$ 处通用，则可按引理
\ref{artin-lemma-single-point} 选取位于 $U \to T$ 之上的 $x \to y$。写成
$U = \Spec(R_0)$ 和 $T = \Spec(R)$。态射 $U \to T$ 对应于满环同态
$R \to R_0$，其核是平方零理想。由假设 (3)，存在某个有限型
$\Lambda$-子代数 $R' \subset R$，使得 $y$ 来自
$U' = \Spec(R')$ 上的对象 $x'$。增大 $R'$ 后，可以且确实假设
$R' \to R_0$ 是满射，因而 $U \subset U'$ 是一阶加厚。于是现在有
$$
x \to y \to x'
\text{ 位于下列对象之上： }
U \to T \to U'
$$
由假设 (1)，存在 $S$ 上的代数空间 $Z$ 表示
$$
(\Sch/U)_{fppf} \times_{x, \mathcal{X}, x'} (\Sch/U')_{fppf}
$$
；参见《代数叠》引理 \ref{algebraic-lemma-representable-diagonal}。由
$2$-纤维积的构造，$Z$ 的一个 $V$-值点对应于三元组
$(a, a', \alpha)$，它由态射 $a : V \to U$、$a' : V \to U'$ 和态射
$\alpha : a^*x \to (a')^*x'$ 组成。得到交换图
$$
\xymatrix{
U \ar[rd] \ar[rdd] \ar[rrd] \\
& Z \ar[r]_{p'} \ar[d]^p & U' \ar[d] \\
& U \ar[r] & S
}
$$
% P11-E0253: the frozen source says "comes the isomorphism", omitting "from".
态射 $i : U \to Z$ 来自同构 $x \to x'|_U$。令 $z_0 = i(u_0) \in Z$。
由引理 \ref{artin-lemma-base-change-versal}，$Z \to U'$ 在 $z_0$ 处光滑。
把 $U$ 替换为 $u_0$ 的一个仿射开邻域，把 $U'$ 替换为相应的开集，再把
$Z$ 替换为这些开集分别在 $p$ 和 $p'$ 下的逆像之交后，便达到
$Z \to U'$ 沿 $i(U)$ 光滑的情形。由于 $u \leadsto u_0$，点 $u$ 位于这个
开集中。缩小 $U$ 显然不破坏引理 \ref{artin-lemma-single-point} 的条件 (3)
（概形 $U$、$T$、$U'$ 具有相同的底拓扑空间）。由于 $U \to U'$ 是仿射
概形的一阶加厚，可以选取态射 $i' : U' \to Z$，使得
$p' \circ i' = \text{id}_{U'}$，且其在 $U$ 上的限制为 $i$（《空间的更多态射》
引理 \ref{spaces-more-morphisms-lemma-smooth-formally-smooth}）。把万有态射
$p^*x \to (p')^*x'$ 沿 $i'$ 拉回，得到态射
$$
x' \to x
$$
，它位于 $p \circ i' : U' \to U$ 之上，并使得复合
$$
x \to x' \to x
$$
是恒等态射。回忆，已有位于态射 $T \to U'$ 之上的 $y \to x'$。作复合便
得到态射 $y \to x$，它的存在与引理 \ref{artin-lemma-single-point} 的条件 (3)
矛盾。这个矛盾完成了证明。
\end{proof}






\section{强形式有效性}
\label{artin-section-strong-formal-effectiveness}

\noindent
本节说明，形式对象有效性的一个强版本如何蕴含通用性的开性。
\cite[Theorem 1.1]{Bhatt-Algebraize} 的证明表明，拟紧且拟分离的代数空间
满足评注 \ref{artin-remark-strong-effectiveness} 中讨论的强形式有效性。此外，
这里发展的理论并非空泛：后文将用它证明凝聚层叠和复形模空间的通用性的开性；
参见《Quot》定理
\ref{quot-theorem-coherent-algebraic-general} 和
\ref{quot-theorem-complexes-algebraic}.

\begin{lemma}
\label{artin-lemma-infinite-sequence-pre}
设 $S$ 是局部 Noether 概形，$\mathcal{X}$ 是 $(\Sch/S)_{fppf}$ 上具有
(RS*) 的群胚纤维化范畴。设 $x$ 是 $\mathcal{X}$ 位于 $S$ 上有限型
仿射概形 $U$ 上的对象。
% P11-E0254: the frozen source says "a finite type points".
设 $u_n \in U$、$n \geq 1$ 是有限型点，并满足：(a) 当 $n \not = m$ 时，
不存在特殊化 $u_n \leadsto u_m$；(b) 对所有 $n$，$x$ 都不在 $u_n$ 处通用。
则存在态射
$$
x \to x_1 \to x_2 \to \ldots
\quad\text{于 }\mathcal{X}\text{ 位于下列对象之上： }\quad
U \to U_1 \to U_2 \to \ldots
$$
，它们位于 $S$ 之上并满足
\begin{enumerate}
\item 对每个 $n$，态射 $U \to U_n$ 是一阶加厚；
\item 对每个 $n$，有短正合列
$$
0 \to \kappa(u_n) \to \mathcal{O}_{U_n} \to \mathcal{O}_{U_{n - 1}} \to 0
$$
，其中当 $n = 1$ 时 $U_0 = U$；
\item 对每个 $n$，{\bf 不}存在二元组 $(W, \alpha)$，它由 $u_n$ 的开邻域
$W \subset U_n$ 和态射 $\alpha : x_n|_W \to x$ 组成，并使得复合
$$
x|_{U \cap W} \xrightarrow{\text{下列对象的限制： }x \to x_n}
x_n|_W \xrightarrow{\alpha} x
$$
是典范态射 $x|_{U \cap W} \to x$。
\end{enumerate}
\end{lemma}

\begin{proof}
由于诸点 $u_n$ 之间没有特殊化（特别地，$u_n$ 两两不同），对每个 $n$
都可以找到开集 $U' \subset U$，使得 $u_n \in U'$，而当
$i = 1, \ldots, n - 1$ 时 $u_i \not \in U'$。由引理
\ref{artin-lemma-single-point}，对每个 $n \geq 1$ 都可以找到
$$
x \to y_n
\quad\text{于 }\mathcal{X}\text{ 位于下列对象之上：}\quad
U \to T_n
$$
，使得
\begin{enumerate}
\item 态射 $U \to T_n$ 是一阶加厚；
\item 有短正合列
$$
0 \to \kappa(u_n) \to \mathcal{O}_{T_n} \to \mathcal{O}_U \to 0
$$
% P11-E0255: the auxiliary condition names the pair with alpha but defines the morphism as beta.
\item {\bf 不}存在二元组 $(W, \alpha)$，它由 $u_n$ 的开邻域
$W \subset T_n$ 和态射 $\beta : y_n|_W \to x$ 组成，并使得复合
$$
x|_{U \cap W} \xrightarrow{\text{下列对象的限制： }x \to y_n}
y_n|_W \xrightarrow{\beta} x
$$
是典范态射 $x|_{U \cap W} \to x$。
\end{enumerate}
因此可以归纳定义
$$
U_1 = T_1, \quad
U_{n + 1} = U_n \amalg_U T_{n + 1}
$$
令 $x_1 = y_1$ 并使用 (RS*)，可归纳找到 $U_{n + 1}$ 上的 $x_{n + 1}$，
其在 $U_n$ 上的限制为 $x_n$，在 $T_{n + 1}$ 上的限制为 $y_{n + 1}$。
$U \to U_n$ 的性质 (1) 来自《更多态射》引理
\ref{more-morphisms-lemma-pushout-along-thickening} 中推出的构造。类似地，
$U_n$ 的性质 (2) 由推出的构造从 $T_n$ 的性质 (2) 得到。按上面的讨论，
缩小到 $u_n$ 的一个开邻域 $U'$ 后，$(U_n, x_n)$ 的性质 (3) 直接由
$(T_n, y_n)$ 的性质 (3) 得到，因为 $T_n$ 和 $U_n$ 的相应开子概形同构。
省略一些细节。
\end{proof}

\begin{remark}[强有效性]
\label{artin-remark-strong-effectiveness}
设 $S$ 是局部 Noether 概形，$\mathcal{X}$ 是 $(\Sch/S)_{fppf}$ 上的
群胚纤维化范畴。假设有
\begin{enumerate}
\item 仿射开集 $\Spec(\Lambda) \subset S$；
\item $\Lambda$-代数的逆系统 $(R_n)$，其转移映射为满射，且核局部幂零；
\item $\mathcal{X}$ 的对象系统 $(\xi_n)$，它位于系统
$(\Spec(R_n))$ 之上。
\end{enumerate}
在这种情形下，令 $R = \lim R_n$。若存在 $\mathcal{X}$ 位于
$\Spec(R)$ 上的对象 $\xi$，其在 $\Spec(R_n)$ 上的限制给出系统
$(\xi_n)$，则称 $(\xi_n)$ {\it 有效}。
\end{remark}

\noindent
并非 $S$ 上每个代数叠 $\mathcal{X}$ 都满足如下形式的强有效性公理：
评注 \ref{artin-remark-strong-effectiveness} 中的每个系统 $(\xi_n)$ 都有效。
《例子》一节 \ref{examples-section-non-formal-effectiveness} 给出了一个反例。

\begin{lemma}
\label{artin-lemma-SGE-implies-openness-versality}
设 $S$ 是局部 Noether 概形，$\mathcal{X}$ 是 $(\Sch/S)_{fppf}$ 上的
群胚纤维化范畴。假设
\begin{enumerate}
\item $\Delta : \mathcal{X} \to \mathcal{X} \times \mathcal{X}$
可由代数空间表示；
\item $\mathcal{X}$ 具有 (RS*)；
\item $\mathcal{X}$ 保极限；
\item 评注 \ref{artin-remark-strong-effectiveness} 中的系统 $(\xi_n)$ 都有效，
其中对所有 $m \geq n$，$\Ker(R_m \to R_n)$ 都是平方零理想。
\end{enumerate}
则 $\mathcal{X}$ 满足通用性的开性。
\end{lemma}

\begin{proof}
选取 $S$ 上局部有限型的概形 $U$、$U$ 的有限型点 $u_0$，以及
$\mathcal{X}$ 位于 $U$ 上的对象 $x$，使得 $x$ 在 $u_0$ 处通用。缩小
$U$ 后，可假设 $U$ 仿射，且 $U$ 映入 $S$ 的仿射开集
$\Spec(\Lambda)$。令 $E \subset U$ 为如下有限型点 $u$ 所成的集合：$x$
不在 $u$ 处通用。由引理 \ref{artin-lemma-generalization-versality}，若
$u \in E$，则 $u_0$ 不是 $u$ 的特殊化。若通用性的开性不成立，则 $u_0$
位于 $E$ 的闭包 $\overline{E}$ 中。由《性质》引理
\ref{properties-lemma-countable-dense-subset}，可选取与 $E$ 具有相同闭包的
可数子集 $E' \subset E$。由《性质》引理
\ref{properties-lemma-maximal-points}，可假设 $E'$ 的诸点之间没有特殊化。
按上面的观察，$u_0$ 不是 $E'$ 中任何点的特殊化，所以 $E'$ 必须为
（可数）无限集。因此可写成 $E' = \{u_1, u_2, u_3, \ldots\}$；诸 $u_i$
之间没有特殊化，且 $u_0$ 位于 $E'$ 的闭包中。

\medskip\noindent
按引理 \ref{artin-lemma-infinite-sequence-pre} 选取位于
$U \to U_1 \to U_2 \to \ldots$ 之上的
$x \to x_1 \to x_2 \to \ldots$。写成 $U_n = \Spec(R_n)$ 和
$U = \Spec(R_0)$。令 $R = \lim R_n$。注意，$R \to R_0$ 是满射，
其核是平方零理想。由假设 (4)，得到 $\Spec(R)$ 上的 $\xi$，其到 $R_n$
的基变换为 $x_n$。由假设 (3)，存在某个有限型 $\Lambda$-子代数
$R' \subset R$，使得 $\xi$ 来自 $U' = \Spec(R')$ 上的对象 $\xi'$。
增大 $R'$ 后，可以且确实假设 $R' \to R_0$ 是满射，因而
$U \subset U'$ 是一阶加厚。于是现在有
$$
x \to x_1 \to x_2 \to \ldots \to \xi'
\text{ 位于下列对象之上： }
U \to U_1 \to U_2 \to \ldots \to U'
$$
由假设 (1)，存在 $S$ 上的代数空间 $Z$ 表示
$$
(\Sch/U)_{fppf} \times_{x, \mathcal{X}, \xi'} (\Sch/U')_{fppf}
$$
；参见《代数叠》引理 \ref{algebraic-lemma-representable-diagonal}。由
$2$-纤维积的构造，$Z$ 的一个 $T$-值点对应于三元组
$(a, a', \alpha)$，它由态射 $a : T \to U$、$a' : T \to U'$ 和态射
$\alpha : a^*x \to (a')^*\xi'$ 组成。得到交换图
$$
\xymatrix{
U \ar[rd] \ar[rdd] \ar[rrd] \\
& Z \ar[r]_{p'} \ar[d]^p & U' \ar[d] \\
& U \ar[r] & S
}
$$
% P11-E0256: the frozen source again says "comes the isomorphism", omitting "from".
态射 $i : U \to Z$ 来自同构 $x \to \xi'|_U$。令 $z_0 = i(u_0) \in Z$。
由引理 \ref{artin-lemma-base-change-versal}，$Z \to U'$ 在 $z_0$ 处光滑。
把 $U$ 替换为 $u_0$ 的一个仿射开邻域，把 $U'$ 替换为相应的开集，再把
$Z$ 替换为这些开集分别在 $p$ 和 $p'$ 下的逆像之交后，便达到
$Z \to U'$ 沿 $i(U)$ 光滑的情形。注意，这还涉及把 $u_n$ 替换为一个
子列，即只保留 $u_n$ 位于该开集中的那些指标。此外，缩小 $U$ 显然不破坏
引理 \ref{artin-lemma-infinite-sequence-pre} 的条件 (3)（概形 $U$、$U_n$、$U'$
具有相同的底拓扑空间）。由于 $U \to U'$ 是仿射概形的一阶加厚，可以选取
态射 $i' : U' \to Z$，使得 $p' \circ i' = \text{id}_{U'}$，且其在 $U$
上的限制为 $i$（《空间的更多态射》引理
\ref{spaces-more-morphisms-lemma-smooth-formally-smooth}）。把万有态射
$p^*x \to (p')^*\xi'$ 沿 $i'$ 拉回，得到态射
$$
\xi' \to x
$$
，它位于 $p \circ i' : U' \to U$ 之上，并使得复合
$$
x \to \xi' \to x
$$
是恒等态射。回忆，已有位于态射 $U_1 \to U'$ 之上的
$x_1 \to \xi'$。作复合便得到态射 $x_1 \to x$，它的存在与引理
\ref{artin-lemma-infinite-sequence-pre} 的条件 (3) 矛盾。这个矛盾完成了证明。
\end{proof}

\begin{remark}
\label{artin-remark-trade-openness-versality-diagonal-with-strong-effectiveness}
% P11-E0257: the frozen source uses "an category" instead of "a category".
可以从一个强形式有效性公理推出群胚纤维化范畴的对角态射满足通用性的开性。
设 $S$ 是局部 Noether 概形，$\mathcal{X}$ 是 $(\Sch/S)_{fppf}$ 上的
群胚纤维化范畴。假设
\begin{enumerate}
\item $\Delta_\Delta : \mathcal{X} \to
\mathcal{X} \times_{\mathcal{X} \times \mathcal{X}} \mathcal{X}$
可由代数空间表示；
\item $\mathcal{X}$ 具有 (RS*)；
\item $\mathcal{X}$ 保极限；
\item 给定评注 \ref{artin-remark-strong-effectiveness} 中的 $S$-代数逆系统
$(R_n)$，其中对所有 $m \geq n$，$\Ker(R_m \to R_n)$ 都是平方零理想时，
函子
$$
\mathcal{X}_{\Spec(\lim R_n)} \longrightarrow
\lim_n \mathcal{X}_{\Spec(R_n)}
$$
全忠实。
\end{enumerate}
则 $\Delta : \mathcal{X} \to \mathcal{X} \times \mathcal{X}$ 满足通用性的
开性。这由把引理 \ref{artin-lemma-SGE-implies-openness-versality} 应用于如下形式的
纤维积得到：
$\mathcal{X} \times_{\Delta, \mathcal{X} \times \mathcal{X}, y}
 (\Sch/V)_{fppf}$，其中 $V$ 是任意 $S$ 上局部有限表示的仿射概形，
而 $y$ 是 $\mathcal{X} \times \mathcal{X}$ 位于 $V$ 上的对象。
若以后需要这个结果，我们会把本评注改成引理并给出详细证明。
\end{remark}











\section{无穷小形变}
\label{artin-section-inf}

\noindent
本节讨论一节 \ref{artin-section-tangent-spaces} 中引入的切空间概念的一种推广。
为妥善处理这一问题，借用《形式形变理论》各节
\ref{formal-defos-section-tangent-spaces-functors},
\ref{formal-defos-section-lifts} 和
\ref{formal-defos-section-infinitesimal-automorphisms} 中的一些记号。

\medskip\noindent
设 $S$ 是概形，$\mathcal{X}$ 是 $(\Sch/S)_{fppf}$ 上的群胚纤维化范畴。
给定 $S$-代数同态 $A' \to A$ 和 $\mathcal{X}$ 位于 $\Spec(A)$ 上的对象
$x$，用 $\textit{Lift}(x, A')$ 表示 $x$ 到 $\Spec(A')$ 的提升所成的范畴。
$\textit{Lift}(x, A')$ 的对象是 $\mathcal{X}$ 的态射 $x \to x'$，它位于
$\Spec(A) \to \Spec(A')$ 之上；$\textit{Lift}(x, A')$ 的态射定义为交换图。
$\textit{Lift}(x, A')$ 的同构类集合记为 $\text{Lift}(x, A')$。参见
《形式形变理论》定义 \ref{formal-defos-definition-lifts} 和评注
\ref{formal-defos-remark-omit-arrow}。若 $A' \to A$ 是具有局部幂零核的满射，
则称 $\text{Lift}(x, A')$ 的元素 $x'$ 为 $x$ 的{\it（无穷小）形变}。
在这种情形下，$x'$ 在 $x$ 之上的{\it 无穷小自同构群}是核
$$
\text{Inf}(x'/x) =
\Ker\left(
\text{Aut}_{\mathcal{X}_{\Spec(A')}}(x') \to
\text{Aut}_{\mathcal{X}_{\Spec(A)}}(x)\right)
$$
注意，$\text{Inf}(x'/x)$ 的元素等同于 $\mathit{Aut}_\mathcal{X}(x')$
（相联系的集合纤维化范畴）中 $\text{id}_x$ 在 $\Spec(A')$ 上的提升。
请与《形式形变理论》定义
\ref{formal-defos-definition-relative-infinitesimal-auts} 和《形式形变理论》评注
\ref{formal-defos-remark-infaut-lifting-equalities} 比较。

\medskip\noindent
% P11-E0258: the frozen source omits "by" in the naming construction for A[M].
若 $M$ 是 $A$-模，把 $A[M]$ 记为如下 $A$-代数：其底 $A$-模为
$A \oplus M$，乘法由
$(a, m) \cdot (a', m') = (aa', am' + a'm)$ 给出。当 $M = A$ 时，这是
$A$ 上的对偶数环，按惯例记为 $A[\epsilon]$。存在 $A$-代数同态
$A[M] \to A$。$x$ 到 $\Spec(A[M])$ 的拉回称为 $x$ 到
$\Spec(A[M])$ 的{\it 平凡形变}。

\begin{lemma}
\label{artin-lemma-functoriality}
设 $S$ 是概形。设 $f : \mathcal{X} \to \mathcal{Y}$ 是
$(\Sch/S)_{fppf}$ 上群胚纤维化范畴的 $1$-态射。设
$$
\xymatrix{
B' \ar[r] & B \\
A' \ar[u] \ar[r] & A \ar[u]
}
$$
是 $S$-代数的交换图。设 $x$ 是 $\mathcal{X}$ 位于 $\Spec(A)$ 上的对象，
$y$ 是 $\mathcal{Y}$ 位于 $\Spec(B)$ 上的对象，而
$\phi : f(x)|_{\Spec(B)} \to y$ 是 $\mathcal{Y}$ 位于 $\Spec(B)$ 上的态射。
则存在典范函子
$$
\textit{Lift}(x, A') \longrightarrow \textit{Lift}(y, B')
$$
，它由 $f$ 和 $\phi$ 在提升范畴之间诱导。此构造以显然的方式与群胚
纤维化范畴的 $1$-态射的复合相容。
\end{lemma}

\begin{proof}
这个引理不证自明。
\end{proof}

\noindent
设 $S$ 是基概形，$\mathcal{X}$ 是 $(\Sch/S)_{fppf}$ 上的群胚纤维化范畴。
定义一个范畴，其对象是二元组 $(x, A' \to A)$，其中
\begin{enumerate}
\item $A' \to A$ 是 $S$-代数的满射，其核为平方零理想；
\item $x$ 是 $\mathcal{X}$ 位于 $\Spec(A)$ 上的对象。
\end{enumerate}
态射 $(y, B' \to B) \to (x, A' \to A)$ 由如下数据给出：交换图
$$
\xymatrix{
B' \ar[r] & B \\
A' \ar[u] \ar[r] & A \ar[u]
}
$$
所示的 $S$-代数同态，以及位于 $\Spec(B)$ 之上的态射
$x|_{\Spec(B)} \to y$。称之为{\it 形变情形}范畴。

\begin{lemma}
\label{artin-lemma-properties-lift-RS-star}
设 $S$ 是概形，$\mathcal{X}$ 是 $(\Sch/S)_{fppf}$ 上的群胚纤维化范畴。
假设 $\mathcal{X}$ 满足条件 (RS*)。设 $A$ 是 $S$-代数，而 $x$ 是
$\mathcal{X}$ 位于 $\Spec(A)$ 上的对象。
\begin{enumerate}
\item 存在 $A$-线性函子
$\text{Inf}_x : \text{Mod}_A \to \text{Mod}_A$
，使得给定形变情形 $(x, A' \to A)$ 和提升 $x'$ 时，存在同构
$\text{Inf}_x(I) \to \text{Inf}(x'/x)$，其中 $I = \Ker(A' \to A)$。
\item 存在 $A$-线性函子
$T_x : \text{Mod}_A \to \text{Mod}_A$
，使得
\begin{enumerate}
\item 给定 $\text{Mod}_A$ 中的 $M$，存在双射
$T_x(M) \to \text{Lift}(x, A[M])$；
\item 给定形变情形 $(x, A' \to A)$，存在作用
$$
T_x(I) \times \text{Lift}(x, A') \to \text{Lift}(x, A')
$$
，其中 $I = \Ker(A' \to A)$。若
$\text{Lift}(x, A') \not = \emptyset$，则此作用单传递。
\end{enumerate}
\end{enumerate}
\end{lemma}

\begin{proof}
把 $\text{Inf}_x$ 定义为函子
$$
\text{Mod}_A \longrightarrow \textit{Sets},\quad
M \longrightarrow
\text{Inf}(x'_M/x) = \text{Lift}(\text{id}_x, A[M])
$$
；它把 $M$ 映到 $x$ 到 $\Spec(A[M])$ 的平凡形变 $x'_M$ 的无穷小
自同构群，等价地，映到 $\mathit{Aut}_\mathcal{X}(x'_M)$ 中
$\text{id}_x$ 的提升所成的群。把 $T_x$ 定义为函子
$$
\text{Mod}_A \longrightarrow \textit{Sets},\quad
M \longrightarrow \text{Lift}(x, A[M])
$$
；它取 $x$ 到 $\Spec(A[M])$ 的无穷小形变的同构类。把《形式形变理论》
引理 \ref{formal-defos-lemma-linear-functor} 应用于 $\text{Inf}_x$ 和
$T_x$。该引理可用，因为 (RS*) 表明
$$
\textit{Lift}(x, A[M \times N]) =
\textit{Lift}(x, A[M]) \times \textit{Lift}(x, A[N])
$$
是范畴等式（平凡形变也彼此对应）。

\medskip\noindent
设 $(x, A' \to A)$ 是形变情形。考虑由规则
$g(a_1, a_2) = \overline{a_1} \oplus a_2 - a_1$ 定义的环同态
$g : A' \times_A A' \to A[I]$。存在同构
$$
A' \times_A A' \longrightarrow A' \times_A A[I]
$$
，由 $(a_1, a_2) \mapsto (a_1, g(a_1, a_2))$ 给出。此同构与到第一因子
$A'$ 的投影交换，因而也与到 $A$ 的投影交换。因此，两次应用 (RS*)，
得到范畴等价
\begin{align*}
\textit{Lift}(x, A') \times \textit{Lift}(x, A')
& =
\textit{Lift}(x, A' \times_A A') \\
& =
\textit{Lift}(x, A' \times_A A[I]) \\
& =
\textit{Lift}(x, A') \times \textit{Lift}(x, A[I])
\end{align*}
使用这些映射以及到最后一个积的末因子的投影，得到“差映射”
$$
\text{Inf}(x'/x) \times  \text{Inf}(x'/x)
\longrightarrow
\text{Inf}_x(I)
\quad\text{且}\quad
\text{Lift}(x, A') \times \text{Lift}(x, A')
\longrightarrow
T_x(I)
$$
这些差映射满足传递律
“$(x'_1 - x'_2) + (x'_2 - x'_3) = x'_1 - x'_3$”，因为
$$
\xymatrix{
A' \times_A A' \times_A A'
\ar[rrrrr]_-{(a_1, a_2, a_3) \mapsto (g(a_1, a_2), g(a_2, a_3))}
\ar[rrrrrd]_{(a_1, a_2, a_3) \mapsto g(a_1, a_3)} & & & & &
A[I] \times_A A[I] = A[I \times I] \ar[d]^{+} \\
& & & & & A[I]
}
$$
交换。把上面的一串等价取逆，便得到一个作用；只要
$\text{Inf}(x'/x)$，相应地 $\text{Lift}(x, A')$ 非空，该作用就是自由且
传递的。注意，$\text{Inf}(x'/x)$ 是群，故总是非空。
\end{proof}

\begin{remark}[函子性]
\label{artin-remark-functoriality}
假设和记号如引理 \ref{artin-lemma-properties-lift-RS-star}。设 $A \to B$ 是
环同态，且 $y = x|_{\Spec(B)}$。设 $M \in \text{Mod}_A$、
$N \in \text{Mod}_B$。
% P11-E0259: the frozen source omits "be" in "let M -> N an A-linear map".
再设 $M \to N$ 是 $A$-线性映射。则存在典范映射
$\text{Inf}_x(M) \to \text{Inf}_y(N)$ 和
$T_x(M) \to T_y(N)$，因为存在拉回函子
$$
\textit{Lift}(x, A[M]) \to \textit{Lift}(y, B[N])
$$
，它来自环同态 $A[M] \to B[N]$。类似地，给定形变情形的态射
$(y, B' \to B) \to (x, A' \to A)$，得到拉回函子
$\textit{Lift}(x, A') \to \textit{Lift}(y, B')$。由于 $\text{Inf}_x$
和 $T_x$ 上的作用、加法和标量乘法的构造只使用提升范畴中的态射（参见
《形式形变理论》引理 \ref{formal-defos-lemma-linear-functor} 的证明），
上述构造具有函子性。换言之，得到 $A$-线性映射
$$
\text{Inf}_x(M) \to \text{Inf}_y(N)
\quad\text{且}\quad
T_x(M) \to T_y(N)
$$
，使得下列图
$$
\vcenter{
\xymatrix{
\text{Inf}_y(J) \ar[r] & \text{Inf}(y'/y) \\
\text{Inf}_x(I) \ar[r] \ar[u] & \text{Inf}(x'/x) \ar[u]
}
}
\quad\text{且}\quad
\vcenter{
\xymatrix{
T_y(J) \times \text{Lift}(y, B') \ar[r] & \text{Lift}(y, B') \\
T_x(I) \times \text{Lift}(x, A') \ar[r] \ar[u] & \text{Lift}(x, A') \ar[u]
}
}
$$
交换。这里 $I = \Ker(A' \to A)$、$J = \Ker(B' \to B)$；$x'$ 是 $x$
到 $A'$ 的提升（它不一定存在），而 $y' = x'|_{\Spec(B')}$。
\end{remark}

\begin{remark}[自同构]
\label{artin-remark-automorphisms}
假设和记号如引理 \ref{artin-lemma-properties-lift-RS-star}。设 $x', x''$ 是 $x$
到 $A'$ 的提升。则有复合映射
$$
\text{Inf}(x'/x) \times
\Mor_{\textit{Lift}(x, A')}(x', x'') \times \text{Inf}(x''/x)
\longrightarrow
\Mor_{\textit{Lift}(x, A')}(x', x'').
$$
由于 $\textit{Lift}(x, A')$ 是群胚，若
$\Mor_{\textit{Lift}(x, A')}(x', x'')$ 非空，则这定义了
$\text{Inf}(x'/x)$ 在 $\Mor_{\textit{Lift}(x, A')}(x', x'')$ 上的单传递
左作用，以及 $\text{Inf}(x''/x)$ 的单传递右作用。现在，引理断言
$\text{Inf}(x'/x) = \text{Inf}_x(I) = \text{Inf}(x''/x)$。我们断言，
经这些认同后，上述两个作用一致。事实上，或者 $x' \not \cong x''$，
此时断言显然；或者 $x' \cong x''$，此时可假设 $x'' = x'$，而结论来自
$\text{Inf}(x'/x)$ 的交换性。特别地，得到良定义的作用
$$
\text{Inf}_x(I) \times \Mor_{\textit{Lift}(x, A')}(x', x'')
\longrightarrow
\Mor_{\textit{Lift}(x, A')}(x', x'')
$$
；只要 $\Mor_{\textit{Lift}(x, A')}(x', x'')$ 非空，此作用就单传递。
\end{remark}

\begin{remark}
\label{artin-remark-short-exact-sequence-thickenings}
设 $S$ 是概形，$\mathcal{X}$ 是 $(\Sch/S)_{fppf}$ 上的群胚纤维化范畴，
$A$ 是 $S$-代数。可以如下定义形变情形的{\it 短正合列}
$$
(x, A_1' \to A) \to (x, A_2' \to A) \to (x, A_3' \to A)
$$
：要求核 $I_i = \Ker(A_i' \to A)$ 之间的相应映射
% P11-E0260: the frozen source omits "that" after "we ask".
给出 $A$-模的短正合列
$$
0 \to I_3 \to I_2 \to I_1 \to 0
$$
。注意，在这种情形下，映射 $A_3' \to A_1'$ 分解经过 $A$，故有典范同构
$A_1' = A[I_1]$。
\end{remark}

\begin{lemma}
\label{artin-lemma-ses-inf-and-T}
设 $S$ 是概形。设 $p : \mathcal{X} \to \mathcal{Y}$ 和
$q : \mathcal{Z} \to \mathcal{Y}$ 是 $(\Sch/S)_{fppf}$ 上群胚纤维化
范畴的 $1$-态射。假设 $\mathcal{X}$、$\mathcal{Y}$、$\mathcal{Z}$
满足 (RS*)。设 $A$ 是 $S$-代数，$w$ 是
$\mathcal{W} = \mathcal{X} \times_\mathcal{Y} \mathcal{Z}$ 位于 $A$
之上的对象。
% P11-E0261: the frozen source omits "by" in the three-object naming construction.
把从 $w$ 得到的 $\mathcal{X}, \mathcal{Y}, \mathcal{Z}$ 的对象分别记为
$x, y, z$。对任意 $A$-模 $M$，有 $A$-模的 $6$-项正合列
$$
\xymatrix{
0 \ar[r] &
\text{Inf}_w(M) \ar[r] &
\text{Inf}_x(M) \oplus \text{Inf}_z(M) \ar[r] &
\text{Inf}_y(M) \ar[lld] \\
 &
T_w(M) \ar[r] &
T_x(M) \oplus T_z(M) \ar[r] &
T_y(M)
}
$$
\end{lemma}

\begin{proof}
由引理 \ref{artin-lemma-fibre-product-RS-star}，$\mathcal{W}$ 满足 (RS*)，
因而 $T_w(M)$ 和 $\text{Inf}_w(M)$ 有定义。横向箭头利用引理
\ref{artin-lemma-functoriality} 的函子性定义。

\medskip\noindent
定义“边界”映射 $\delta : \text{Inf}_y(M) \to T_w(M)$。选取同构
$p(x) \to y$ 和 $y \to q(z)$，使得在《范畴》引理
\ref{categories-lemma-2-product-fibred-categories}，更确切地说《范畴》引理
\ref{categories-lemma-2-product-categories-over-C} 对 $2$-纤维积的描述中，
$w = (x, z, p(x) \to y \to q(z))$。令 $x', y', z', w'$ 分别表示
$x, y, z, w$ 在 $A[M]$ 上的平凡形变。通过拉回得到同构
$y' \to p(x')$ 和 $q(z') \to y'$。元素 $\alpha \in \text{Inf}_y(M)$
等同于 $A[M]$ 上的自同构 $\alpha : y' \to y'$，它在 $A$ 上限制为 $y$
的恒等态射。因此，令
$$
\delta(\alpha) =
(x', z', p(x') \to y' \xrightarrow{\alpha} y' \to q(z'))
$$
，便得到 $T_w(M)$ 的对象。由《形式形变理论》引理
\ref{formal-defos-lemma-morphism-linear-functors}，这是 $A$-模的映射。

\medskip\noindent
证明的其余部分与《形式形变理论》引理
\ref{formal-defos-lemma-deformation-categories-fiber-product-morphisms}
的证明完全相同。
\end{proof}

\begin{remark}[与前述切空间的相容性]
\label{artin-remark-compare-deformation-spaces}
设 $S$ 是局部 Noether 概形，$\mathcal{X}$ 是 $(\Sch/S)_{fppf}$ 上的
群胚纤维化范畴。假设 $\mathcal{X}$ 具有 (RS*)。设 $k$ 是 $S$ 上的
有限型域，$x_0$ 是 $\mathcal{X}$ 位于 $\Spec(k)$ 上的对象。则有
$k$-向量空间的等式
$$
T\mathcal{F}_{\mathcal{X}, k, x_0} = T_{x_0}(k)
\quad\text{且}\quad
\text{Inf}(\mathcal{F}_{\mathcal{X}, k, x_0}) =
\text{Inf}_{x_0}(k)
$$
其中等号左边的空间由 (\ref{artin-equation-tangent-space}) 和
(\ref{artin-equation-infinitesimal-automorphisms}) 给出，右边的空间由引理
\ref{artin-lemma-properties-lift-RS-star} 给出。
\end{remark}

\begin{remark}[典范元素]
\label{artin-remark-canonical-element}
假设和记号如引理 \ref{artin-lemma-properties-lift-RS-star}。选取仿射开集
$\Spec(\Lambda) \subset S$，使得 $\Spec(A) \to S$ 对应于环同态
$\Lambda \to A$。考虑环同态
$$
A \longrightarrow A[\Omega_{A/\Lambda}],
\quad
a \longmapsto (a, \text{d}_{A/\Lambda}(a))
$$
把 $x$ 沿相应态射 $\Spec(A[\Omega_{A/\Lambda}]) \to \Spec(A)$ 拉回，
得到 $x$ 在 $A[\Omega_{A/\Lambda}]$ 上的形变 $x_{can}$。称之为
{\it 典范元素}
$$
x_{can} \in T_x(\Omega_{A/\Lambda}) = \text{Lift}(x, A[\Omega_{A/\Lambda}]).
$$
其次，假设 $\Lambda$ 是 Noether 环，且 $\Lambda \to A$ 有限型。令
$k = \kappa(\mathfrak p)$ 为 $U = \Spec(A)$ 的有限型点 $u_0$ 处的剩余域。
令 $x_0 = x|_{u_0}$。由 (RS*) 以及 $A[k] = A \times_k k[k]$，空间
$T_x(k)$ 是形变函子 $\mathcal{F}_{\mathcal{X}, k, x_0}$ 的切空间。借助
$$
T\mathcal{F}_{U, k, u_0} =
\text{Der}_\Lambda(A, k) = \Hom_A(\Omega_{A/\Lambda}, k)
$$
（参见《形式形变理论》例
\ref{formal-defos-example-tangent-space-prorepresentable-functor}）和 $T_x$
的函子性，典范元素给出 $U$ 上的对象 $x$ 所诱导的切空间映射。具体地，
$\theta \in T\mathcal{F}_{U, k, u_0}$ 映到
$T_x(k) = T\mathcal{F}_{\mathcal{X}, k, x_0}$ 中的 $T_x(\theta)(x_{can})$。
\end{remark}

\begin{remark}[典范自同构]
\label{artin-remark-canonical-isomorphism}
设 $S$ 是局部 Noether 概形，$\mathcal{X}$ 是 $(\Sch/S)_{fppf}$ 上的
群胚纤维化范畴。假设 $\mathcal{X}$ 满足条件 (RS*)。设 $A$ 是 $S$-代数，
$\Spec(A) \to S$ 映入某个仿射开集，并设 $x, y$ 是 $\mathcal{X}$ 位于
$\Spec(A)$ 上的对象。再设 $A \to B$ 是环同态，而
$\alpha : x|_{\Spec(B)} \to y|_{\Spec(B)}$ 是 $\mathcal{X}$ 位于
$\Spec(B)$ 上的态射。考虑环同态
$$
B \longrightarrow B[\Omega_{B/A}],
\quad
b \longmapsto (b, \text{d}_{B/A}(b))
$$
把 $\alpha$ 沿相应态射 $\Spec(B[\Omega_{B/A}]) \to \Spec(B)$ 拉回，
得到 $B[\Omega_{B/A}]$ 上 $x$ 与 $y$ 的拉回之间的态射 $\alpha_{can}$。
% P11-E0262: the frozen source uses the noun "pullback" as a verb instead of "pull back".
另一方面，可以沿如下态射拉回 $\alpha$：
$\Spec(B[\Omega_{B/A}]) \to \Spec(B)$，它对应于 $B$ 到
$B[\Omega_{B/A}]$ 第一直和项的嵌入。由评注 \ref{artin-remark-automorphisms}
中的讨论，可以取差
$$
\varphi(x, y, \alpha) = \alpha_{can} - \alpha|_{\Spec(B[\Omega_{B/A}])} \in
\text{Inf}_{x|_{\Spec(B)}}(\Omega_{B/A}).
$$
称之为{\it 典范自同构}。它依赖于所有数据 $A$、$x$、$y$、$A \to B$
和 $\alpha$。
\end{remark}





\section{障碍理论}
\label{artin-section-obstruction-theory}

\noindent
本节说明何谓障碍理论。与无穷小形变空间和无穷小自同构空间不同，
障碍理论是一项附加数据。这一表述受引理
\ref{artin-lemma-properties-lift-RS-star} 和评注 \ref{artin-remark-functoriality}
的结果启发。

\begin{definition}
\label{artin-definition-obstruction-theory}
设 $S$ 是局部 Noether 基概形，$\mathcal{X}$ 是 $(\Sch/S)_{fppf}$ 上的
群胚纤维化范畴。一个{\it 障碍理论}由下列数据给出：
\begin{enumerate}
\item 对每个满足 $\Spec(A) \to S$ 映入某个仿射开集的 $S$-代数 $A$，
以及 $\mathcal{X}$ 位于 $\Spec(A)$ 上的每个对象 $x$，给定一个 $A$-线性函子
$$
\mathcal{O}_x : \text{Mod}_A \to \text{Mod}_A
$$
，称为{\it 障碍模}函子；
% P11-E0263: the frozen source omits the existential verb before the induced map.
\item 对如 (1) 中的 $(x, A)$、环同态 $A \to B$、
$M \in \text{Mod}_A$、$N \in \text{Mod}_B$ 和 $A$-线性映射
$M \to N$，给定诱导的 $A$-线性映射
$\mathcal{O}_x(M) \to \mathcal{O}_y(N)$，其中 $y = x|_{\Spec(B)}$；
\item 对每个形变情形 $(x, A' \to A)$，给定一个{\it 障碍}元素
$o_x(A') \in \mathcal{O}_x(I)$，其中 $I = \Ker(A' \to A)$。
\end{enumerate}
这些数据须满足下列条件：
\begin{enumerate}
\item[(i)] 函子性映射使障碍模成为从三元组 $(x, A, M)$ 的范畴到集合范畴的
函子；
\item[(ii)] 对每个形变情形的态射
$(y, B' \to B) \to (x, A' \to A)$，元素 $o_x(A')$ 映到
$o_y(B')$；
\item[(iii)] 对每个形变情形 $(x, A' \to A)$，有
$$
\text{Lift}(x, A') \not = \emptyset
\Leftrightarrow
o_x(A') = 0
$$
。
\end{enumerate}
\end{definition}

\noindent
最后一个条件解释了上述术语。
% P11-E0264: the frozen source calls O_x(A') the obstruction module although the datum is O_x(I).
模 $\mathcal{O}_x(A')$ 称为{\it 障碍模}，元素 $o_x(A')$ 称为{\it 障碍}。
大多数障碍理论还有附加性质；要使它们有用，还需施加附加条件。
此外，这只是一种示例性定义；例如，定义中可以只考虑 $S$ 上有限型的形变情形。

\medskip\noindent
引入障碍理论的主要理由之一，是检验通用性的开性。下面的引理
\ref{artin-lemma-get-openness-obstruction-theory} 就是这类结果的一个例子。
这样做的最初思想出自 Artin；参见引言中提到的 Artin 论文。例如，Flenner
\cite{Flenner}、Hall \cite{Hall-coherent}、Hall 和 Rydh
\cite{rydh_axioms}、Olsson \cite{olsson_deformation}、Olsson 和 Starr
\cite{olsson-starr} 以及 Lieblich \cite{lieblich-complexes} 的工作都采用了
这一思想（文献次序随意）。此外，对特定的群胚纤维化范畴，作者往往会发展少量
适合手头问题的理论。
% P11-E0265: the frozen source retains a literal future-reference placeholder.
我们稍后将发展这一理论（此处插入将来的文献引用）。

\begin{lemma}
\label{artin-lemma-get-openness-obstruction-theory}
\begin{reference}
这是 \cite[Theorem 4.4]{Hall-coherent}。
\end{reference}
设 $S$ 是局部 Noether 概形，$\mathcal{X}$ 是 $(\Sch/S)_{fppf}$ 上的
群胚纤维化范畴。假设
\begin{enumerate}
\item $\Delta : \mathcal{X} \to \mathcal{X} \times \mathcal{X}$
可由代数空间表示；
\item $\mathcal{X}$ 具有 (RS*)；
\item $\mathcal{X}$ 保极限；
\item 存在一个障碍理论\footnote{分析证明可见，实际上只需对满足
$B = A$ 和 $y = x$ 的映射
$(y, B' \to B) \to (x, A' \to A)$，检验定义
\ref{artin-definition-obstruction-theory} 中障碍类的函子性 (ii)。}；
\item 对 $\mathcal{X}$ 位于 $\Spec(A)$ 上的对象 $x$ 以及 $A$-模
$M_n$（$n \geq 1$），有
\begin{enumerate}
\item $T_x(\prod M_n) = \prod T_x(M_n)$,
\item $\mathcal{O}_x(\prod M_n) \to \prod \mathcal{O}_x(M_n)$
是单射。
\end{enumerate}
\end{enumerate}
则 $\mathcal{X}$ 满足通用性的开性。
\end{lemma}

\begin{proof}
我们通过验证引理 \ref{artin-lemma-SGE-implies-openness-versality} 的条件 (4)
来证明。设 $(\xi_n)$ 和 $(R_n)$ 如评注 \ref{artin-remark-strong-effectiveness}，
且对所有 $m \geq n$，$\Ker(R_m \to R_n)$ 都是平方零理想。令
$A = R_1$、$x = \xi_1$。
% P11-E0266: the frozen source uses the incomplete naming construction "Denote M_n = ...".
记 $M_n = \Ker(R_n \to R_1)$。则 $M_n$ 是 $A$-模。令
$R = \lim R_n$，并令
$$
% P11-E0270: the frozen source omits the comma between r_3 and the ellipsis.
\tilde R = \{(r_1, r_2, r_3 \ldots) \in \prod R_n
\text{ such that all have the same image in }A\}
$$
则 $\tilde R \to A$ 是满射，其核为 $M = \prod M_n$。存在映射
$R \to \tilde R$ 和映射 $\tilde R \to A[M]$；后者由
$(r_1, r_2, r_3, \ldots) \mapsto
(r_1, r_2 - r_1, r_3 - r_2, \ldots)$ 给出。二者合在一起给出形变情形的短正合列
$$
(x, R \to A) \to (x, \tilde R \to A) \to (x, A[M])
$$
；参见评注 \ref{artin-remark-short-exact-sequence-thickenings}。相应的核序列
$0 \to \lim M_n \to M \to M \to 0$，正是计算模系统 $(M_n)$ 的极限的
典范序列。

\medskip\noindent
设 $o_x(\tilde R) \in \mathcal{O}_x(M)$ 为障碍元素。由于已有提升
$\xi_n$，可知 $o_x(\tilde R)$ 在 $\mathcal{O}_x(M_n)$ 中映为零。由假设
(5)(b)，得到 $o_x(\tilde R) = 0$。选取 $x$ 到 $\Spec(\tilde R)$ 的提升
$\tilde \xi$，并令 $\tilde \xi_n$ 为 $\tilde \xi$ 到 $\Spec(R_n)$ 的限制。
% P11-E0271: the frozen source uses singular "There exists" with plural "elements".
由引理 \ref{artin-lemma-properties-lift-RS-star} 的 (2)(b)，存在元素
$t_n \in T_x(M_n)$，使得 $t_n \cdot \tilde \xi_n = \xi_n$。由假设
(5)(a)，可以找到 $t \in T_x(M)$，它在 $T_x(M_n)$ 中映到 $t_n$。
把 $\tilde \xi$ 替换为 $t \cdot \tilde \xi$ 后，对所有 $n$，
$\tilde \xi$ 在 $\Spec(R_n)$ 上限制为 $\xi_n$。特别地，由于
$\xi_{n + 1}$ 在 $\Spec(R_n)$ 上限制为 $\xi_n$，$\tilde \xi$ 到
$\Spec(A[M])$ 的限制 $\overline{\xi}$ 具有如下性质：对所有 $n$，它在
$\Spec(A[M_n])$ 上限制为平凡形变。因此由假设 (5)(a)，
$\overline{\xi}$ 是 $x$ 的平凡形变。把公理 (RS*) 应用于
$R = \tilde R \times_{A[M]} A$，便知 $\tilde \xi$ 是 $x$ 在 $R$ 上某个形变
$\xi$ 的拉回。证明完毕。
\end{proof}

\begin{example}
\label{artin-example-global-sections}
设 $S = \Spec(\Lambda)$，其中 $\Lambda$ 是 Noether 环。设 $W \to S$
为概形态射，$\mathcal{F}$ 为在 $S$ 上平坦的拟凝聚 $\mathcal{O}_W$-模。
考虑函子
$$
F : (\Sch/S)_{fppf}^{opp} \longrightarrow \textit{Sets},
\quad
T/S \longrightarrow H^0(W_T, \mathcal{F}_T)
$$
，其中 $W_T = T \times_S W$ 是基变换，$\mathcal{F}_T$ 是 $\mathcal{F}$
到 $W_T$ 的拉回。若 $T = \Spec(A)$，则记 $W_T = W_A$，其余类同。
设 $\mathcal{X} \to (\Sch/S)_{fppf}$ 为与 $F$ 相联系的群胚纤维化范畴。
则 $\mathcal{X}$ 具有一个障碍理论。具体地，
\begin{enumerate}
\item 给定 $\Lambda$-代数 $A$ 和
$x \in H^0(W_A, \mathcal{F}_A)$，令
$\mathcal{O}_x(M) = H^1(W_A, \mathcal{F}_A \otimes_A M)$,
% P11-E0267: the frozen source places o_x(A') in O_x(A), while the boundary target is O_x(I).
\item 给定形变情形 $(x, A' \to A)$，令
$o_x(A') \in \mathcal{O}_x(A)$ 为 $x$ 在边界映射
$$
H^0(W_A, \mathcal{F}_A) \longrightarrow H^1(W_A, \mathcal{F}_A \otimes_A I)
$$
下的像；该边界映射来自模的短正合列
$$
0 \to \mathcal{F}_A \otimes_A I \to
\mathcal{F}_{A'} \to \mathcal{F}_A \to 0.
$$
\end{enumerate}
我们略去了一些细节，尤其是上述短正合列的构造（它使用 $W_A$ 和
$W_{A'}$ 具有相同底拓扑空间这一事实），以及为何 $\mathcal{F}$ 在 $S$
上的平坦性蕴含上述序列短正合的说明。
\end{example}

\begin{example}[关键例]
\label{artin-example-key}
设 $S = \Spec(\Lambda)$，其中 $\Lambda$ 是 Noether 环。取
$\mathcal{X} = (\Sch/X)_{fppf}$，其中 $X = \Spec(R)$ 且
$R = \Lambda[x_1, \ldots, x_n]/J$。朴素余切复形 $\NL_{R/\Lambda}$
（典范地）同伦等价于
$$
J/J^2
\longrightarrow
\bigoplus\nolimits_{i = 1, \ldots, n} R\text{d}x_i,
$$
参见《代数》引理 \ref{algebra-lemma-NL-homotopy}。考虑形变情形
$(x, A' \to A)$。
% P11-E0268: the frozen source omits "by" in "Denote I the kernel".
把 $A' \to A$ 的核记为 $I$。对象 $x$ 对应于 $(a_1, \ldots, a_n)$，
其中 $a_i \in A$，且对所有 $f \in J$，$f(a_1, \ldots, a_n) = 0$ 在
$A$ 中成立。令
\begin{align*}
% P11-E0269: the frozen source writes O_x(A') although the displayed module depends on I.
\mathcal{O}_x(A')
& =
\Hom_R(J/J^2, I)/\Hom_R(R^{\oplus n}, I) \\
& =
\Ext^1_R(\NL_{R/\Lambda}, I) \\
& =
\Ext^1_A(\NL_{R/\Lambda} \otimes_R A, I).
\end{align*}
选取 $a_i' \in A'$，使其在 $A$ 中的像为 $a_i$。则 $o_x(A')$ 是映射
$J/J^2 \to I$ 的类；该映射把 $f \in J$ 送到
$f(a_1', \ldots, a'_n) \in I$。我们略去 $o_x(A')$ 与选取无关的验证。
显然，若 $o_x(A') = 0$，则该映射可以提升。最后，函子性是直接的。
于是得到一个障碍理论。注意，$o_x(A')$ 还可稍微更典范地描述为 $D(A)$
中的复合
$$
\NL_{R/\Lambda} \to \NL_{A/\Lambda} \to \NL_{A/A'} = I[1]
$$
；最后一个认同参见《代数》引理 \ref{algebra-lemma-NL-surjection}。
\end{example}









\section{朴素障碍理论}
\label{artin-section-naive-obstruction-theory}

\noindent
本节标题指的是：本节将使用朴素余切复形。设 $(x, A' \to A)$ 是给定的
局部 Noether 概形 $S$ 上群胚纤维化范畴的一个形变情形。关键例
\ref{artin-example-key} 表明，任何障碍理论都应当与 $D(A)$ 中以 $A$ 的朴素余切
复形为目标的映射密切相关。展开这一思路，引理
\ref{artin-lemma-characterize-versal} 给出通用性的判别准则，进而在引理
\ref{artin-lemma-openness} 中得到通用性的开性判据。为了进一步形式化这个概念，
我们在定义 \ref{artin-definition-naive-obstruction-theory} 中引入朴素障碍理论的概念。

\medskip\noindent
下文使用《代数》一节 \ref{algebra-section-netherlander} 中定义的朴素余切复形。
特别地，若 $A' \to A$ 是 $\Lambda$-代数的满射，且其平方零核为 $I$，则有映射
$$
\NL_{A'/\Lambda} \to \NL_{A/\Lambda} \to \NL_{A/A'}
$$
，其复合同伦等价于零（参见《代数》评注
\ref{algebra-remark-composition-homotopy-equivalent-to-zero}）。由于这里使用的是
朴素余切复形而非完整余切复形，一般而言这并不构成可区别三角。存在同伦等价
$\NL_{A/A'} \to I[1]$（即把 $I$ 置于次数 $-1$ 所得的复形；参见《代数》
引理 \ref{algebra-lemma-NL-surjection}）。最后，注意存在典范映射
$\NL_{A/\Lambda} \to \Omega_{A/\Lambda}$。

\begin{lemma}
\label{artin-lemma-compute-ext-into-field}
设 $A \to k$ 是环同态，其中 $k$ 是域。设 $E \in D^-(A)$。则
$\Ext^i_A(E, k) = \Hom_k(H^{-i}(E \otimes^\mathbf{L} k), k)$。
\end{lemma}

\begin{proof}
从略。提示：把 $E$ 替换为由自由 $A$-模组成的上有界复形，并计算等式两边。
\end{proof}

\begin{lemma}
\label{artin-lemma-construct-essential-surjection}
设 $\Lambda \to A \to k$ 是 Noether 环的有限型环同态，且对 $A$ 的某个素理想
$\mathfrak p$ 有 $k = \kappa(\mathfrak p)$。
% P11-E0272: the frozen source omits the article in "be morphism of D^-(A)".
设 $\xi : E \to \NL_{A/\Lambda}$ 是 $D^{-}(A)$ 中的态射，并且
$H^{-1}(\xi \otimes^{\mathbf{L}} k)$ 不是满射。则存在 $\Lambda$-代数的满射
$A' \to A$，使得
\begin{enumerate}
\item[(a)] $I = \Ker(A' \to A)$ 平方为零，并且作为 $A$-模同构于 $k$；
\item[(b)] $\Omega_{A'/\Lambda} \otimes k = \Omega_{A/\Lambda} \otimes k$；
\item[(c)] $E \to \NL_{A/A'}$ 为零。
\end{enumerate}
\end{lemma}

\begin{proof}
设 $f \in A$、$f \not \in \mathfrak p$。假设对诱导映射
$E \otimes_A A_f \to \NL_{A_f/\Lambda}$，$A'' \to A_f$ 满足 (a)、(b)、(c)；
参见《代数》引理 \ref{algebra-lemma-localize-NL}。则可令
$A' = A'' \times_{A_f} A$，从而得到一个解。确实，
% P11-E0273: the frozen source identifies the kernel with Ker(A'' -> A), although only A'' -> A_f is given.
$A' \to A$ 显然满足 (a)，因为
$\Ker(A' \to A) = \Ker(A'' \to A) = I$。选取提升 $f$ 的
$f'' \in A''$。则 $A'$ 在 $(f'', f)$ 处的局部化同构于 $A''$（例如由
《代数进阶》引理 \ref{more-algebra-lemma-diagram-localize}）。因此 (b) 和 (c)
对 $A'$ 也显然成立。由此可见，可以把 $A$ 替换为局部化 $A_f$（有限多次）。
特别地，在作此替换后，可以假设 $\mathfrak p$ 是 $A$ 的极大理想；参见
《态射》引理 \ref{morphisms-lemma-point-finite-type}。

\medskip\noindent
选取表示 $A = \Lambda[x_1, \ldots, x_n]/J$。则 $\NL_{A/\Lambda}$
（典范地）同伦等价于
$$
J/J^2
\longrightarrow
\bigoplus\nolimits_{i = 1, \ldots, n} A\text{d}x_i,
$$
参见《代数》引理 \ref{algebra-lemma-NL-homotopy}。必要时作局部化后（使用
Nakayama 引理），可以选取 $J$ 的生成元 $f_1, \ldots, f_m$，使得
$f_j \otimes 1$ 构成 $J/J^2 \otimes_A k$ 的一组基。进一步重新编号后，
可假设 $\text{d}f_1, \ldots, \text{d}f_r$ 的像构成映射
$J/J^2 \otimes k \to \bigoplus k\text{d}x_i$ 的像的一组基，并且
$\text{d}f_{r + 1}, \ldots, \text{d}f_m$ 在
$\bigoplus k\text{d}x_i$ 中映为零。在这些选取下，空间
$$
H^{-1}(\NL_{A/\Lambda} \otimes^{\mathbf{L}}_A k) =
H^{-1}(\NL_{A/\Lambda} \otimes_A k)
$$
以 $f_{r + 1} \otimes 1, \ldots, f_m \otimes 1$ 为一组基。再次换基后，
可以假设 $H^{-1}(\xi \otimes^{\mathbf{L}} k)$ 的像包含在
$f_{r + 1} \otimes 1, \ldots, f_{m - 1} \otimes 1$ 的 $k$-线性张成中。令
$$
A' = \Lambda[x_1, \ldots, x_n]/(f_1, \ldots, f_{m - 1}, \mathfrak pf_m)
$$
按构造，$A' \to A$ 满足 (a)。由于 $\text{d}f_m$ 在
$\bigoplus k\text{d}x_i$ 中映为零，可知 (b) 成立。最后，按构造，诱导映射
$E \to \NL_{A/A'} = I[1]$ 在
$H^{-1}(E \otimes_A^\mathbf{L} k) \to I \otimes_A k$ 上诱导零映射。
由引理 \ref{artin-lemma-compute-ext-into-field}，可知该复合为零。
\end{proof}

\noindent
下一个引理是我们的关键技术结果。

\begin{lemma}
\label{artin-lemma-characterize-versal}
设 $S$ 是局部 Noether 概形，$\mathcal{X}$ 是 $(\Sch/S)_{fppf}$ 上满足
(RS*) 的群胚纤维化范畴。设 $U = \Spec(A)$ 是 $S$ 上有限型仿射概形，
且它映入仿射开集 $\Spec(\Lambda)$。设 $x$ 是 $\mathcal{X}$ 位于 $U$ 上的对象，
$\xi : E \to \NL_{A/\Lambda}$ 是 $D^{-}(A)$ 中的态射。假设
\begin{enumerate}
\item[(i)] 对每个形变情形 $(x, A' \to A)$，$x$ 可提升到 $\Spec(A')$，
当且仅当 $E \to \NL_{A/\Lambda} \to \NL_{A/A'}$ 为零；
\item[(ii)] 存在函子同构
$T_x(-) \to \Ext^0_A(E, -)$
，使得 $E \to \NL_{A/\Lambda} \to \Omega^1_{A/\Lambda}$ 对应于典范元素
（参见评注 \ref{artin-remark-canonical-element}）。
\end{enumerate}
设 $u_0 \in U$ 是剩余域为 $k = \kappa(u_0)$ 的有限型点。考虑下列陈述：
\begin{enumerate}
\item $x$ 在 $u_0$ 处通用；
\item $\xi : E \to \NL_{A/\Lambda}$ 诱导满射
$H^{-1}(E \otimes_A^{\mathbf{L}} k) \to
H^{-1}(\NL_{A/\Lambda} \otimes_A^{\mathbf{L}} k)$
和单射
$H^0(E \otimes_A^{\mathbf{L}} k) \to
H^0(\NL_{A/\Lambda} \otimes_A^{\mathbf{L}} k)$.
\end{enumerate}
则总有 (2) $\Rightarrow$ (1)；若 $u_0$ 是闭点，则还有
(1) $\Rightarrow$ (2)。
\end{lemma}

\begin{proof}
令 $\mathfrak p = \Ker(A \to k)$ 为与 $u_0$ 对应的素理想。

\medskip\noindent
% P11-E0274: the frozen source omits "is" in "Assume that x versal".
假设 $x$ 在 $u_0$ 处通用，且 $u_0$ 是 $U$ 的闭点。若
$H^{-1}(\xi \otimes_A^{\mathbf{L}} k)$ 不是满射，则取引理
\ref{artin-lemma-construct-essential-surjection} 中核为 $I$ 的扩张 $A' \to A$。
因为 $u_0$ 是闭点，$I$ 是有限 $A$-模，因而 $A'$ 是有限型
$\Lambda$-代数（若 $u_0$ 非闭，这一点会失效）。特别地，$A'$ 是 Noether 环。
由 $A'$ 的性质 (c) 和 $\xi$ 的性质 (i)，$x$ 可提升为 $A'$ 上的对象 $x'$。
% P11-E0275: the frozen source omits the article before "kernel".
令 $\mathfrak p' \subset A'$ 为到 $k$ 的满同态之核。由 Artin--Rees 引理
（《代数》引理 \ref{algebra-lemma-Artin-Rees}），存在 $n > 1$，使得
$(\mathfrak p')^n \cap I = 0$。于是
$$
B' = A'/(\mathfrak p')^n \longrightarrow A/\mathfrak p^n = B
$$
是 Artin 局部环的小的本质扩张；参见《形式形变理论》引理
\ref{formal-defos-lemma-essential-surjection}。另一方面，由于 $x$ 在 $u_0$
处通用，而 $x'|_{\Spec(B')}$ 是 $x|_{\Spec(B)}$ 的提升，存在整数
$m \geq n$ 和映射 $q : A/\mathfrak p^m \to B'$，使得复合
$A/\mathfrak p^m \to B' \to B$ 是商映射。由于 $B'$ 的极大理想的
$n$ 次幂为零，$q$ 分解经过 $B$，这与 $B' \to B$ 是本质满射矛盾。
这一矛盾说明 $H^{-1}(\xi \otimes_A^{\mathbf{L}} k)$ 是满射。

\medskip\noindent
% P11-E0276: the frozen source again omits "is" in "Assume that x versal".
假设 $x$ 在 $u_0$ 处通用。由引理 \ref{artin-lemma-compute-ext-into-field}，映射
$H^0(\xi \otimes_A^{\mathbf{L}} k)$ 对偶于映射
$\Ext^0_A(\NL_{A/\Lambda}, k) \to \text{Ext}^0_A(E, k)$。注意
$$
\Ext^0_A(\NL_{A/\Lambda}, k) = \text{Der}_\Lambda(A, k)
\quad\text{且}\quad
T_x(k) = \Ext^0_A(E, k)
$$
条件 (ii) 保证映射
$\Ext^0_A(\NL_{A/\Lambda}, k) \to \text{Ext}^0_A(E, k)$ 把 $U$ 在 $u_0$
处的切向量 $\theta$ 送到 $x_0$ 的相应无穷小形变；参见评注
\ref{artin-remark-canonical-element}。因此，若 $x$ 通用，则此映射为满射；参见
《形式形变理论》引理 \ref{formal-defos-lemma-versal-criterion}。于是
$H^0(\xi \otimes_A^{\mathbf{L}} k)$ 是单射。这就证明了当 $u_0$ 是闭点时的
(1) $\Rightarrow$ (2)。

\medskip\noindent
在证明的余下部分，假设
$H^{-1}(E \otimes_A^\mathbf{L} k) \to
H^{-1}(\NL_{A/\Lambda} \otimes_A^\mathbf{L} k)$ 是满射，并且
% P11-E0277: the frozen source omits "is" before "injective".
$H^0(E \otimes_A^\mathbf{L} k) \to
H^0(\NL_{A/\Lambda} \otimes_A^\mathbf{L} k)$
是单射。令 $R = A_\mathfrak p^\wedge$，并令 $\eta$ 为与
$x|_{\Spec(R)}$ 相联系的 $R$ 上形式对象。切空间上的映射
$d\underline{\eta}$ 是满射，因为它被认同为单射
$H^0(E \otimes_A^{\mathbf{L}} k) \to
H^0(\NL_{A/\Lambda} \otimes_A^{\mathbf{L}} k)$
的对偶（参见上一段）。根据《形式形变理论》引理
\ref{formal-defos-lemma-versal-criterion}，只须证明如下断言：设
$C' \to C$ 是剩余域为 $k$ 的有限型 Artin 局部 $\Lambda$-代数的小扩张，
$R \to C$ 是与剩余域认同相容的 $\Lambda$-代数同态。令
$y = x|_{\Spec(C)}$，并设 $y'$ 是 $y$ 到 $C'$ 的提升。要证：可以把
$\Lambda$-代数同态 $R \to C$ 提升为 $R \to C'$。

\medskip\noindent
注意，只须提升 $\Lambda$-代数同态 $A \to C$。令
$I = \Ker(C' \to C)$。注意，$I$ 是 $1$ 维 $k$-向量空间。提升
$A \to C$ 的障碍 $ob$ 是 $\Ext^1_A(\NL_{A/\Lambda}, I)$ 的元素；参见例
\ref{artin-example-key}。由引理 \ref{artin-lemma-compute-ext-into-field} 和我们的假设，
映射 $\xi$ 诱导单射
$$
\Ext^1_A(\NL_{A/\Lambda}, I)
\longrightarrow
\Ext^1_A(E, I)
$$
按 $ob$ 的构造及 (i)，$ob$ 在 $\Ext^1_A(E, I)$ 中的像就是把 $x$ 提升到
$A \times_C C'$ 的障碍。由 (RS*)，$y/C$ 可提升到 $y'/C'$ 这一事实蕴含
$x$ 可提升到 $A \times_C C'$。因此 $ob = 0$，结论得证。
\end{proof}

\noindent
上面的关键引理使我们能在若干情形下推出通用性的开性。

\begin{lemma}
\label{artin-lemma-openness}
设 $S$ 是局部 Noether 概形，$\mathcal{X}$ 是 $(\Sch/S)_{fppf}$ 上满足
(RS*) 的群胚纤维化范畴。设 $U = \Spec(A)$ 是 $S$ 上有限型仿射概形，
且它映入仿射开集 $\Spec(\Lambda)$。设 $x$ 是 $\mathcal{X}$ 位于 $U$ 上的对象，
$\xi : E \to \NL_{A/\Lambda}$ 是 $D^{-}(A)$ 中的态射。假设
\begin{enumerate}
\item[(i)] 对每个形变情形 $(x, A' \to A)$，$x$ 可提升到 $\Spec(A')$，
当且仅当 $E \to \NL_{A/\Lambda} \to \NL_{A/A'}$ 为零；
\item[(ii)] 存在函子同构
$T_x(-) \to \Ext^0_A(E, -)$
，使得 $E \to \NL_{A/\Lambda} \to \Omega^1_{A/\Lambda}$ 对应于典范元素
（参见评注 \ref{artin-remark-canonical-element}）；
\item[(iii)] $E$ 的上同调群都是有限 $A$-模。
\end{enumerate}
若 $x$ 在闭点 $u_0 \in U$ 处通用，则存在开邻域
$u_0 \in U' \subset U$，使得 $x$ 在 $U'$ 的每个有限型点处通用。
\end{lemma}

\begin{proof}
令 $C$ 为 $\xi$ 的锥，于是有可区别三角
$$
E \to \NL_{A/\Lambda} \to C \to E[1]
$$
属于 $D^{-}(A)$。由引理 \ref{artin-lemma-characterize-versal}，$x$ 在 $u_0$
处通用这一假设蕴含 $H^{-1}(C \otimes^\mathbf{L} k) = 0$。由《代数进阶》
引理 \ref{more-algebra-lemma-cut-complex-in-two}，存在不属于与 $u_0$ 对应素理想的
$f \in A$，使得对任意 $A_f$-模 $M$，都有
$H^{-1}(C \otimes^\mathbf{L}_A M) = 0$。再次使用引理
\ref{artin-lemma-characterize-versal}，可知在开集 $D(f) \subset U$ 的所有有限型点处
都有通用性。
\end{proof}

\noindent
上述技术引理提示了下面的定义。

\begin{definition}
\label{artin-definition-naive-obstruction-theory}
设 $S$ 是局部 Noether 基概形，$\mathcal{X}$ 是 $(\Sch/S)_{fppf}$ 上的
群胚纤维化范畴。假设 $\mathcal{X}$ 满足 (RS*)。一个{\it 朴素障碍理论}
由下列数据给出：
\begin{enumerate}
\item
\label{artin-item-map}
对每个满足 $\Spec(A) \to S$ 映入仿射开集
$\Spec(\Lambda) \subset S$ 的 $S$-代数 $A$，以及 $\mathcal{X}$ 位于
$\Spec(A)$ 上的每个对象 $x$，给定对象 $E_x \in D^-(A)$ 和映射
% P11-E0278: the frozen source writes E rather than the just-defined E_x in xi_x.
$\xi_x : E \to \NL_{A/\Lambda}$；
\item
\label{artin-item-inf}
给定如 (\ref{artin-item-map}) 中的 $(x, A)$，存在函子变换
$$
\text{Inf}_x( - ) \to \Ext^{-1}_A(E_x, -)
\quad\text{且}\quad
T_x(-) \to \Ext^0_A(E_x, -)
$$
\item
\label{artin-item-functoriality}
对如 (\ref{artin-item-map}) 中的 $(x, A)$ 和环同态 $A \to B$，令
$y = x|_{\Spec(B)}$，则 $D(A)$ 中存在函子性映射 $E_x \to E_y$。
\end{enumerate}
这些数据须满足下列条件：
\begin{enumerate}
\item[(i)]
在 (\ref{artin-item-functoriality}) 的情形中，图
$$
\xymatrix{
E_y \ar[r]_{\xi_y} & \NL_{B/\Lambda} \\
E_x \ar[u] \ar[r]^{\xi_x} & \NL_{A/\Lambda} \ar[u]
}
$$
在 $D(A)$ 中交换；
\item[(ii)]
给定如 (\ref{artin-item-map}) 中的 $(x, A)$ 和 $A \to B \to C$，令
$y = x|_{\Spec(B)}$、$z = x|_{\Spec(C)}$，则函子性映射 $E_x \to E_y$
与 $E_y \to E_z$ 的复合就是函子性映射 $E_x \to E_z$；
\item[(iii)]
(\ref{artin-item-inf}) 中的映射是同构，并与函子性映射及评注
\ref{artin-remark-functoriality} 中的映射相容；
\item[(iv)]
复合 $E_x \to \NL_{A/\Lambda} \to \Omega_{A/\Lambda}$ 对应于
$T_x(\Omega_{A/\Lambda}) = \Ext^0(E_x, \Omega_{A/\Lambda})$ 的典范元素；
参见评注 \ref{artin-remark-canonical-element}；
\item[(v)]
给定形变情形 $(x, A' \to A)$，其中 $I = \Ker(A' \to A)$，复合
$E_x \to \NL_{A/\Lambda} \to \NL_{A/A'}$ 在
$$
% P11-E0279: the frozen source's first Hom target is NL_{A/Lambda}, contradicting the composition and following terms.
\Hom_A(E_x, \NL_{A/\Lambda}) = \Ext^0_A(E_x, \NL_{A/A'}) =
\Ext^1_A(E_x, I)
$$
中为零，当且仅当 $x$ 可提升到 $A'$。
\end{enumerate}
\end{definition}

\noindent
特别地可见，令 $\mathcal{O}_x( - ) = \Ext^1_A(E_x, -)$，便得到一节
\ref{artin-section-obstruction-theory} 意义下的障碍理论。

\begin{lemma}
\label{artin-lemma-naive-obstruction-theory-qis}
设 $S$ 和 $\mathcal{X}$ 如定义 \ref{artin-definition-naive-obstruction-theory}，
并为 $\mathcal{X}$ 配备一个朴素障碍理论。设 $A \to B$ 和 $y \to x$
如 (\ref{artin-item-functoriality})。设 $k$ 是作为域的 $B$-代数。则函子性映射
$E_x \to E_y$ 诱导双射
$$
H^i(E_x \otimes_A^{\mathbf{L}} k) \to H^i(E_y \otimes_B^{\mathbf{L}} k)
$$
，其中 $i = 0, 1$。
\end{lemma}

\begin{proof}
令 $z = x|_{\Spec(k)}$。由 (RS*)，
$$
\textit{Lift}(x, A[k]) = \textit{Lift}(z, k[k])
\quad\text{且}\quad
\textit{Lift}(y, B[k]) = \textit{Lift}(z, k[k])
$$
，因为 $A[k] = A \times_k k[k]$ 且 $B[k] = B \times_k k[k]$。因此，
朴素障碍理论的性质蕴含函子性映射 $E_x \to E_y$ 对 $i = -1, 0$ 诱导双射
% P11-E0280: the frozen source inconsistently typesets the target Ext operator with \text{Ext}.
$\Ext^i_A(E_x, k) \to \text{Ext}^i_B(E_y, k)$。由引理
\ref{artin-lemma-compute-ext-into-field}，映射
$H^i(E_x \otimes_A^{\mathbf{L}} k) \to H^i(E_y \otimes_B^{\mathbf{L}} k)$,
$i = 0, 1$ 在对偶向量空间上诱导同构，因而本身也是同构。
\end{proof}

\begin{lemma}
\label{artin-lemma-naive-obstruction-theory-gives-openness}
设 $S$ 是局部 Noether 概形。
% P11-E0281: the frozen source gives the base of the fibred category an opp, unlike the chapter's definitions and conventions.
设 $p : \mathcal{X} \to (\Sch/S)_{fppf}^{opp}$ 是群胚纤维化范畴。
假设 $\mathcal{X}$ 满足 (RS*)，并且 $\mathcal{X}$ 具有朴素障碍理论。若对
$A$ 在 $S$ 上有限型的二元组 $(A, x)$，定义
\ref{artin-definition-naive-obstruction-theory} 中的复形 $E_x$ 具有有限生成上同调群，
则 $\mathcal{X}$ 满足通用性的开性。
\end{lemma}

\begin{proof}
设 $U$ 是 $S$ 上局部有限型概形，$x$ 是 $\mathcal{X}$ 位于 $U$ 上的对象，
$u_0$ 是 $U$ 的有限型点，且 $x$ 在 $u_0$ 处通用。首先可把 $U$ 缩小为
仿射概形，使 $u_0$ 为闭点，并使 $U \to S$ 映入仿射开集
$\Spec(\Lambda)$。记 $U = \Spec(A)$。设
$\xi_x : E_x \to \NL_{A/\Lambda}$ 为障碍映射。此时应用引理
\ref{artin-lemma-openness} 即得结论。
\end{proof}









\section{一个对偶概念}
\label{artin-section-dual}

\noindent
设 $(x, A' \to A)$ 是给定的局部 Noether 概形 $S$ 上群胚纤维化范畴
$\mathcal{X}$ 的形变情形。假设 $\mathcal{X}$ 具有障碍理论；参见定义
\ref{artin-definition-obstruction-theory}。实践中，往往有 $A$-模复形
$K^\bullet$ 和函子同构
$$
\text{Inf}_x(-) \to H^0(K^\bullet \otimes_A^\mathbf{L} -),\quad
T_x(-) \to H^1(K^\bullet \otimes_A^\mathbf{L} -),\quad
\mathcal{O}_x(-) \to H^2(K^\bullet \otimes_A^\mathbf{L} -)
$$
本节稍加形式化这一情形，并说明它如何在若干情形下导出通用性的开性。

\begin{example}
\label{artin-example-global-sections-dual}
设 $\Lambda, S, W, \mathcal{F}$ 如例 \ref{artin-example-global-sections}。
假设 $W \to S$ 是真态射，且 $\mathcal{F}$ 凝聚。由《概形的上同调》评注
\ref{coherent-remark-explain-perfect-direct-image}，存在由有限射影
$\Lambda$-模组成的有限复形 $N^\bullet$，它普遍计算 $\mathcal{F}$ 的上同调。
特别地，例 \ref{artin-example-global-sections} 中的障碍空间为
$\mathcal{O}_x(M) = H^1(N^\bullet \otimes_\Lambda M)$。因此，取
$K^\bullet = N^\bullet \otimes_\Lambda A[-1]$，便有
$\mathcal{O}_x(M) = H^2(K^\bullet \otimes_A^\mathbf{L} M)$。
\end{example}

\begin{situation}
\label{artin-situation-dual}
设 $S$ 是局部 Noether 概形，$\mathcal{X}$ 是 $(\Sch/S)_{fppf}$ 上的
群胚纤维化范畴。假设 $\mathcal{X}$ 具有 (RS*)，因而可以谈论函子
$T_x(-)$；参见引理 \ref{artin-lemma-properties-lift-RS-star}。设
$U = \Spec(A)$ 是 $S$ 上有限型仿射概形，且它映入仿射开集
$\Spec(\Lambda)$。设 $x$ 是 $\mathcal{X}$ 位于 $U$ 上的对象。假设给定
\begin{enumerate}
\item $A$-模复形 $K^\bullet$；
\item 函子变换 $T_x(-) \to H^1(K^\bullet \otimes_A^\mathbf{L} -)$；
\item 对每个核为 $I = \Ker(A' \to A)$ 的形变情形 $(x, A' \to A)$，
给定元素
$o_x(A') \in H^2(K^\bullet \otimes_A^\mathbf{L} I)$
\end{enumerate}
，并满足下列（最小）条件：
\begin{enumerate}
\item[(i)] 变换
$T_x(-) \to H^1(K^\bullet \otimes_A^\mathbf{L} -)$
是同构；
\item[(ii)] 给定形变情形的态射
$(x, A'' \to A) \to (x, A' \to A)$，元素 $o_x(A')$ 经映射
$H^2(K^\bullet \otimes_A^\mathbf{L} I) \to
H^2(K^\bullet \otimes_A^\mathbf{L} I')$
映到元素 $o_x(A'')$，其中 $I' = \Ker(A'' \to A)$；
\item[(iii)] $x$ 可提升为 $\Spec(A')$ 上的对象，当且仅当
$o_x(A') = 0$。
\end{enumerate}
也可以把无穷小自同构纳入其中，但为了得到尽可能强的结果，我们不这样做。
\end{situation}

\noindent
在情形 \ref{artin-situation-dual} 中，
$K^\bullet \otimes_A^\mathbf{L} \NL_{A/\Lambda}$ 将起重要作用。假设给定元素
$\xi \in H^1(K^\bullet \otimes_A^\mathbf{L} \NL_{A/\Lambda})$。于是：
(1) 对任意核为平方零理想 $I$ 的 $\Lambda$-代数满射 $A' \to A$，典范映射
$\NL_{A/\Lambda} \to \NL_{A/A'} = I[1]$ 把 $\xi$ 送到元素
$\xi_{A'} \in H^2(K^\bullet \otimes_A^\mathbf{L} I)$；(2) 映射
$\NL_{A/\Lambda} \to \Omega_{A/\Lambda}$ 把 $\xi$ 送到
$H^1(K^\bullet \otimes_A^\mathbf{L} \Omega_{A/\Lambda})$ 的元素 $\xi_{can}$。

\begin{lemma}
\label{artin-lemma-dual-obstruction}
% P11-E0285: the frozen source leaves "In Situation ..." as a sentence fragment.
在情形 \ref{artin-situation-dual} 中，进一步假设
\begin{enumerate}
\item[(iv)] 给定评注 \ref{artin-remark-short-exact-sequence-thickenings} 中那样的
形变情形短正合列以及提升 $x'_2 \in \text{Lift}(x, A_2')$，则
$o_x(A_3') \in H^2(K^\bullet \otimes_A^\mathbf{L} I_3)$
等于 $\partial\theta$，其中
$\theta \in H^1(K^\bullet \otimes_A^\mathbf{L} I_1)$
是经 $A_1' = A[I_1]$ 及给定映射
$T_x(-) \to H^1(K^\bullet \otimes_A^\mathbf{L} -)$ 与
$x'_2|_{\Spec(A_1')}$ 对应的元素。
\end{enumerate}
在这种情形下，存在元素
$\xi \in H^1(K^\bullet \otimes_A^\mathbf{L} \NL_{A/\Lambda})$
，使得
\begin{enumerate}
\item 对每个形变情形 $(x, A' \to A)$，都有
$\xi_{A'} = o_x(A')$；
\item 经给定变换 $T_x(-) \to H^1(K^\bullet \otimes_A^\mathbf{L} -)$，
$\xi_{can}$ 与评注 \ref{artin-remark-canonical-element} 的典范元素相匹配。
\end{enumerate}
\end{lemma}

\begin{proof}
% P11-E0282: the frozen source uses "a" before alpha instead of "an".
选取核为 $J$ 的 $\alpha : \Lambda[x_1, \ldots, x_n] \to A$。记
$P = \Lambda[x_1, \ldots, x_n]$。在证明的余下部分，使用
$$
\NL(\alpha) = (J/J^2 \longrightarrow \bigoplus A \text{d}x_i)
$$
；《代数》引理 \ref{algebra-lemma-NL-homotopy} 和《代数进阶》引理
\ref{more-algebra-lemma-derived-tor-homotopy} 保证这样做是允许的。考虑元素
$o_x(P/J^2) \in H^2(K^\bullet \otimes_A^\mathbf{L} J/J^2)$ 以及商
$$
C = (P/J^2 \times \bigoplus A \text{d}x_i)/(J/J^2)
$$
，其中 $J/J^2$ 对角嵌入。注意，$C \to A$ 是核为
$\bigoplus A\text{d}x_i$ 的满射。此外，$C \to A$ 有截面 $A \to C$，
它把 $f \in P$ 的类映到推出中 $(f, \text{d}f)$ 的类。
% P11-E0283: the frozen source omits "by" in "denote x_C the pullback".
为后用，把 $x$ 沿相应态射 $\Spec(C) \to \Spec(A)$ 的拉回记为 $x_C$。
因此 $o_x(C) = 0$。由此，$o_x(P/J^2)$ 在
$H^2(K^\bullet \otimes_A^\mathbf{L} \bigoplus A\text{d}x_i)$ 中映为零。
于是存在元素 $\xi \in H^1(K^\bullet \otimes_A^\mathbf{L} \NL(\alpha))$，
它映到 $o_x(P/J^2)$。

\medskip\noindent
注意，对任意形变情形 $(x, A' \to A)$，存在与到 $A$ 的增广相容的
$\Lambda$-代数同态 $P/J^2 \to A'$。因此，按构造及情形
\ref{artin-situation-dual} 的性质 (ii)，元素 $\xi$ 满足引理的第一个性质。

\medskip\noindent
注意，$\xi$ 的选取只在相差
$H^1(K^\bullet \otimes_A^\mathbf{L} \bigoplus A\text{d}x_i)$ 的元素的意义下
良定义。我们将证明，用上述群的一个元素修正 $\xi$ 后，可使引理的第二个断言
成立。令 $C' \subset C$ 为 $P/J^2$ 在 $\Lambda$-代数同态
$P/J^2 \to C$（第一因子的包含）下的像。注意
$\Ker(C' \to A) = \Im(J/J^2 \to \bigoplus A\text{d}x_i)$.
令 $\overline{C} = A[\Omega_{A/\Lambda}]$。映射
$P/J^2 \times \bigoplus A \text{d}x_i \to \overline{C}$,
$(f, \sum f_i \text{d}x_i) \mapsto (f \bmod J, \sum f_i \text{d}x_i)$
分解经过满射 $C \to \overline{C}$。于是
$$
(x, \overline{C} \to A) \to (x, C \to A) \to (x, C' \to A)
$$
是形变情形的短正合列。与之相伴的分裂
$\overline{C} = A[\Omega_{A/\Lambda}]$（来自评注
\ref{artin-remark-short-exact-sequence-thickenings}）等于上面给定的分裂。此外，
截面 $A \to C$ 与映射 $C \to \overline{C}$ 的复合，就是评注
\ref{artin-remark-canonical-element} 中的映射
$(1, \text{d}) : A \to A[\Omega_{A/\Lambda}]$。因此，$x_C$ 限制为
$T_x(\Omega_{A/\Lambda}) = \text{Lift}(x, A[\Omega_{A/\Lambda}])$ 的典范元素
$x_{can}$。由条件 (iv)，$o_x(P/J^2)$ 在
$$
H^1(K^\bullet \otimes_A^\mathbf{L} \Im(J/J^2 \to \bigoplus A\text{d}x_i))
$$
中映到 $\partial x_{can}$。按构造，$\xi$ 映到 $o_x(P/J^2)$。因此，
$x_{can}$ 和 $\xi_{can}$ 在上面显示的群中映到同一元素；由上同调长正合列，
这意味着二者相差 $H^1(K^\bullet \otimes_A^\mathbf{L}
\bigoplus A\text{d}x_i)$ 的一个元素，正如所需。
\end{proof}

\begin{lemma}
\label{artin-lemma-dual-openness}
在情形 \ref{artin-situation-dual} 中，假设引理 \ref{artin-lemma-dual-obstruction} 的
(iv) 成立，并且 $K^\bullet$ 是 $D(A)$ 的完美对象。在这种情形下，若 $x$
在闭点 $u_0 \in U$ 处通用，则存在开邻域 $u_0 \in U' \subset U$，使得
$x$ 在 $U'$ 的每个有限型点处通用。
\end{lemma}

\begin{proof}
可以假设 $K^\bullet$ 是由有限射影 $A$-模组成的有限复形。于是，与
$K^\bullet$ 作导出张量积就等于直接与 $K^\bullet$ 作张量积。令
$E^\bullet$ 为 $K^\bullet$ 的对偶完美复形；参见《代数进阶》引理
\ref{more-algebra-lemma-dual-perfect-complex}。（因此
$E^n = \Hom_A(K^{-n}, A)$，其微分是 $K^\bullet$ 的微分的转置。）令
$E \in D^{-}(A)$ 表示由复形 $E^\bullet[-1]$ 表示的对象。令
$\xi \in H^1(\text{Tot}(K^\bullet \otimes_A \NL_{A/\Lambda}))$ 为引理
\ref{artin-lemma-dual-obstruction} 中构造的元素。
% P11-E0284: the frozen source omits "by" in the naming construction for the corresponding map xi.
把相应映射记为 $\xi : E = E^\bullet[-1] \to \NL_{A/\Lambda}$（同上）。
我们断言，二元组 $(E, \xi)$ 满足引理 \ref{artin-lemma-openness} 的全部假设，
从而证明完毕。

\medskip\noindent
确实，引理 \ref{artin-lemma-openness} 的假设 (i) 来自引理
\ref{artin-lemma-dual-obstruction} 的结论 (1) 以及同上所述的事实
$H^2(K^\bullet \otimes_A^\mathbf{L} -) = \Ext^1(E, -)$。引理
\ref{artin-lemma-openness} 的假设 (ii) 来自引理 \ref{artin-lemma-dual-obstruction} 的结论
(2) 以及同上所述的事实
$H^1(K^\bullet \otimes_A^\mathbf{L} -) = \Ext^0(E, -)$。引理
\ref{artin-lemma-openness} 的假设 (iii) 是显然的。
\end{proof}







\section{Noether 概形上的保极限函子}
\label{artin-section-noetherian}

\noindent
有时，只考虑定义在（局部）Noether 概形全子范畴上的函子或叠会很方便。
本节讨论代数空间的情形。

\medskip\noindent
设 $S$ 是局部 Noether 概形。为了正确对齐各个范畴，我们稍微严格一些；
忽略集合论问题的读者，可以在下文中把所选的概形集合直接替换为所有概形的类。
如《拓扑》评注 \ref{topologies-remark-choice-sites}，选取包含 $S$ 的概形范畴
$\Sch_\alpha$，使得得到大站点 $(\Sch/S)_{Zar}$、
$(\Sch/S)_\etale$, $(\Sch/S)_{smooth}$, $(\Sch/S)_{syntomic}$,
和 $(\Sch/S)_{fppf}$，它们具有相同的底范畴 $\Sch_\alpha/S$。
% P11-E0286: the frozen source omits "by" in the naming construction for Noetherian_alpha.
把
$$
\textit{Noetherian}_\alpha \subset \Sch_\alpha
$$
记为由局部 Noether 概形组成的全子范畴。这确定了全子范畴
$$
\textit{Noetherian}_\alpha/S \subset \Sch_\alpha/S
$$
对 $\tau \in \{Zariski, \etale, smooth, syntomic, fppf\}$，有
\begin{enumerate}
\item 若 $f : X \to Y$ 是 $\Sch_\alpha/S$ 的态射，
$Y$ 属于 $\textit{Noetherian}_\alpha/S$，且 $f$ 局部有限型，则
$X$ 属于 $\textit{Noetherian}_\alpha/S$；
\item 若 $f : X \to Y$ 和 $g : Z \to Y$ 是
$\textit{Noetherian}_\alpha/S$ 的态射，且 $f$ 局部有限型，则
$\textit{Noetherian}_\alpha/S$ 中的纤维积 $X \times_Y Z$ 存在，
并与 $\Sch_\alpha/S$ 中的纤维积一致；
\item 若 $\{X_i \to X\}_{i \in I}$ 是 $(\Sch/S)_\tau$ 的覆盖，且
$X$ 属于 $\textit{Noetherian}_\alpha/S$，则对象 $X_i$ 都属于
$\textit{Noetherian}_\alpha/S$；
\item 在范畴 $\textit{Noetherian}_\alpha/S$ 上赋予
$(\Sch/S)_\tau$ 中对象均属于 $\textit{Noetherian}_\alpha/S$ 的覆盖族，
便得到一个站点，记为 $(\textit{Noetherian}/S)_\tau$；
\item 包含函子
$(\textit{Noetherian}/S)_\tau \to (\Sch/S)_\tau$
是全忠实、连续且余连续的。
\end{enumerate}
由《站点》引理 \ref{sites-lemma-cocontinuous-morphism-topoi} 和
\ref{sites-lemma-when-shriek}，得到拓扑斯态射
$$
g_\tau : \Sh((\textit{Noetherian}/S)_\tau) \longrightarrow \Sh((\Sch/S)_\tau)
$$
；其拉回函子就是沿包含函子
$(\textit{Noetherian}/S)_\tau \to (\Sch/S)_\tau$ 限制层。

\begin{remark}[警告]
\label{artin-remark-no-fibre-products}
站点 $(\textit{Noetherian}/S)_\tau$ 不具有纤维积。因此，处理层时必须谨慎。
例如，连续包含函子
$(\textit{Noetherian}/S)_\tau \to (\Sch/S)_\tau$
并不定义站点态射。当 $\tau = fppf$ 时的一个例子，参见《例》一节
\ref{examples-section-sheaves-locally-Noetherian}。
\end{remark}

\noindent
设 $F : (\textit{Noetherian}/S)_\tau^{opp} \to \textit{Sets}$ 为函子。
若对 $(\textit{Noetherian}/S)_\tau$ 中仿射概形的任意有向极限
$X = \lim X_i$，都有 $F(X) = \colim F(X_i)$，则称 $F$ {\it 保极限}。

\begin{lemma}
\label{artin-lemma-canonical-extension}
设 $\tau \in \{Zariski, \etale, smooth, syntomic, fppf\}$。沿包含函子
$(\textit{Noetherian}/S)_\tau \to (\Sch/S)_\tau$ 作限制，定义下列两范畴之间
的等价：
\begin{enumerate}
\item $(\Sch/S)_\tau$ 上的保极限层所成的范畴；
\item $(\textit{Noetherian}/S)_\tau$ 上的保极限层所成的范畴。
\end{enumerate}
\end{lemma}

\begin{proof}
设 $F : (\textit{Noetherian}/S)_\tau^{opp} \to \textit{Sets}$ 是既保极限
又为层的函子。由《拓扑》引理 \ref{topologies-lemma-extend} 和
\ref{topologies-lemma-extend-sheaf-general}，存在唯一的函子
$F' : (\Sch/S)_\tau^{opp} \to \textit{Sets}$，它保极限、是层，且限制为
$F$。事实上，《拓扑》引理 \ref{topologies-lemma-extend} 中 $F'$ 的构造关于
$F$ 具有函子性，而且该构造是限制函子的拟逆。略去一些细节。
\end{proof}

\begin{lemma}
\label{artin-lemma-representable-limit-preserving}
设 $X$ 是 $(\textit{Noetherian}/S)_\tau$ 的对象。若点函子
$h_X : (\textit{Noetherian}/S)_\tau^{opp} \to \textit{Sets}$ 保极限，
则 $X$ 在 $S$ 上局部有限表示。
\end{lemma}

\begin{proof}
设 $V \subset X$ 是映入仿射开集 $U \subset S$ 的仿射开子概形。可以把
$V = \lim V_i$ 写成 $U$ 上有限表示仿射概形 $V_i$ 的有向极限；参见
《代数》引理 \ref{algebra-lemma-ring-colimit-fp}。按假设，对某个 $i$，箭头
$V \to X$ 分解为 $V \to V_i \to X$。增大 $i$ 后，可以假设
$V_i \to X$ 分解经过 $V$，因为由《极限》引理
\ref{limits-lemma-descend-opens}，$V \subset X$ 在 $V_i$ 中的逆像最终等于
$V_i$。于是，在 $U$ 上概形的范畴中，对某个 $i$，恒等态射 $V \to V$
分解经过 $V_i$。因此 $V \to U$ 有限表示；相应的代数事实是：若 $B$ 是
$A$-代数，且 $\text{id} : B \to B$ 分解经过某个有限表示 $A$-代数，
则 $B$ 在 $A$ 上有限表示（一个不错的练习）。所以 $X$ 在 $S$ 上局部有限表示。
\end{proof}

\noindent
下一个引理有一个适用于可由代数空间表示的变换的版本。

\begin{lemma}
\label{artin-lemma-representable}
设 $\tau \in \{Zariski, \etale, smooth, syntomic, fppf\}$。设
$F', G' : (\Sch/S)_\tau^{opp} \to \textit{Sets}$ 是保极限层，
$a' : F' \to G'$ 是函子变换。
% P11-E0287: the frozen source omits "by" in the naming construction for a : F -> G.
把 $a' : F' \to G'$ 到 $(\textit{Noetherian}/S)_\tau$ 的限制记为
$a : F \to G$。下列条件等价：
\begin{enumerate}
\item $a'$ 可表示（作为函子变换；参见《范畴》定义
\ref{categories-definition-representable-morphism}）；
\item 对 $(\textit{Noetherian}/S)_\tau$ 的每个对象 $V$ 和每个映射
$V \to G$，纤维积
$F \times_G V : (\textit{Noetherian}/S)_\tau^{opp} \to \textit{Sets}$
是可表示函子；
\item 与 (2) 相同，但只要求 $V$ 是映入 $S$ 的某个仿射开集的 $S$ 上有限型
仿射概形。
\end{enumerate}
\end{lemma}

\begin{proof}
假设 (1)。由《空间的极限》引理
\ref{spaces-limits-lemma-locally-finite-presentation-permanence}，变换 $a'$ 保极限
\footnote{即使 $\tau \not = fppf$，这也有意义，因为 $(\Sch/S)_\tau$ 的
底范畴等于 $(\Sch/S)_{fppf}$ 的底范畴，而该陈述与拓扑无关。}。取 (2) 中的
$\xi : V \to G$。
% P11-E0288: the frozen source uses the incomplete naming construction "Denote V' = V but viewed ...".
记 $V' = V$，但把它看作 $(\Sch/S)_\tau$ 的对象。由于 $G$ 是 $G'$ 到
$(\textit{Noetherian}/S)_\tau$ 的限制，$\xi \in G(V)$ 对应于
$\xi' \in G'(V')$。按假设，$V' \times_{\xi', G'} F'$ 可由概形 $U'$ 表示。
与投影 $V' \times_{\xi', G'} F' \to V'$ 对应的概形态射 $U' \to V'$，由
《空间的极限》引理
\ref{spaces-limits-lemma-base-change-locally-finite-presentation} 和《极限》命题
\ref{limits-proposition-characterize-locally-finite-presentation}，是局部有限表示的。
因此 $U'$ 是局部 Noether 概形，从而 $U'$ 同构于
$(\textit{Noetherian}/S)_\tau$ 的某个对象 $U$。于是 $U$ 表示
$F \times_G V$，正如所需。

\medskip\noindent
(2) $\Rightarrow$ (3) 是立即的。假设 (3)。我们来证明 (1)。设 $T$ 是
$(\Sch/S)_\tau$ 的对象，并设 $T \to G'$ 为态射。需要证明函子
$F' \times_{G'} T$ 可由 $T$ 上的某个概形 $X$ 表示。令 $\mathcal{B}$ 为
$T$ 中映入 $S$ 的某个仿射开集的仿射开集所成的集合。这是 $T$ 的拓扑的一组基。
下面证明：对 $W \in \mathcal{B}$，纤维积 $F' \times_{G'} W$ 可由 $W$ 上的
某个概形 $X_W$ 表示。若 $\mathcal{B}$ 中 $W_1 \subset W_2$，则得到同构
$X_{W_1} \to X_{W_2} \times_{W_2} W_1$，因为 $X_{W_1}$ 和
$X_{W_2} \times_{W_2} W_1$ 都表示函子 $F' \times_{G'} W_1$。这些同构是典范的，
并满足《构造》引理 \ref{constructions-lemma-relative-glueing} 中提到的余圈条件。
因此，可以把概形 $X_W$ 黏合成 $T$ 上的概形 $X$。
% P11-E0289: the frozen source has subject-verb disagreement in "Compatibility ... provide".
黏合映射与映射 $X_W \to F'$ 的相容性给出映射 $X \to F'$。所得映射
$X \to F' \times_{G'} T$ 是同构，因为这可在 $T$ 上局部检验（该箭头的源与
目标都是 Zariski 拓扑的层）。

\medskip\noindent
设 $W$ 是映入仿射开集 $U \subset S$ 的仿射概形，并设 $W \to G'$ 为映射。
仍假设 (3)，需要证明 $F' \times_{G'} W$ 可由概形表示。可以把
$W = \lim V'_i$ 写成 $U$ 上有限表示仿射概形 $V'_i$ 的有向极限；参见
《代数》引理 \ref{algebra-lemma-ring-colimit-fp}。
% P11-E0290: the frozen source uses "an Noetherian scheme" instead of "a Noetherian scheme".
由于 $V'_i$ 在某个 Noether 概形上有限型，$V'_i$ 是 Noether 概形。
% P11-E0291: the frozen source uses the incomplete naming construction "Denote V_i = V'_i but viewed ...".
记 $V_i = V'_i$，但把它看作 $(\textit{Noetherian}/S)_\tau$ 的对象。
% P11-E0292: the frozen source omits "we" in "As G' is limit preserving can choose".
由于 $G'$ 保极限，可以选取某个 $i$ 和映射 $V'_i \to G'$，使得
$W \to G'$ 是复合 $W \to V'_i \to G'$。由于 $G$ 是 $G'$ 到
$(\textit{Noetherian}/S)_\tau$ 的限制，态射 $V'_i \to G'$ 就等同于态射
$V_i \to G$（见上文）。由假设 (3)，函子 $F \times_G V_i$ 可由
$(\textit{Noetherian}/S)_\tau$ 的某个对象 $X_i$ 表示。函子
$F \times_G V_i$ 保极限，因为它是 $F' \times_{G'} V'_i$ 的限制，而后一个
函子由《空间的极限》引理
\ref{spaces-limits-lemma-fibre-product-locally-finite-presentation}、$F'$ 和 $G'$
保极限的假设，以及《极限》评注 \ref{limits-remark-limit-preserving}（它说明
$V'_i$ 的点函子保极限）而保极限。由引理
\ref{artin-lemma-representable-limit-preserving}，$X_i$ 在 $S$ 上局部有限表示。
% P11-E0293: the frozen source uses the incomplete naming construction "Denote X'_i = X_i but viewed ...".
记 $X'_i = X_i$，但把它看作 $(\Sch/S)_\tau$ 的对象。于是
$F' \times_{G'} V'_i$ 和点函子 $h_{X'_i}$ 都是
$h_{X_i} : (\textit{Noetherian}/S)_\tau^{opp} \to \textit{Sets}$ 到
$(\Sch/S)_\tau$ 上保极限层的扩张。由引理 \ref{artin-lemma-canonical-extension}
的范畴等价，推出 $X'_i$ 表示 $F' \times_{G'} V'_i$。最后，
$$
F' \times_{G'} W = F' \times_{G'} V'_i \times_{V'_i} W =
X'_i \times_{V'_i} W
$$
可表示，正如所需。
\end{proof}






\section{Noether 情形中的代数空间}
\label{artin-section-algebraic-spaces-noetherian}

\noindent
设 $S$ 是局部 Noether 概形。以
$(\textit{Noetherian}/S)_\etale \subset (\Sch/S)_\etale$ 表示一节
\ref{artin-section-noetherian} 中研究的站点。设
$F : (\textit{Noetherian}/S)_\etale^{opp} \to \textit{Sets}$ 为函子；
换言之，$F$ 是 $(\textit{Noetherian}/S)_\etale$ 上的预层。在此设置下，
一节 \ref{artin-section-axioms-functors} 的所有公理 [-1]、[0]、[1]、[2]、[3]、
[4]、[5] 都有意义。下面逐一回顾，以明确我们的确切含义。

\medskip\noindent
公理 [-1]。这是一个集合论条件；不关心集合论问题的读者可以忽略它。
它对 $F$ 有意义，因为它涉及在 $S$ 上有限型域的谱处计算 $F$，而这些谱都是
$(\textit{Noetherian}/S)_\etale$ 的对象。

\medskip\noindent
% P11-E0295: the frozen source says the sheaf is on the opposite category rather than on the site.
公理 [0]。这是说 $F$ 是 $(\textit{Noetherian}/S)_\etale^{opp}$ 上的层，
即满足 \'etale 覆盖的层条件。

\medskip\noindent
公理 [1]。这是说 $F$ 按一节 \ref{artin-section-noetherian} 的定义保极限：
对 $(\textit{Noetherian}/S)_\etale$ 中仿射概形的任意有向极限
$X = \lim X_i$，都有 $F(X) = \colim F(X_i)$。

\medskip\noindent
公理 [2]。这是说 $F$ 满足 Rim--Schlessinger 条件 (RS)。考察定义
\ref{artin-definition-RS} 中条件 (RS) 的定义及一节
\ref{artin-section-axioms-functors} 中的讨论，可知其含义如下：给定 $S$ 上有限型
概形的任意推出 $Y' = Y \amalg_X X'$，其中 $Y, X, X'$ 是 Artin 局部环的谱，
则
$$
F(Y \amalg_X X') \to F(Y) \times_{F(X)} F(X')
$$
为双射。此条件有意义，因为概形 $X$、$X'$、$Y$、$Y'$ 在 $S$ 上有限型，
故都属于 $(\text{Noetherian}/S)_\etale$。

\medskip\noindent
公理 [3]。这是说每个切空间 $TF_{k, x_0}$ 都是有限维的。此条件有意义，
因为切空间 $TF_{k, x_0}$ 由 $F$ 在 $\Spec(k)$ 和
$\Spec(k[\epsilon])$ 处的取值构造，其中 $k$ 是 $S$ 上有限型域；因此，
这些取值都在范畴 $(\textit{Noetherian}/S)_\etale$ 的对象处计算。

\medskip\noindent
% P11-E0296: the frozen source has both a missing article before axiom and a spurious "the" before every.
公理 [4]。这是说每个形式对象都有效。考察各节
\ref{artin-section-formal-objects} 和 \ref{artin-section-axioms-functors} 中的讨论，
可知这里只涉及在 Noether 概形处计算函子，因此此条件对 $F$ 有意义。

\medskip\noindent
公理 [5]。这是说 $F$ 满足通用性的开性。回忆，其含义如下：给定 $S$ 上局部
有限型概形 $U$、$x \in F(U)$，以及有限型点 $u_0 \in U$，且 $x$ 在
$u_0$ 处通用，则存在开邻域 $u_0 \in U' \subset U$，使得 $x$ 在 $U'$
的每个有限型点处通用。与前面一样，验证此条件只涉及在 Noether 概形处计算函子。

\begin{proposition}
\label{artin-proposition-spaces-diagonal-representable-noetherian}
设 $S$ 是局部 Noether 概形，
$F : (\textit{Noetherian}/S)_\etale^{opp} \to \textit{Sets}$ 是函子。
假设
\begin{enumerate}
\item $\Delta : F \to F \times F$ 可表示（作为函子变换；参见《范畴》定义
\ref{categories-definition-representable-morphism}）；
\item $F$ 满足公理 [-1]、[0]、[1]、[2]、[3]、[4]、[5]（见上文）；
\item 对 $S$ 的所有有限型点 $s$，$\mathcal{O}_{S, s}$ 都是 G-环。
\end{enumerate}
则存在唯一的代数空间
$F' : (\Sch/S)_{fppf}^{opp} \to \textit{Sets}$
，其到 $(\textit{Noetherian}/S)_\etale$ 的限制为 $F$（说明见证明）。
\end{proposition}

\begin{proof}
回忆，站点 $(\Sch/S)_{fppf}$ 和 $(\Sch/S)_\etale$ 具有相同的底范畴；
参见一节 \ref{artin-section-noetherian} 中的讨论。类似地，站点
$(\textit{Noetherian}/S)_\etale$ 和
$(\textit{Noetherian}/S)_{fppf}$ 具有相同的底范畴。
% P11-E0297: the frozen source uses plural "underlying categories" for one shared underlying category.
由公理 [0] 和 [1]，函子 $F$ 是层且保极限。设
$F' : (\Sch/S)_\etale^{opp} \to \textit{Sets}$ 为 $F$ 的唯一扩张，
它是层（关于 \'etale 拓扑）且保极限；参见引理
\ref{artin-lemma-canonical-extension}。于是 $F'$ 满足一节
\ref{artin-section-axioms-functors} 中给出的公理 [0] 和 [1]。由引理
\ref{artin-lemma-representable}，$\Delta' : F' \to F' \times F'$ 可表示（由概形表示）。
另一方面，$F'$ 显然满足一节 \ref{artin-section-axioms-functors} 的公理
[-1]、[2]、[3]、[4]、[5]，因为这些公理都只涉及在
$(\textit{Noetherian}/S)_\etale$ 的对象处计算 $F'$，而我们已假设 $F$
满足相应条件。因此，由命题 \ref{artin-proposition-spaces-diagonal-representable}，
$F'$ 是代数空间。
\end{proof}












\section{Artin 收缩定理}
\label{artin-section-contractions}

\noindent
本节将自由使用形式代数空间的语言；参见《形式空间》一节
\ref{formal-spaces-section-introduction}。Artin 收缩定理是 Artin 论文
\cite{ArtinII} 的两个主要定理之一；第一个定理是他的膨胀定理，我们已在
《形式空间的代数化》一节 \ref{restricted-section-dilatations} 中陈述并证明。

\begin{situation}
\label{artin-situation-contractions}
设 $S$ 是局部 Noether 概形，$X'$ 是 $S$ 上局部有限型的代数空间，
$T' \subset |X'|$ 是闭子集。设 $U' \subset X'$ 为满足
$|U'| = |X'| \setminus T'$ 的开子空间。设 $W$ 是 $S$ 上局部 Noether 的
形式代数空间，且 $W_{red}$ 在 $S$ 上局部有限型。最后，设
$$
g : X'_{/T'} \longrightarrow W
$$
为形式修正；参见《形式空间的代数化》定义
\ref{restricted-definition-formal-modification}。回忆，$X'_{/T'}$
表示 $X'$ 沿 $T'$ 的形式完备化；参见《形式空间》一节
\ref{formal-spaces-section-completion}。
\end{situation}

\noindent
在上述情形中，我们的目标是证明：存在 $S$ 上代数空间的真态射
$f : X' \to X$、闭子集 $T \subset |X|$，以及形式代数空间的同构
$a : X_{/T} \to W$，使得
\begin{enumerate}
\item $T'$ 是 $T$ 在 $|f| : |X'| \to |X|$ 下的逆像；
\item $f : X' \to X$ 将 $U'$ 同构地映到 $X$ 的某个开子空间 $U$；以及
\item $g = a \circ f_{/T}$，其中
$f_{/T} : X'_{/T'} \to X_{/T}$ 是诱导态射。
\end{enumerate}
称 $(f : X' \to X, T, a)$ 为一个{\it 解}。

\medskip\noindent
我们将依照 Artin 的策略：在 $S$ 上局部 Noether 概形的范畴上构造一个函子
$F$，利用命题 \ref{artin-proposition-spaces-diagonal-representable-noetherian}
证明 $F$ 是代数空间，并证明取 $X = F$ 即可。

\begin{remark}
\label{artin-remark-G-rings}
特别地，我们无法证明所求结果对情形 \ref{artin-situation-contractions} 的每个实例
都成立，因为我们将需要假设 $S$ 的局部环都是 G-环。若读者能在一般情形下证明
该结果，或有反例，请通过
\href{mailto:stacks.project@gmail.com}{stacks.project@gmail.com}
告知我们。
\end{remark}

\noindent
在情形 \ref{artin-situation-contractions} 中，设 $V$ 是 $S$ 上局部 Noether 概形。
我们的函子 $F$ 在 $V$ 上的取值由所有如下三元组组成：
$$
(Z, u' : V \setminus Z \to U', \hat x : V_{/Z} \to W)
$$
它们满足下列条件：
\begin{enumerate}
\item $Z \subset V$ 是闭子集；
\item $u' : V \setminus Z \to U'$ 是 $S$ 上的态射；
\item $\hat x : V_{/Z} \to W$ 是 $S$ 上形式代数空间的 adic 态射；
\item $u'$ 与 $\hat x$ 相容（见下文）。
\end{enumerate}
相容性条件如下：拉回形式修正 $g$，得到形式修正
$$
X'_{/T'} \times_{g, W, \hat x} V_{/Z} \longrightarrow V_{/Z}
$$
；参见《形式空间的代数化》引理
\ref{restricted-lemma-base-change-formal-modification}。由膨胀的主定理
（《形式空间的代数化》定理 \ref{restricted-theorem-dilatations}），存在唯一的
代数空间真态射 $V' \to V$，它在 $V \setminus Z$ 上为同构，并且
$V'_{/Z} \to V_{/Z}$ 与上面显示的箭头同构。换言之，对某个态射
$\hat x' : V'_{/Z} \to X'_{/T'}$，有形式代数空间的 Cartesian 图
$$
\xymatrix{
V'_{/Z} \ar[r] \ar[d]_{\hat x'} & V_{/Z} \ar[d]^{\hat x} \\
X'_{/T'} \ar[r]^g & W
}
$$
以后不再特别说明地把 $V \setminus Z$ 视为 $V'$ 的开子空间。相容性条件是：
应存在态射 $x' : V' \to X'$，它在 $V \setminus Z \subset V'$ 和
$V'_{/Z}$ 上分别限制为 $u'$ 和 $\hat x$。
% P11-E0298: the frozen prose says the restriction on V'_{/Z} is \hat x, while the displayed diagram and preceding definition use \hat x'.
换言之，使图
$$
\xymatrix{
V \setminus Z \ar[r] \ar[d]_{u'} &
V' \ar[d]^{x'} &
V'_{/Z} \ar[l] \ar[d]^{\hat x'} \ar[r] &
V_{/Z} \ar[d]^{\hat x} \\
U' \ar[r] &
X' &
X'_{/T'} \ar[r]^g \ar[l] &
W
}
$$
交换。注意，由《形式空间的代数化》引理
\ref{restricted-lemma-faithful}，若态射 $x'$ 存在，则它是唯一的。
我们将用
“{\it $V' \to V$、$\hat x'$ 和 $x'$ 见证 $u'$ 与
$\hat x$ 之间的相容性}”表示这种情形。

\begin{remark}
\label{artin-remark-how-to-think-compatibility}
在情形 \ref{artin-situation-contractions} 中，设 $V$ 是 $S$ 上局部 Noether 概形，
$(Z, u', \hat x)$ 是满足上述 (1)、(2)、(3) 的三元组。我们想说明一种理解
相容性条件 (4) 的方式。由于我们将使用 $S$ 上解析空间的虚构范畴
$\textit{An}_S$ 以及虚构的解析化函子，所述方式在数学上并不精确。
% P11-E0299: the frozen source omits "to" in "we are going use".
$$
\left\{
\begin{matrix}
\text{locally Noetherian formal} \\
\text{algebraic spaces over }S
\end{matrix}
\right\}
\longrightarrow
\textit{An}_S,
\quad\quad
Y \longmapsto Y^{an}
$$
例如，若在 $S = \Spec(k)$ 上有 $Y = \text{Spf}(k[[t]])$，则应把 $Y^{an}$
理解为开单位圆盘。若 $Y = \Spec(k)$，则 $Y^{an}$ 是单点。范畴
$\textit{An}_S$ 应有开浸入和闭浸入，并且我们应能取闭集的开补。给定 $Y$，
态射 $Y_{red} \to Y$ 应诱导闭浸入 $Y_{red}^{an} \to Y^{an}$。
% P11-E0300: the frozen source has the agreement error "should induces".
令 $Y^{rig} = Y^{an} \setminus Y_{red}^{an}$ 为其开补。若 $Y$ 是代数空间且
$Z \subset Y$ 为闭集，则态射 $Y_{/Z} \to Y$ 应诱导开浸入
$Y_{/Z}^{an} \to Y^{an}$，继而诱导开浸入
$$
can : (Y_{/Z})^{rig} \longrightarrow (Y \setminus Z)^{an}
$$
此外，给定局部 Noether 形式代数空间的形式修正 $g : Y' \to Y$，诱导态射
$g^{rig} : (Y')^{rig} \to Y^{rig}$ 应为同构。给定 $\text{An}_S$ 和解析化
函子，可以考虑如下要求：
$$
\xymatrix{
(V_{/Z})^{rig} \ar[rr]_{can} \ar[d]_{(g^{rig})^{-1} \circ \hat x^{an}} & &
(V \setminus Z)^{an} \ar[d]^{(u')^{an}} \\
(X'_{/T'})^{rig} \ar[rr]^{can} & & (X' \setminus T')^{an}
}
$$
交换。
% P11-E0301: the frozen source writes (X'_{T'})^{rig}; the formal completion used throughout is X'_{/T'}.
此要求有意义，因为 $g^{rig} : (X'_{T'})^{rig} \to W^{rig}$ 是同构，且
$U' = X' \setminus T'$。最后，在解析化函子满足某些忠实性假设时，此要求将
等价于上面表述的相容性条件。我们希望这能促使读者把 $u'$ 与 $\hat x$ 的
相容性理解为某些映射相等这一要求，而不是要求某个交换图存在。
\end{remark}

\begin{lemma}
\label{artin-lemma-functor}
在情形 \ref{artin-situation-contractions} 中，规则 $F$ 把 $S$ 上局部 Noether 概形
$V$ 映到满足相容性条件的三元组 $(Z, u', \hat x)$ 所成的集合，并把
$S$ 上局部 Noether 概形的态射 $\varphi : V_2 \to V_1$ 映到映射
% P11-E0302: the frozen source has the duplicated article "a a morphism".
$$
F(\varphi) : F(V_1) \longrightarrow F(V_2)
$$
；后者把 $F(V_1)$ 的元素 $(Z_1, u'_1, \hat x_1)$ 映到 $F(V_2)$ 中由下列
规则给出的 $(Z_2, u'_2, \hat x_2)$：
\begin{enumerate}
\item $Z_2 \subset V_2$ 是 $Z_1$ 在 $\varphi$ 下的逆像；
\item $u'_2$ 是 $u'_1$ 与
$\varphi|_{V_2 \setminus Z_2} : V_2 \setminus Z_2 \to V_1 \setminus Z_1$
的复合；
\item $\hat x_2$ 是 $\hat x_1$ 与
$\varphi_{/Z_2} : V_{2, /Z_2} \to V_{1, /Z_1}$
的复合。
\end{enumerate}
则该规则是反变函子。
\end{lemma}

\begin{proof}
为看出 $u'_2$ 与 $\hat x_2$ 之间的相容性，设
$V'_1 \to V_1$、$\hat x'_1$ 和 $x'_1$ 见证 $u'_1$ 与 $\hat x_1$ 之间的
相容性。令 $V'_2 = V_2 \times_{V_1} V'_1$；令 $\hat x'_2$ 为
$\hat x'_1$ 与 $V'_{2, /Z_2} \to V'_{1, /Z_1}$ 的复合；令 $x'_2$
为 $x'_1$ 与 $V'_2 \to V'_1$ 的复合。于是
$V'_2 \to V_2$、$\hat x'_2$ 和 $x'_2$ 见证 $u'_2$ 与 $\hat x_2$
之间的相容性。略去详细验证。
\end{proof}

\begin{lemma}
\label{artin-lemma-solution}
在情形 \ref{artin-situation-contractions} 中，若存在解
$(f : X' \to X, T, a)$，则在 $S$ 上局部 Noether 概形 $V$ 的范畴上有函子性
双射 $F(V) = \Mor_S(V, X)$。
\end{lemma}

\begin{proof}
设 $V$ 是 $S$ 上局部 Noether 概形，$x : V \to X$ 是 $S$ 上的态射。
于是按如下方式得到 $F(V)$ 中的元素 $(Z, u', \hat x)$：
\begin{enumerate}
\item $Z \subset V$ 是 $T$ 在 $x$ 下的逆像；
\item $u' : V \setminus Z \to U' = U$ 是 $x$ 到 $V \setminus Z$ 的限制；
\item $\hat x : V_{/Z} \to W$ 是 $x_{/Z} : V_{/Z} \to X_{/T}$
与同构 $a : X_{/T} \to W$ 的复合。
\end{enumerate}
此三元组满足相容性条件，因为可以取 $V' = V \times_{x, X} X'$，并取
$\hat x'$ 为投影 $x' : V' \to X'$ 的完备化。
% P11-E0303: the frozen source omits "to be" or "as" after "we can take \hat x'".

\medskip\noindent
反之，设给定 $F(V)$ 的元素 $(Z, u', \hat x)$。我们断言，存在唯一的态射
$x : V \to X$ 与 $u'$ 和 $\hat x$ 相容。事实上，设
$V' \to V$、$\hat x'$ 和 $x'$ 见证 $u'$ 与 $\hat x$ 之间的相容性。
《形式空间的代数化》命题 \ref{restricted-proposition-glue-modification}
恰好给出我们所需的唯一态射 $x : V \to X$：它在 $V_{/Z}$ 上与 $\hat x$
一致，在 $V'$ 上与 $x'$ 一致（从而在 $V \setminus Z$ 上与 $u'$ 一致）。

\medskip\noindent
略去验证上述两个构造在各自定义域之间给出互逆双射。
\end{proof}

\begin{lemma}
\label{artin-lemma-functor-is-solution}
在情形 \ref{artin-situation-contractions} 中，若存在 $S$ 上局部有限型的代数空间
$X$，并且在 $S$ 上局部 Noether 概形 $V$ 的范畴上有函子性双射
$F(V) = \Mor_S(V, X)$，则 $X$ 是一个解。
\end{lemma}

\begin{proof}
我们必须构造真态射 $f : X' \to X$、闭子集 $T \subset |X|$ 以及同构
$a : X_{/T} \to W$，使其具有情形 \ref{artin-situation-contractions} 之后列出的
性质 (1)、(2)、(3)。

\medskip\noindent
本证明中的讨论略显拘泥，因为我们想仔细匹配底范畴。本段说明愿意略去这些细节的
读者如何更大胆地推进。读者可以把函子 $F$ 的定义扩张到 $S$ 上所有局部
Noether 代数空间。这样便可断言 $F$ 与 $X$ 作为这些代数空间范畴上的函子一致，
即 $X$ 表示 $F$。然后考虑 $F(X)$ 中的通用对象 $(T, u', \hat x)$。
读者会发现：对于见证 $u'$ 与 $\hat x$ 之间相容性的三元组
$X'' \to X$、$\hat x'$、$x'$，态射 $x' : X'' \to X'$ 是同构；将 $x'$
取逆便得到 $f : X' \to X$。最后，我们已有 $T \subset |X|$，而读者可以证明
$\hat x$ 是同构，它可作为最后一个要素，即 $a$。
% P11-E0304: the frozen source has "which can served" instead of "which can serve".

\medskip\noindent
以 $h_X(-) = \Mor_S(-, X)$ 表示 $X$ 的点函子到一节
\ref{artin-section-noetherian} 的范畴 $(\textit{Noetherian}/S)_\etale$ 的限制。
% P11-E0305: the frozen source uses the incomplete naming construction "Denote h_X(...) the functor".
由《空间的极限》注 \ref{spaces-limits-remark-limit-preserving}，代数空间 $X$
和 $X'$ 保极限，故其限制 $h_X$ 和 $h_{X'}$ 也保极限。因此，为构造 $f$，
只需构造变换 $h_{X'} \to h_X = F$；参见引理
\ref{artin-lemma-canonical-extension}。于是，设 $V \to S$ 是
$(\textit{Noetherian}/S)_\etale$ 的对象，且
$\tilde x : V \to X'$ 属于 $h_{X'}(V)$。按如下方式得到 $F(V)$ 的元素
$(Z, u', \hat x)$：
\begin{enumerate}
\item $Z \subset V$ 是 $T'$ 在 $\tilde x$ 下的逆像；
\item $u' : V \setminus Z \to U'$ 是 $\tilde x$ 到 $V \setminus Z$ 的限制；
\item $\hat x : V_{/Z} \to W$ 是
$x_{/Z} : V_{/Z} \to X'_{/T'}$ 与 $g : X'_{/T'} \to W$ 的复合。
% P11-E0306: the frozen source uses x_{/Z}, although the morphism introduced here is \tilde x.
\end{enumerate}
此三元组满足相容性条件：首先，我们总能自然得到 $V' \to V$ 和
$\hat x' : V'_{/Z'} \to X'_{/T'}$（参见引理 \ref{artin-lemma-functor} 之前的讨论）。
% P11-E0307: the frozen source introduces Z' only in V'_{/Z'} here; the construction elsewhere uses the inverse image of Z without this notation.
然后只需把 $x' : V' \to X'$ 定义为 $V' \to V$ 与态射
$\tilde x : V \to X'$ 的复合。略去验证此构造确实可行。

\medskip\noindent
若 $\xi : V \to X$ 是 \'etale 态射且 $V$ 是概形，则得到
$\xi = (Z, u', \hat x) \in F(V) = h_X(V) = X(V)$。当然，若
$\varphi : V' \to V$ 还是概形的一个 \'etale 态射，则
$(Z, u', \hat x)$ 到 $F(V')$ 的拉回对应于 $\xi \circ \varphi$。
闭子集 $T \subset |X|$ 就定义为满足如下条件的闭子集：当
$\xi = (Z, u', \hat x)$ 时，$\xi : V \to X$ 把 $T$ 拉回到 $Z$。
% P11-E0308: the frozen source omits terminal punctuation after "pulls T back to Z".

\medskip\noindent
考虑 $S$ 上 Noether 概形 $V$，以及如上对应于 $(Z, u', \hat x)$ 的态射
$\xi : V \to X$。可见，$\xi(V)$ 在集合论意义下包含于 $T$，当且仅当
$V = Z$（作为拓扑空间）。因此，$X_{/T}$ 与 $W$ 作为函子一致。这给出同构
$a : X_{/T} \to W$。（这里略去了一个小细节：对局部 Noether 形式代数空间
$X_{/T}$ 和 $W$，只需检验它们在 $S$ 上局部 Noether 概形处取值后得到同构。）

\medskip\noindent
略去条件 (1)、(2)、(3) 的证明。
\end{proof}

\begin{remark}
\label{artin-remark-diagonal}
在情形 \ref{artin-situation-contractions} 中，设 $V$ 是 $S$ 上局部 Noether 概形。
% P11-E0309: the frozen source leaves "In Situation ..." as a sentence fragment.
对 $i = 1, 2$，设 $(Z_i, u'_i, \hat x_i) \in F(V)$，并设
$V'_i \to V$、$\hat x'_i$ 和 $x'_i$ 见证 $u'_i$ 与 $\hat x_i$ 之间的相容性
（$i = 1, 2$）。

\medskip\noindent
令 $V' = V'_1 \times_V V'_2$。以 $E' \to V'$ 表示下列两个态射的等化子：
$$
V' \to V'_1 \xrightarrow{x'_1} X'
\quad\text{且}\quad
V' \to V'_2 \xrightarrow{x'_2} X'
$$
令 $Z = Z_1 \cap Z_2$。以 $E_W \to V_{/Z}$ 表示下列两个态射的等化子：
$$
V_{/Z} \to V_{/Z_1} \xrightarrow{\hat x_1} W
\quad\text{且}\quad
V_{/Z} \to V_{/Z_2} \xrightarrow{\hat x_2} W
$$
注意，$E' \to V$ 是分离且局部有限型的，而 $E_W$ 是在 $V$ 上分离的局部
Noether 形式代数空间。所涉及各态射之间的相容性表明：
\begin{enumerate}
\item $\Im(E' \to V) \cap (Z_1 \cup Z_2)$ 包含于 $Z = Z_1 \cap Z_2$；
\item 态射 $E' \times_V (V \setminus Z) \to V \setminus Z$
是单态射，并且等于 $u'_1$ 与 $u'_2$ 到
$V \setminus (Z_1 \cup Z_2)$ 的限制的等化子；
\item 态射 $E'_{/Z} \to V_{/Z}$ 通过 $E_W$ 分解，并且图
$$
\xymatrix{
E'_{/Z} \ar[r] \ar[d] & X'_{/T'} \ar[d]^g \\
E_W \ar[r] & W
}
$$
为 Cartesian 图。特别地，态射 $E'_{/Z} \to E_W$ 作为 $g$ 的基变换，是形式修正；
\item $E'$、$(E' \to V)^{-1}Z$ 和 $E'_{/Z} \to E_W$ 构成情形
\ref{artin-situation-contractions} 中的三元组，其基概形为局部 Noether 概形 $V$；
% P11-E0310: the frozen source uses "E', ..., and ... is a triple"; "form a triple" is grammatically clearer.
\item 给定局部 Noether 概形的态射 $\varphi : A \to V$，下列条件等价：
\begin{enumerate}
\item $(Z_1, u'_1, \hat x_1)$ 和 $(Z_2, u'_2, \hat x_2)$
限制为 $F(A)$ 的同一个元素；
\item $A \setminus \varphi^{-1}(Z) \to V \setminus Z$
通过 $E' \times_V (V \setminus Z)$ 分解，且
$A_{/\varphi^{-1}(Z)} \to V_{/Z}$ 通过 $E_W$ 分解。
\end{enumerate}
\end{enumerate}
由引理 \ref{artin-lemma-solution} 和 \ref{artin-lemma-functor-is-solution} 可知：
若 (4) 中的三元组有一个解 $E \to V$，则在 $V$ 上局部 Noether 概形的范畴上，
$E$ 表示 $F \times_{\Delta, F \times F} V$。
\end{remark}

\begin{lemma}
\label{artin-lemma-closed-immersion}
在情形 \ref{artin-situation-contractions} 中，假设给定闭子集 $Z \subset S$，使得
\begin{enumerate}
\item $Z$ 在 $X'$ 中的逆像是 $T'$；
\item $U' \to S \setminus Z$ 是闭浸入；
\item $W \to S_{/Z}$ 是闭浸入。
\end{enumerate}
则存在解 $(f : X' \to X, T, a)$，而且 $X \to S$ 是闭浸入。
\end{lemma}

\begin{proof}
假设有闭子概形 $X \subset S$，使得
$X \cap (S \setminus Z) = U'$ 且 $X_{/Z} = W$。则 $X$ 表示函子 $F$
（略去一些细节），因而是一个解。寻找 $X$ 显然是 $S$ 上的局部问题。
这样便约化到下一段讨论的情形。

\medskip\noindent
假设 $S = \Spec(A)$ 为仿射概形。设 $I \subset A$ 是定义 $Z$ 的根理想，
写作 $I = (f_1, \ldots, f_r)$。由假设，给定
\begin{enumerate}
\item 闭浸入 $U' \to S \setminus Z$ 确定理想
$J_i \subset A[1/f_i]$，使得 $J_i$ 和 $J_j$ 在 $A[1/f_if_j]$ 中生成同一理想；
\item 闭浸入 $W \to S_{/Z}$ 是映射
$\text{Spf}(A^\wedge/J') \to \text{Spf}(A^\wedge)$，其中
$J' \subset A^\wedge$ 是理想，而 $A^\wedge$ 是 $A$ 的 $I$-adic 完备化。
\end{enumerate}
为完成证明，需要找到理想 $J \subset A$，使得 $J_i = J[1/f_i]$ 且
$J' = JA^\wedge$。由《代数续篇》命题
\ref{more-algebra-proposition-equivalence}，只需证明对所有 $i$，
$J_i$ 和 $J'$ 在 $A^\wedge[1/f_i]$ 中生成同一理想。

\medskip\noindent
回忆，$A' = H^0(X', \mathcal{O})$ 是有限 $A$-代数，且其构造与平坦基变换交换
（《空间的上同调》引理
\ref{spaces-cohomology-lemma-proper-over-affine-cohomology-finite} 和
\ref{spaces-cohomology-lemma-flat-base-change-cohomology}）。令
$J'' = \Ker(A \to A')$\footnote{与读者可能预期的不同，理想 $J$ 和 $J''$
一般并不相等。}。
% P11-E0311: the frozen source uses the incomplete naming construction "Denote J'' = ...".
对 $A[1/f_i]$ 的谱作基变换便得到 $J_i = J''A[1/f_i]$。注意，有交换图
$$
\xymatrix{
X' \ar[d] &
X'_{/T'} \times_{S_{/Z}} \text{Spf}(A^\wedge) \ar[l] \ar[d] &
X'_{/T'} \times_W \text{Spf}(A^\wedge/J') \ar@{=}[l] \ar[d] \\
\Spec(A) &
\text{Spf}(A^\wedge) \ar[l] &
\text{Spf}(A^\wedge/J') \ar[l]
}
$$
中间的竖直箭头是左侧竖直箭头沿显然闭子集的完备化。由形式函数定理，有
$$
(A')^\wedge = \Gamma(X' \times_S \Spec(A^\wedge), \mathcal{O}) =
\lim H^0(X' \times_S \Spec(A/I^n), \mathcal{O})
$$
；参见《空间的上同调》定理
\ref{spaces-cohomology-theorem-formal-functions}。由该图可知，$J'$ 在
$(A')^\wedge$ 中映为零。因此 $J' \subset J'' A^\wedge$。考虑箭头
$$
X'_{/T'} \to
\text{Spf}(A^\wedge/J''A^\wedge) \to
\text{Spf}(A^\wedge/J') = W
$$
已知复合 $g$ 是形式修正（特别地，它是 rig-\'etale 且 rig-满射的），而第二个
箭头是闭浸入（特别地，它是 adic 单态射）。因此
$X'_{/T'} \to \text{Spf}(A^\wedge/J''A^\wedge)$ 是 rig-满射且
rig-\'etale 的；参见《形式空间的代数化》引理
\ref{restricted-lemma-rig-surjective-alternative-permanence} 和
\ref{restricted-lemma-rig-etale-alternative-permanence}。
应用《形式空间的代数化》引理
\ref{restricted-lemma-rig-etale-descent} 和
\ref{restricted-lemma-permanence-rig-surjective}
可知，$\text{Spf}(A^\wedge/J''A^\wedge) \to W$ 是 rig-\'etale 且 rig-满射的。
由《形式空间的代数化》引理
\ref{restricted-lemma-closed-immersion-rig-smooth-rig-surjective}
可知，对某个 $n > 0$ 有 $I^n J'' A^\wedge \subset J'$。从而
$J'' A^\wedge[1/f_i] = J' A^\wedge[1/f_i]$，于是对所有 $i$ 都有
$J_i A^\wedge[1/f_i] = J' A^\wedge[1/f_i]$，正如所需。
\end{proof}

\begin{lemma}
\label{artin-lemma-diagonal-contractions}
在情形 \ref{artin-situation-contractions} 中，假设 $X' \to S$ 和 $W \to S$ 分离。
则对角态射 $\Delta : F \to F \times F$ 可由闭浸入表示。
\end{lemma}

\begin{proof}
结合引理 \ref{artin-lemma-closed-immersion} 与注 \ref{artin-remark-diagonal} 中的讨论。
\end{proof}

\begin{lemma}
\label{artin-lemma-sheaf}
在情形 \ref{artin-situation-contractions} 中，函子 $F$ 对 $S$ 上局部 Noether 概形的
所有 \'etale 覆盖都满足层条件。
\end{lemma}

\begin{proof}
略。提示：态射可以在 \'etale 局部定义。
\end{proof}

\begin{lemma}
\label{artin-lemma-limit-preserving}
在情形 \ref{artin-situation-contractions} 中，函子 $F$ 保极限：对 $S$ 上 Noether
仿射概形的任意有向极限 $V = \lim V_\lambda$，都有
$F(V) = \colim F(V_\lambda)$。
\end{lemma}

\begin{proof}
这是一个长得离谱的证明。其中很大一部分由关于极限和 \'etale 局部化的标准论证
组成。我们建议读者直接跳到证明的最后部分，那里才会出现有趣的内容。

\medskip\noindent
设 $V = \lim_{\lambda \in \Lambda} V_i$ 是 $S$ 上概形的有向极限，其中 $V$
和 $V_\lambda$ 都是 Noether 概形，转移态射为仿射态射。
% P11-E0312: the frozen source writes V_i in a limit indexed by lambda and then consistently uses V_lambda.
关于概形极限的材料，参见《极限》一节 \ref{limits-section-limits}。
我们将证明 $\colim F(V_\lambda) \to F(V)$ 为双射。

\medskip\noindent
单射性证明：记号。设 $\lambda \in \Lambda$，并设
$\xi_{\lambda, 1}, \xi_{\lambda, 2} \in F(V_\lambda)$ 是限制为 $F(V)$
中同一个元素的两个元素。写作
$\xi_{\lambda, 1} = (Z_{\lambda, 1}, u'_{\lambda, 1},
\hat x_{\lambda, 1})$ 以及
$\xi_{\lambda, 2} = (Z_{\lambda, 2}, u'_{\lambda, 2}, \hat x_{\lambda, 2})$。

\medskip\noindent
单射性证明：$Z_{\lambda, i}$ 的一致。由于 $Z_{\lambda, 1}$ 和
$Z_{\lambda, 2}$ 限制为 $V$ 的同一个闭子集，增大 $i$ 后可假设
$Z_{\lambda, 1} = Z_{\lambda, 2}$；参见《极限》引理
\ref{limits-lemma-inverse-limit-top} 和《拓扑》引理
\ref{topology-lemma-describe-limits}。
% P11-E0313: the frozen source says "after increasing i"; the directed-system index here is lambda.
以 $Z_\lambda \subset V_\lambda$ 表示这个共同值；对
$\mu \geq \lambda$，以 $Z_\mu \subset V_\mu$ 表示其在 $V_\mu$ 中的逆像，
并以 $Z$ 表示其在 $V$ 中的逆像。
% P11-E0314: the frozen source duplicates "and" before "denote Z".
下面将使用如下事实：若把 $Z$ 和 $Z_\mu$ 看作既约闭子概形，则作为概形有
$Z = \lim_{\mu \geq \lambda} Z_\mu$。

\medskip\noindent
单射性证明：$u'_{\lambda, i}$ 的一致。由于 $U'$ 在 $S$ 上局部有限型，
且 $u'_{\lambda, 1}$ 与 $u'_{\lambda, 2}$ 到 $V \setminus Z$ 的限制相同，
增大 $\lambda$ 后可假设 $u'_{\lambda, 1} = u'_{\lambda, 2}$；参见《极限》命题
\ref{limits-proposition-characterize-locally-finite-presentation}。
以 $u'_\lambda$ 表示共同值，并以 $u'$ 表示其到 $V \setminus Z$ 的限制。

\medskip\noindent
单射性证明：重述。此时有
$\xi_{\lambda, 1} = (Z_\lambda, u'_\lambda, \hat x_{\lambda, 1})$，以及
$\xi_{\lambda, 2} = (Z_\lambda, u'_\lambda, \hat x_{\lambda, 2})$。
本部分证明面临的主要问题是证明：增大 $\lambda$ 后，态射
$\hat x_{\lambda, 1}$ 与 $\hat x_{\lambda, 2}$ 变为相同。

\medskip\noindent
单射性证明：$\hat x_{\lambda, i}|_{Z_\lambda}$ 的一致。考虑态射
% P11-E0315: the frozen source misspells "morphisms" as "morpisms".
$\hat x_{\lambda, 1}|_{Z_\lambda}, \hat x_{\lambda, 2}|_{Z_\lambda} :
Z_\lambda \to W_{red}$.
这些态射限制为同一个态射 $Z \to W_{red}$。由于 $W_{red}$ 是 $S$ 上局部
有限型的概形，由《极限》命题
\ref{limits-proposition-characterize-locally-finite-presentation}
可知，把 $\lambda$ 换成更大的指标后，可假设
$\hat x_{\lambda, 1}|_{Z_\lambda} = \hat x_{\lambda, 2}|_{Z_\lambda} :
Z_\lambda \to W_{red}$.

\medskip\noindent
单射性证明：结束。接下来，把注 \ref{artin-remark-diagonal} 中的讨论应用于
$V_\lambda$ 以及两个元素
$\xi_{\lambda, 1}, \xi_{\lambda, 2} \in F(V_\lambda)$，得到
\begin{enumerate}
\item $e_\lambda : E_\lambda' \to V_\lambda$
分离且局部有限型；
\item $e_\lambda^{-1}(V_\lambda \setminus Z_\lambda) \to
V_\lambda \setminus Z_\lambda$ 是同构；
\item 单态射 $E_{W, \lambda} \to V_{\lambda, /Z_\lambda}$，它是
$\hat x_{\lambda, 1}$ 与 $\hat x_{\lambda, 2}$ 的等化子；
\item 形式修正 $E'_{\lambda, /Z_\lambda} \to E_{W, \lambda}$。
\end{enumerate}
由注 \ref{artin-remark-diagonal} 的断言 (2) 以及上面得到
$u'_{\lambda, 1} = u'_{\lambda, 2} = u'_\lambda$ 的准备工作，断言 (2) 成立。
由于
$\hat x_{\lambda, 1}|_{Z_\lambda} = \hat x_{\lambda, 2}|_{Z_\lambda}$
，$Z_\lambda = (V_{\lambda, /Z_\lambda})_{red}$ 通过
$E_{W, \lambda}$ 分解。因此，由《形式空间》引理
\ref{formal-spaces-lemma-monomorphism-iso-over-red} 可知，
$E_{W, \lambda} \to V_{\lambda, /Z_\lambda}$ 是闭浸入。再把注
\ref{artin-remark-diagonal} 的断言 (4) 和引理 \ref{artin-lemma-closed-immersion}
应用于 $V_\lambda$ 上的三元组
$E_\lambda'$、$e_\lambda^{-1}(Z_\lambda)$、
$E'_{\lambda, /Z_\lambda} \to E_{W, \lambda}$，可知存在闭浸入
$E_\lambda \to V_\lambda$，它是此三元组的解。
接下来，使用注 \ref{artin-remark-diagonal} 的断言 (5)，并结合引理
\ref{artin-lemma-solution}，可知 $E_\lambda$ 是 $\xi_{\lambda, 1}$ 与
$\xi_{\lambda, 2}$ 的“等化子”。特别地，$V \to V_\lambda$ 通过
$E_\lambda$ 分解。再次使用《极限》命题
\ref{limits-proposition-characterize-locally-finite-presentation}
，可找到 $\mu \geq \lambda$，使得 $V_\mu \to V_\lambda$ 通过
$E_\lambda$ 分解，因而 $\xi_{\lambda, 1}$ 和 $\xi_{\lambda, 2}$ 到
$V_\mu$ 的拉回相同，正如所需。

\medskip\noindent
满射性证明：陈述。设 $\xi = (Z, u', \hat x)$ 是 $F(V)$ 的元素。
必须找到 $\lambda \in \Lambda$ 和元素 $\xi_\lambda \in F(V_\lambda)$，
使其限制为 $\xi$。

\medskip\noindent
满射性证明：问题是 \'etale 局部的。由证明前一部分得到的唯一性以及引理
\ref{artin-lemma-sheaf} 中 $F$ 的层条件，此问题在 \'etale 拓扑下是 $V$ 上局部的。
更确切地，设 $v \in V$。我们断言，只需找到 \'etale 态射
$(\tilde V, \tilde v) \to (V, v)$、某个 $\lambda$、某个 \'etale 态射
$\tilde V_\lambda \to V_\lambda$，以及某个元素
$\tilde \xi_\lambda \in F(\tilde V_\lambda)$，使得
% P11-E0316: the frozen source has the malformed phrase "some an etale morphism".
$\tilde V = \tilde V_\lambda \times_{V_\lambda} V$
且 $\xi|_U = \tilde \xi_\lambda|_U$。
% P11-E0317: the frozen source uses U in both restrictions without defining U in this paragraph; the intended comparison appears to be after pullback to \tilde V.
略去此断言的详细证明\footnote{为证明这一点，把若干态射
$\tilde V \to V$ 组成一个有限 \'etale 覆盖，并证明相应态射
$\tilde V_\lambda \to V_\lambda$ 在增大 $\lambda$ 后也组成 \'etale 覆盖。
然后利用单射性看出诸元素 $\tilde \xi_\lambda$ 在增大 $\lambda$ 后可以黏合，
并利用 $F$ 的层条件把这些元素下降为 $F(V_\lambda)$ 的一个元素。}。

\medskip\noindent
满射性证明：重述问题。回忆，若 \'etale 态射
$(\tilde V, \tilde v) \to (V, v)$ 中的 $\tilde V$ 为仿射概形，则对某个
$\lambda$，它是某个 \'etale 态射 $\tilde V_\lambda \to V_\lambda$
的基变换，其中 $\tilde V_\lambda$ 为仿射概形；例如参见《拓扑》引理
\ref{topologies-lemma-limit-fppf-topology}。给定 $\tilde V_\lambda$，有
$\tilde V = \lim_{\mu \geq \lambda} \tilde V_\lambda \times_{V_\lambda} V_\mu$.
因此，给定 \'etale 态射 $(\tilde V, \tilde v) \to (V, v)$ 且
$\tilde V$ 仿射后，可以用 $(\tilde V, \tilde v)$ 替换 $(V, v)$，并用
$\xi$ 到 $\tilde V$ 的限制替换 $\xi$。

\medskip\noindent
满射性证明：约化到基为仿射的情形。特别地，设 $\tilde S \subset S$ 是仿射开
子概形，使得 $v \in V$ 映到 $\tilde S$ 的一个点。由上一段，可以用
$\tilde V = \tilde S \times_S V$ 替换 $V$。当然，此时只需对 $F$ 到
$\tilde S$ 上局部 Noether 概形范畴的限制解决问题。这个函子正是整个情形基变换到
$\tilde S$ 后所对应的函子。因此，在证明的余下部分中可以并且确实假设
$S = \Spec(R)$ 是 Noether 仿射概形。

\medskip\noindent
满射性证明：容易的情形。若 $v \in V \setminus Z$，则可取
$\tilde V = V \setminus Z$。由《极限》引理 \ref{limits-lemma-descend-opens}，
它下降为某个 $\lambda$ 下的开子概形
$\tilde V_\lambda \subset V_\lambda$。接着，增大 $\lambda$ 后可假设存在态射
$u'_\lambda : \tilde V_\lambda \to U'$，其限制为 $u'$。取
$\tilde \xi_\lambda = (\emptyset, u'_\lambda, \emptyset)$，便得到所需的
$F(\tilde V_\lambda)$ 元素。

\medskip\noindent
满射性证明：困难情形及到仿射 $W$ 的约化。最困难的情形来自考虑
$v \in Z \subset V$。我们断言可以约化到 $W$ 为仿射形式概形的情形；
建议读者跳过此论证\footnote{Artin 证明本引理的方法是绕过这一步，因而可以避免
先证明单射性。具体地，在证明过程中，Artin 始终使用所涉及所有空间的有限仿射
\'etale 覆盖，并追踪它们之间的映射。事后看来，这或许比我们这里的做法更可取。}。
% P11-E0318: the frozen source says "a finite affine etale coverings", mixing a singular article with a plural noun.
事实上，可以选取 \'etale 态射 $\tilde W \to W$，其中 $\tilde W$ 是仿射形式
代数空间，使得 $v$ 在 $\hat x : V_{/Z} \to W$ 下的像包含于
$\tilde W \to W$ 在既约空间上的像。于是态射
$$
p : \tilde W \times_{W, g} X'_{/T'} \longrightarrow X'_{/T'}
$$
以及
$$
q : \tilde W \times_{W, \hat x} V_{/Z} \to V_{/Z}
$$
都是局部 Noether 形式代数空间的 \'etale 态射。由《形式空间的代数化》定理
\ref{restricted-theorem-dilatations-general}
的一个容易情形，存在代数空间态射 $\tilde X' \to X'$；它局部有限型，在
$U'$ 上为同构，并且 $\tilde X'_{/T'} \to X'_{/T'}$ 与 $p$ 同构。
由《形式空间的代数化》引理 \ref{restricted-lemma-output-etale}，态射
$\tilde X' \to X'$ 是 \'etale 的。以
$\tilde T' \subset |\tilde X'|$ 表示 $T'$ 的逆像，以
$\tilde U' \subset \tilde X'$ 表示其补开子空间，并以
$\tilde g' : \tilde X'_{/\tilde T'} \to \tilde W$ 表示由
$\tilde W \to W$ 对 $g$ 作基变换所得的形式修正。
% P11-E0319: the frozen source uses three consecutive incomplete "Denote ... the ..." constructions.
于是
$$
\tilde X',\ \tilde T',\ \tilde U',\ \tilde W,
\ \tilde g : \tilde X'_{/\tilde T'} \to \tilde W
$$
是情形 \ref{artin-situation-contractions} 的另一个实例。
% P11-E0320: the frozen source defines \tilde g' in prose but uses \tilde g in the displayed tuple and subsequently.
以 $\tilde F$ 表示由此三元组构造的函子。
% P11-E0321: the frozen source uses the incomplete construction "Denote \tilde F the functor".
有函子变换
$$
\tilde F \longrightarrow F
$$
，它以显然的方式利用态射 $\tilde X' \to X'$ 和 $\tilde W \to W$ 构造；
略去细节。

\medskip\noindent
满射性证明：困难情形及到仿射 $W$ 的约化，第 2 部分。由上面使用的同一定理，
存在代数空间态射 $\tilde V \to V$；它局部有限型，在 $V \setminus Z$ 上为同构，
并且 $\tilde V_{/Z} \to V_{/Z}$ 与 $q$ 同构。以
$\tilde Z \subset \tilde V$ 表示 $Z$ 的逆像。
% P11-E0322: the frozen source uses the incomplete naming construction "Denote \tilde Z ... the inverse image".
由《形式空间的代数化》引理
\ref{restricted-lemma-output-etale} 和 \ref{restricted-lemma-output-separated}
，态射 $\tilde V \to V$ 是 \'etale 且分离的。特别地，$\tilde V$ 是（局部
Noether）概形；例如参见《空间的态射》命题
\ref{spaces-morphisms-proposition-locally-quasi-finite-separated-over-scheme}.
已有态射 $u'$，可将其视为态射
$$
\tilde u' : \tilde V \setminus \tilde Z \longrightarrow \tilde U'
$$
，其中 $\tilde U' \subset \tilde X'$ 是同构地映到 $U'$ 的开集。还有态射
$$
\tilde {\hat x} :
\tilde V_{/\tilde Z} = \tilde W \times_{W, \hat x} V_{/Z}
\longrightarrow
\tilde W
$$
，这里使用的就是投影。因此有三元组
$$
\tilde \xi = (\tilde Z, \tilde u', \tilde {\hat x}) \in \tilde F(\tilde V)
$$
略去相容性条件的证明。提示：若 $V' \to V$、$\hat x'$ 和 $x'$ 见证
$u'$ 与 $\hat x$ 之间的相容性，则令
$\tilde V' = V' \times_V \tilde V$；它带有态射
$\tilde{\hat x}'$ 和 $\tilde x'$，而这些数据确实可行。
% P11-E0323: the frozen source misspells "morphisms" as "morphsms" and uses "show" where the subject requires "shows".
$\tilde \xi$ 在变换 $\tilde F \to F$ 下的像是 $\xi$ 到 $\tilde V$ 的限制。

\medskip\noindent
满射性证明：困难情形及到仿射 $W$ 的约化，第 3 部分。由我们对
$\tilde W \to W$ 的选择，存在仿射开集
$\tilde V_{open} \subset \tilde V$（记号快不够用了），它在 $V$ 中的像包含
所选点 $v \in V$。由下一段研究的情形及前面的说明，可以把
$\tilde \xi|_{\tilde V_{open}}$ 下降为 $\tilde F$ 在
$\tilde V_{\lambda, open}$ 上的某个元素 $\tilde \xi_\lambda$；这里有某个
\'etale 态射 $\tilde V_{\lambda, open} \to V_\lambda$，其到 $V$ 的基变换为
$\tilde V_{open}$。应用函子变换 $\tilde F \to F$，便得到我们所求的
$F(\tilde V_{\lambda, open})$ 元素。这样便约化到下一段讨论的情形。

\medskip\noindent
满射性证明：$W$ 仿射的情形。此时 $v \in Z \subset V$，且 $W$ 是仿射形式
代数空间。回忆，
$$
\xi = (Z, u', \hat x) \in F(V)
$$
仍可用 $v$ 的 \'etale 邻域替换 $V$。特别地，可以并且确实假设 $V$ 和
$V_\lambda$ 都是仿射的。

\medskip\noindent
满射性证明：下降 $Z$。可以找到 $\lambda$ 以及闭子概形
$Z_\lambda \subset V_\lambda$，使得 $Z$ 是 $Z_\lambda$ 到 $V$ 的基变换；
参见《极限》引理 \ref{limits-lemma-descend-finite-presentation}。警告：我们不知道
$Z_\lambda$ 是否是 $V_\lambda$ 的既约闭子概形（一般而言也并非如此）。
对 $\mu \geq \lambda$，以 $Z_\mu \subset V_\mu$ 表示其在 $V_\mu$ 中的
概形论逆像。下面将使用作为概形的等式
$Z = \lim_{\mu \geq \lambda} Z_\mu$。

\medskip\noindent
满射性证明：下降 $u'$。由于 $U'$ 在 $S$ 上局部有限型，增大 $\lambda$ 后
可假设存在态射 $u'_\lambda : V_\lambda \setminus Z_\lambda \to U'$，
其到 $V \setminus Z$ 的限制为 $u'$；参见《极限》命题
\ref{limits-proposition-characterize-locally-finite-presentation}。
对 $\mu \geq \lambda$，以 $u'_\mu$ 表示 $u'_\lambda$ 到
$V_\mu \setminus Z_\mu$ 的限制。

\medskip\noindent
满射性证明：下降一个见证。设 $V' \to V$、$\hat x'$ 和 $x'$ 见证
$u'$ 与 $\hat x$ 之间的相容性。使用与上面相同的参考结果，增大
$\lambda$ 后可假设存在有限型态射 $V'_\lambda \to V_\lambda$，它到 $V$
的基变换为 $V' \to V$。再次增大 $\lambda$ 后，可假设
$V'_\lambda \to V_\lambda$ 为真态射（《极限》引理
\ref{limits-lemma-eventually-proper}）。接着，可假设
$V'_\lambda \to V_\lambda$ 在 $V_\lambda \setminus Z_\lambda$ 上为同构
（《极限》引理 \ref{limits-lemma-descend-isomorphism}）。再者，可假设存在态射
$x'_\lambda : V'_\lambda \to X'$，其到 $V'$ 的限制为 $x'$。再次增大
$\lambda$，可假设 $x'_\lambda$ 在 $V_\lambda \setminus Z_\lambda$
上与 $u'_\lambda$ 一致。对 $\mu \geq \lambda$，以 $V'_\mu$ 表示
$V'_\lambda$ 的基变换，并以 $x'_\mu$ 表示 $x'_\lambda$ 的限制。

\medskip\noindent
满射性证明：代数。写作 $W = \text{Spf}(B)$、$V = \Spec(A)$，并且对
$\mu \geq \lambda$ 写作 $V_\mu = \Spec(A_\mu)$。以
$I_\mu \subset A_\mu$ 和 $I \subset A$ 分别表示定义 $Z_\mu$ 和 $Z$ 的理想。
% P11-E0324: the frozen source uses the incomplete construction "Denote I_mu ... and I ... the ideals".
于是 $I_\lambda A_\mu = I_\mu$ 且 $I_\lambda A = I$。态射 $\hat x$
确定且由如下连续环同态确定：
$$
(\hat x)^\sharp : B \longrightarrow A^\wedge
$$
，其中 $A^\wedge$ 是 $A$ 的 $I$-adic 完备化。为完成证明，必须证明：
对某个充分大的 $\mu$，此映射下降为取值于 $A_\mu^\wedge$ 的映射，其中
$A_\mu^\wedge$ 是 $A_\mu$ 的 $I_\mu$-adic 完备化。这并非平凡事实；
Artin 在论文
\cite{ArtinII} 中写道：“由于数据 (3.5) 涉及 $I$-adic 完备化，而完备化不与
有向极限交换，验证颇为微妙。这是分析中收敛证明的一个代数类比。”

\medskip\noindent
满射性证明：代数，更多环。记
$$
C_\mu = \Gamma(V'_\mu, \mathcal{O})
\quad\text{且}\quad
C = \Gamma(V', \mathcal{O})
$$
注意，$A \to C$ 和 $A_\mu \to C_\mu$ 是有限环同态，因为
$V' \to V$ 和 $V'_\mu \to V_\mu$ 是真态射；参见《空间的上同调》引理
\ref{spaces-cohomology-lemma-proper-over-affine-cohomology-finite}.
由于 $V = \lim V_\mu$ 且 $V' = \lim V'_\mu$，有
$$
A = \colim A_\mu
\quad\text{且}\quad
C = \colim C_\mu
$$
；这是《极限》引理 \ref{limits-lemma-descend-section} 的结论\footnote{由于各个
态射并不平坦，我们不知道是否有
$C_\mu = C_\lambda \otimes_{A_\lambda} A_\mu$。}。对元素
$a \in I$（相应地，$a \in I_\mu$），由平坦基变换，映射 $A_a \to C_a$
（相应地，$(A_\mu)_a \to (C_\mu)_a$）是同构（《空间的上同调》引理
\ref{spaces-cohomology-lemma-flat-base-change-cohomology}）。
因此，$A \to C$ 的核与余核支撑在 $V(I)$ 上；$A_\mu \to C_\mu$ 也类似。
% P11-E0325: the frozen source uses singular "is supported" for the compound subject "kernel and cokernel".
由此可知，$A \to C$ 的核与余核被 $I$ 的某个幂零化，而
$A_\mu \to C_\mu$ 的核与余核被 $I_\mu$ 的某个幂零化；参见《代数》引理
\ref{algebra-lemma-Noetherian-power-ideal-kills-module}。

\medskip\noindent
满射性证明：代数，更多环同态。以 $Z_n \subset V$ 表示 $Z$ 的第 $n$ 阶
无穷小邻域，以 $Z_{\mu, n} \subset V_\mu$ 表示 $Z_\mu$ 的第 $n$ 阶
无穷小邻域。
% P11-E0326: the frozen source uses two incomplete "Denote ... the nth infinitesimal neighbourhood" constructions.
由形式函数定理（《空间的上同调》定理
\ref{spaces-cohomology-theorem-formal-functions})
，有
$$
C^\wedge = \lim_n H^0(V' \times_V Z_n, \mathcal{O})
\quad\text{且}\quad
C_\mu^\wedge =
\lim_n H^0(V'_\mu \times_{V_\mu} Z_{\mu, n}, \mathcal{O})
$$
，其中 $C^\wedge$ 和 $C_\mu^\wedge$ 分别是关于 $I$ 和 $I_\mu$ 的完备化。
% P11-E0327: the frozen source says two objects "are the completion" instead of "are the completions".
把态射 $x'_\mu : V'_\mu \to X'$ 的完备化与态射
$g : X'_{/T'} \to W$ 复合，得到
$$
g \circ x'_{\mu, /Z_\mu} :
V'_{\mu, /Z_\mu} = \colim V_\mu' \times_{V_\mu} Z_{\mu, n}
\longrightarrow
W
$$
，因而由上面对完备化 $C_\mu^\wedge$ 的描述，得到连续环同态
$$
(g \circ x'_{\mu, /Z_\mu})^\sharp : B \longrightarrow C_\mu^\wedge
$$
$V' \to V$、$\hat x'$、$x'$ 见证 $u'$ 与 $\hat x$ 之间的相容性这一事实
蕴含下图交换：
% P11-E0328: the frozen source uses singular "witnesses" for the compound subject V' -> V, \hat x', x'.
$$
\xymatrix{
C_\mu^\wedge \ar[r] &
C^\wedge \\
B \ar[u]^{(g \circ x'_{\mu, /Z_\mu})^\sharp} \ar[r]^{(\hat x)^\sharp} &
A^\wedge \ar[u]
}
$$

\medskip\noindent
满射性证明：更多代数论证。回忆，有限 $A$-模 $\Ker(A \to C)$ 和
$\Coker(A \to C)$ 被 $I$ 的某个幂零化；类似地，有限 $A_\mu$-模
$\Ker(A_\mu \to C_\mu)$ 和 $\Coker(A_\mu \to C_\mu)$ 被 $I_\mu$
的某个幂零化。这说明这些模等于其完备化。由于在有限 $A$-模范畴上作
$I$-adic 完备化是正合的（参见《代数》一节
\ref{algebra-section-completion-noetherian}），故有
$$
\Coker(A^\wedge \to C^\wedge) = \Coker(A \to C)
$$
；核以及映射 $A_\mu \to C_\mu$ 也类似。当然，还有
$$
\Ker(A \to C) = \colim \Ker(A_\mu \to C_\mu)
\quad\text{且}\quad
\Coker(A \to C) = \colim \Coker(A_\mu \to C_\mu)
$$
回忆，$S = \Spec(R)$ 为仿射概形。由于所有态射都是 $S$ 上的态射，上述所有
环同态都是 $R$-代数同态。由《形式空间的代数化》引理
\ref{restricted-lemma-Noetherian-finite-type-red}
可知，$B$ 在 $R$ 上拓扑有限型。设 $B$ 由 $b_1, \ldots, b_n$ 拓扑生成。
取某个 $\mu$（例如 $\lambda$），并考虑下列元素：
$$
\text{images of }
(g \circ x'_{\mu, /Z_\mu})^\sharp(b_1)
, \ldots,
(g \circ x'_{\mu, /Z_\mu})^\sharp(b_n)
\text{ 于 }\Coker(A_\mu \to C_\mu)
$$
由上面方块的交换性，这些元素在 $\Coker(\alpha)$ 中的像为零。
% P11-E0329: the frozen source refers to Coker(alpha), but no morphism alpha is defined in this proof.
由于 $\Coker(A \to C) = \colim \Coker(A_\mu \to C_\mu)$，且这些余核等于
其完备化，增大 $\mu$ 后可假设这些像全都为零。这意味着连续同态
$(g \circ x'_{\mu, /Z_\mu})^\sharp$ 的像包含于
$\Im(A_\mu \to C_\mu)$。选取元素 $a_{\mu, j} \in (A_\mu)^\wedge$，
使其在 $(C_\mu)^\wedge$ 中映到
$(g \circ x'_{\mu, /Z_\mu})^\sharp(b_1)$。
% P11-E0330: the frozen source chooses a_{\mu,j} but maps every one to the expression with b_1; the subsequent sentence uses b_j.
由上面方块的交换性，$a_{\mu, j} \in A_\mu^\wedge$ 和
$(\hat x)^\sharp(b_j) \in A^\wedge$ 映到 $C^\wedge$ 的同一个元素。
由于 $\Ker(A \to C) = \colim \Ker(A_\mu \to C_\mu)$，且这些核等于其
完备化，增大 $\mu$ 后可以调整 $a_{\mu, j}$ 的选择，使
$a_{\mu, j}$ 在 $A^\wedge$ 中的像等于 $(\hat x)^\sharp(b_j)$。

\medskip\noindent
满射性证明：最后的代数论证。设 $\mathfrak b \subset B$ 是拓扑幂零元素所成的
理想。设 $J \subset R[x_1, \ldots, x_n]$ 是由满足
$h(b_1, \ldots, b_n) \in \mathfrak b$ 的那些
$h(x_1, \ldots, x_n)$ 组成的理想。于是得到拓扑 $R$-代数的连续满射
$$
\Phi : R[x_1, \ldots, x_n]^\wedge \longrightarrow B,\quad
x_j \longmapsto b_j
$$
，其中左边关于 $J$ 完备化。由于 $R[x_1, \ldots, x_n]$ 是 Noether 环，
可为 $J$ 选取生成元 $h_1, \ldots, h_m$。由上面方块的交换性，
$h_j(a_{\mu, 1}, \ldots, a_{\mu, n})$ 是 $A_\mu^\wedge$ 的元素，且其在
$A^\wedge$ 中的像包含于 $IA^\wedge$。事实上，环同态
$(\hat x)^\sharp$ 连续，并且因为 $A^\wedge/IA^\wedge = A/I$ 既约，
$IA^\wedge$ 是 $A^\wedge$ 的拓扑幂零元素所成的理想。（关于 Noether 环中
完备化的结果，参见《代数》一节
\ref{algebra-section-completion-noetherian}。）由于
$A/I = \colim A_\mu/I_\mu$，增大 $\mu$ 后可假设
$h_j(a_{\mu, 1}, \ldots, a_{\mu, n})$ 属于 $I_\mu A_\mu^\wedge$。
特别地，$A_\mu^\wedge$ 的元素
$h_j(a_{\mu, 1}, \ldots, a_{\mu, n})$ 在 $A_\mu^\wedge$ 中拓扑幂零。
因此得到连续 $R$-代数同态
$$
\Psi : R[x_1, \ldots, x_n]^\wedge \longrightarrow A_\mu^\wedge,\quad
x_j \longmapsto a_{\mu, j}
$$
为得到所需结论，需要检验 $\Ker(\Phi)$ 是否被 $\Psi$ 零化。这未必立即成立，
但增大 $\mu$ 后可以做到。事实上，由于
$R[x_1, \ldots, x_n]^\wedge$ 是 Noether 环，可以为理想 $\Ker(\Phi)$
选取生成元 $g_1, \ldots, g_l$。于是
$$
\Psi(g_1), \ldots, \Psi(g_l) \in
\Ker(A_\mu^\wedge \to C_\mu^\wedge) = \Ker(A_\mu \to C_\mu)
$$
在 $\Ker(A \to C) = \colim \Ker(A_\mu \to C_\mu)$ 中映为零。
因此，像前面一样增大 $\mu$，便得到所需结果。

\medskip\noindent
满射性证明：收尾。上面构造的连续环同态 $B \to (A_\mu)^\wedge$ 确定态射
$\hat x_\mu : V_{\mu, /Z_\mu} \to W$。$\hat x_\mu$ 与 $u'_\mu$ 的相容性
来自如下事实：按构造，环同态 $B \to (A_\mu)^\wedge$ 与环同态
$A_\mu \to C_\mu$ 相容。事实上，我们在构造中使用的真态射
$V'_\mu \to V_\mu$ 以及态射 $x'_\mu$ 和
$\hat x'_\mu = x'_{\mu, /Z_\mu}$ 将见证这一相容性。证明完成。
\end{proof}

\begin{lemma}
\label{artin-lemma-rs}
在情形 \ref{artin-situation-contractions} 中，函子 $F$ 满足
Rim--Schlessinger 条件 (RS)。
\end{lemma}

\begin{proof}
回忆，此条件只涉及函子 $F$ 在 $S$ 上概形 $V$ 处的取值 $F(V)$，这些概形是
Artin 局部环的谱；以及由 $S$ 上概形态射 $V_1 \to V_2$ 诱导的限制映射
$F(V_2) \to F(V_1)$，这些概形同样是 Artin 局部环的谱。
因此，设 $V/S$ 是某个 Artin 局部环的谱。
% P11-E0331: the frozen source misspells "spectrum" as "spetruim".
若 $\xi = (Z, u', \hat x) \in F(V)$，则在集合论意义下要么
$Z = \emptyset$，要么 $Z = V$。在第一种情形中，$\hat x$ 是从空形式代数空间
到 $W$ 的态射。在第二种情形中，$u'$ 是从空概形到 $X'$ 的态射，而
$\hat x : V \to W$ 是到 $W$ 的态射。因此
$$
F(V) = U'(V) \amalg W(V)
$$
，而且对上述 $V_1 \to V_2$，诱导映射 $F(V_2) \to F(V_1)$ 与此分解相容。
因此，只需证明 $U'$ 和 $W$ 都满足 Rim--Schlessinger 条件。对 $U'$，
这由引理 \ref{artin-lemma-algebraic-stack-RS} 得到。为证明 $W$ 的情形，按
《形式空间》引理 \ref{formal-spaces-lemma-structure-locally-noetherian}
写作 $W = \colim W_n$。设 $V = \Spec(A)$，其中
$(A, \mathfrak m)$ 是 Artin 局部环。取 $n \geq 1$ 使
$\mathfrak m^n = 0$，则 $W(V) = W_n(V)$。故 $W$ 的 Rim--Schlessinger
条件由如下事实推出：对所有 $n$，$W_n$ 都满足 Rim--Schlessinger 条件；
而后者又由引理
\ref{artin-lemma-algebraic-stack-RS} 推出。
\end{proof}

\begin{lemma}
\label{artin-lemma-finite-dim}
在情形 \ref{artin-situation-contractions} 中，函子 $F$ 的切空间是有限维的。
\end{lemma}

\begin{proof}
在引理 \ref{artin-lemma-rs} 的证明中已经看到：若 $V$ 是 Artin 局部环的谱，则
$F(V) = U'(V) \amalg W(V)$。切空间完全由 $F$ 在 $S$ 上此类概形处的取值
计算。因此，只需证明函子 $U'$ 和 $W$ 的切空间有限维。对 $U'$，这由引理
\ref{artin-lemma-finite-dimension} 得到。像引理 \ref{artin-lemma-rs} 的证明中那样写作
$W = \colim W_n$。于是 $W$ 的切空间等于 $W_2$ 的切空间，因为为得到切空间，
只需在满足 $\mathfrak m^2 = 0$ 的 Artin 局部环
$(A, \mathfrak m)$ 的谱处计算 $W$。再次由引理
\ref{artin-lemma-finite-dimension} 可知，$W_2$ 的切空间有限维。
\end{proof}

\begin{lemma}
\label{artin-lemma-formal-object-effective}
在情形 \ref{artin-situation-contractions} 中，假设 $X' \to S$ 分离。则 $F$
的每个形式对象都有效。
\end{lemma}

\begin{proof}
$F$ 的形式对象 $\xi = (R, \xi_n)$ 由如下数据组成：Noether 完备局部
$S$-代数 $R$，其剩余域在 $S$ 上有限型；以及对所有 $n$ 的元素
$\xi_n \in F(\Spec(R/\mathfrak m^n))$，满足
$\xi_{n + 1}|_{\Spec(R/\mathfrak m^n)} = \xi_n$。由引理
\ref{artin-lemma-rs} 证明中的讨论，$\xi$ 要么是 $U'$ 的形式对象，要么是 $W$
的形式对象。在第一种情形中，由引理 \ref{artin-lemma-effective}，$\xi$ 有效。
第二种情形才有趣。令 $V = \Spec(R)$。我们将构造元素
$(Z, u', \hat x) \in F(V)$，使其在
$F(\Spec(R/\mathfrak m^n))$ 中的像对所有 $n \geq 1$ 都是 $\xi_n$。

\medskip\noindent
可以把元素族 $\xi_n$ 看作态射
$$
\xi : \text{Spf}(R) \longrightarrow W
$$
，这是 $S$ 上局部 Noether 形式代数空间的态射。注意，一般而言 $\xi$
{\it 不是} adic 态射。为修正这一点，设 $I \subset R$ 是对应于形式闭子空间
$$
\text{Spf}(R) \times_{\xi, W} W_{red} \subset \text{Spf}(R)
$$
的理想。注意 $I \subset \mathfrak m_R$。令
$Z = V(I) \subset V = \Spec(R)$。由于 $R$ 是 $\mathfrak m_R$-adic
完备的，它更是 $I$-adic 完备的（《代数》引理
\ref{algebra-lemma-complete-by-sub}）。此外，我们断言，对每个 $n \geq 1$，
态射
$$
\xi|_{\text{Spf}(R/I^n)} :
\text{Spf}(R/I^n)
\longrightarrow
W
$$
实际上来自态射
$$
\xi'_n :  \Spec(R/I^n) \longrightarrow W
$$
事实上，像引理 \ref{artin-lemma-rs} 的证明中那样写作 $W = \colim W_n$，
注意 $\xi|_{\text{Spf}(R/I^n)}$ 映入 $W_n$，再应用《形式空间》引理
\ref{formal-spaces-lemma-map-into-algebraic-space}
把它代数化为态射 $\Spec(R/I^n) \to W_n$，正如所需。以
$\text{Spf}'(R) = V_{/Z}$ 表示赋予 $I$-adic 拓扑的 $R$ 的形式谱——
等价地，它是 $V$ 沿 $Z$ 的形式完备化。利用态射 $\xi'_n$，得到 adic 态射
$$
\hat x = (\xi'_n) : \text{Spf}'(R) \longrightarrow W
$$
，它是 $S$ 上局部 Noether 形式代数空间的态射。考虑基变换
$$
\text{Spf}'(R) \times_{\hat x, W, g} X'_{/T'} \longrightarrow \text{Spf}'(R)
$$
由《形式空间的代数化》引理
\ref{restricted-lemma-base-change-formal-modification}，这是形式修正。
因此，由膨胀的主定理（《形式空间的代数化》定理
\ref{restricted-theorem-dilatations}），得到真态射
$$
V' \longrightarrow V = \Spec(R)
$$
，它在 $\Spec(R) \setminus V(I)$ 上为同构，且其完备化恢复上面的形式修正。
换言之，
$$
V' \times_{\Spec(R)} \Spec(R/I^n) =
\Spec(R/I^n) \times_{\xi'_n, W, g} X'_{/T'}
$$
特别地，这说明我们有相容的态射系
$$
V' \times_{\Spec(R)} \Spec(R/I^n) \longrightarrow X' \times_S \Spec(R/I^n)
$$
因此，由 Grothendieck 代数化定理（采用《空间的态射续篇》引理
\ref{spaces-more-morphisms-lemma-algebraize-morphism})
的形式），得到态射
$$
x' : V' \to X'
$$
，它在 $S$ 上并恢复上面显示的态射。最后，令
$u' : V \setminus Z \to X'$ 为 $x'$ 到 $V \setminus Z \subset V'$
的限制，便得到所求元素 $(Z, u', \hat x) \in F(V)$ 的第三个分量。
% P11-E0332: the frozen source uses the incomplete construction "setting u' ... the restriction".
\end{proof}

\begin{lemma}
\label{artin-lemma-openness-smoothness}
设 $S$ 是局部 Noether 概形，$V$ 是 $S$ 上局部有限型的概形，
$Z \subset V$ 为闭集。设 $W$ 是 $S$ 上局部 Noether 的形式代数空间，且
$W_{red}$ 在 $S$ 上局部有限型。设 $g : V_{/Z} \to W$ 是 $S$ 上形式代数
空间的 adic 态射。设 $v \in V$ 是闭点，使得 $g$ 在 $v$ 处通用（含义见一节
\ref{artin-section-axioms-functors}）。则用 $v$ 的一个开邻域替换 $V$ 后，态射
$g$ 光滑（见证明）。
\end{lemma}

\begin{proof}
由于 $g$ 是 adic 的，它可由代数空间表示（《形式空间》一节
\ref{formal-spaces-section-adic}）。因此，所谓 $g$ 光滑，是指 $g$ 按
《自举》定义 \ref{bootstrap-definition-property-transformation} 的意义光滑。

\medskip\noindent
像《形式空间》引理 \ref{formal-spaces-lemma-structure-locally-noetherian}
中那样写作 $W = \colim W_n$。令
$V_n = V_{/Z} \times_{\hat x, W} W_n$。
% P11-E0333: the frozen source uses \hat x in this fibre product, but the morphism introduced in the lemma is g.
则 $V_n$ 是底集合为 $Z$ 的闭子概形。$V \to W$ 光滑等价于所有态射
$V_n \to W_n$ 都光滑；这是因为当 $T$ 是拟紧概形时，任意态射
$T \to W$ 都通过某个 $n$ 所对应的 $W_n$ 分解。
% P11-E0334: the frozen source refers to smoothness of V -> W, although only g: V_{/Z} -> W has been defined.
由引理 \ref{artin-lemma-base-change-versal}，态射 $V_n \to W_n$ 在 $v$ 处光滑
\footnote{该引理适用，因为 $W$ 的对角态射可由代数空间表示且局部有限型；
参见《形式空间》引理
\ref{formal-spaces-lemma-diagonal-morphism-formal-algebraic-spaces}
，而且我们已在引理 \ref{artin-lemma-rs} 的证明中看到 $W$ 具有 (RS)。}。
当然，这意味着给定任意 $n$，可以缩小 $V$，使 $V_n \to W_n$ 光滑。
问题在于找到一个同时对所有 $n$ 都有效的开集。

\medskip\noindent
问题在 $V$ 上是局部的，故可假设 $S = \Spec(R)$ 和 $V = \Spec(A)$ 都是仿射的。

\medskip\noindent
本段约化到 $W$ 为仿射形式代数空间的情形。选取仿射形式概形 $W'$ 以及
\'etale 态射 $W' \to W$，使 $v$ 在 $W_{red}$ 中的像属于
$W'_{red} \to W_{red}$ 的像。则
$V_{/Z} \times_{g, W} W' \to V_{/Z}$ 是 $S$ 上形式代数空间的 adic
\'etale 态射，而 $V_{/Z} \times_{g, W} W'$ 是仿射形式代数空间。
由《形式空间的代数化》引理
\ref{restricted-lemma-algebraize-rig-etale-affine}，存在仿射概形的
\'etale 态射 $\varphi : V' \to V$，使得 $V'$ 沿
$Z' = \varphi^{-1}(Z)$ 的完备化在 $V_{/Z}$ 上同构于
$V_{/Z} \times_{g, W} W'$。注意，$v$ 是某个 $v' \in V'$ 的像。
由于光滑性在基变换下保持，对所有 $n$，$V'_n \to W'_n$ 都光滑。
下一段将证明：用 $v'$ 的一个开邻域替换 $V'$ 后，对所有 $n$，
态射 $V'_n \to W'_n$ 都光滑。然后，用 $V' \to V$ 的开像替换 $V$，
由光滑性的 \'etale 下降可知 $V_n \to W_n$ 光滑。略去一些细节。

\medskip\noindent
假设 $S = \Spec(R)$、$V = \Spec(A)$、$Z = V(I)$ 且
$W = \text{Spf}(B)$。设 $v$ 对应于极大理想
$I \subset \mathfrak m \subset A$。给定 adic 连续 $R$-代数同态
$$
B \longrightarrow A^\wedge
$$
设 $\mathfrak b \subset B$ 是拓扑幂零元素所成的理想（这是 Noether adic
拓扑环 $B$ 的最大定义理想）。注意，$\mathfrak b A^\wedge$ 和
$IA^\wedge$ 都是 Noether adic 环 $A^\wedge$ 的定义理想。此外，
$\mathfrak m A^\wedge$ 是 $A^\wedge$ 的极大理想，包含
$\mathfrak b A^\wedge$ 和 $IA^\wedge$。已知
$$
B_n = B/\mathfrak b^n \to A^\wedge/\mathfrak b^n A^\wedge = A_n
$$
对所有 $n$ 都在 $\mathfrak m$ 处光滑。由上面的讨论，可以并且确实假设
$B_1 \to A_1$ 是光滑环同态。以 $\mathfrak m_1 \subset A_1$ 表示对应于
$\mathfrak m$ 的极大理想。
% P11-E0335: the frozen source uses the incomplete construction "Denote m_1 ... the maximal ideal".
由于光滑性蕴含平坦性，对所有 $n \geq 1$，映射
$$
\mathfrak b^n/\mathfrak b^{n + 1} \otimes_{B_1} (A_1)_{\mathfrak m_1}
\longrightarrow
\left(\mathfrak b^nA^\wedge/\mathfrak b^{n + 1}A^\wedge\right)_{\mathfrak m_1}
$$
是同构；参见《代数》引理 \ref{algebra-lemma-what-does-it-mean-again}。
考虑 Rees 代数
$$
B' = \bigoplus\nolimits_{n \geq 0} \mathfrak b^n/\mathfrak b^{n + 1}
$$
，它是 Noether 环 $B_1$ 上有限型的分次代数；再考虑 Rees 代数
$$
A' = \bigoplus\nolimits_{n \geq 0}
\mathfrak b^nA^\wedge/\mathfrak b^{n + 1}A^\wedge
$$
，它是 Noether 环 $A_1$ 上有限型的分次代数。
% P11-E0336: the frozen source duplicates the article in "a a finite type graded algebra".
考虑分次 $A_1$-代数同态
$$
\Psi : B' \otimes_{B_1} A_1 \longrightarrow A'
$$
由上文，此映射在 $A_1$ 的极大理想 $\mathfrak m_1$ 处局部化后是同构。
因此，$\Ker(\Psi)$（相应地，$\Coker(\Psi)$）是
$B' \otimes_{B_1} A_1$（相应地，$A'$）上的有限模，且其在
$\mathfrak m_1$ 处的局部化为零。从而，用一个主局部化替换 $A_1$
（并相应替换 $A$）后，可假设 $\Psi$ 是同构。（这是证明的关键步骤。）
反向推导可知 $B_n \to A_n$ 平坦；参见《代数》引理
\ref{algebra-lemma-what-does-it-mean-again}。因此 $A_n \to B_n$ 光滑
（作为具有光滑纤维的平坦环同态；参见《代数》引理
\ref{algebra-lemma-flat-fibre-smooth}），证明完成。
% P11-E0337: the frozen source reverses the ring-map direction in the final smoothness conclusion; the map proved flat and used throughout is B_n -> A_n.
\end{proof}

\begin{lemma}
\label{artin-lemma-openness-versality}
在情形 \ref{artin-situation-contractions} 中，函子 $F$ 满足通用性的开性。
\end{lemma}

\begin{proof}
必须证明如下事实：给定 $S$ 上局部有限型的概形 $V$、$\xi \in F(V)$，
以及有限型点 $v_0 \in V$，且 $\xi$ 在 $v_0$ 处通用，则用 $v_0$ 的一个
开邻域替换 $V$ 后，$\xi$ 在 $V$ 的每个有限型点处都通用。写作
$\xi = (Z, u', \hat x)$。

\medskip\noindent
第一种情形：$v_0 \not \in Z$。首先可用 $V \setminus Z$ 替换 $V$。
于是 $\xi = (\emptyset, u', \emptyset)$，且态射
$u' : V \to X'$ 在 $v_0$ 处通用。由《空间的态射续篇》引理
\ref{spaces-more-morphisms-lemma-lifting-along-artinian-at-point}
，这意味着 $u' : V \to X'$ 在 $v_0$ 处光滑。由于一个态射的光滑点集合是开集，
缩小 $V$ 后可假设 $u'$ 光滑。
% P11-E0338: the frozen source has the spurious article in "the set of a points".
同一引理于是说明 $\xi$ 在每一点处通用，正如所需。

\medskip\noindent
第二种情形：$v_0 \in Z$。像《形式空间》引理
\ref{formal-spaces-lemma-structure-locally-noetherian} 中那样写作
$W = \colim W_n$。由引理 \ref{artin-lemma-openness-smoothness}，可假设
$\hat x : V_{/Z} \to W$ 是形式代数空间的光滑态射。立即可知，
$\xi = (Z, u', \hat x)$ 在 $Z$ 的所有有限型点处通用。设
$V' \to V$、$\hat x'$ 和 $x'$ 见证 $u'$ 与 $\hat x$ 之间的相容性。
作为 $\hat x$ 的基变换，$\hat x' : V'_{/Z} \to X'_{/T'}$ 光滑。
由于 $\hat x'$ 是 $x' : V' \to X'$ 的完备化，由已经使用过的
《空间的态射续篇》引理
\ref{spaces-more-morphisms-lemma-lifting-along-artinian-at-point}.
可知，$x' : V' \to X'$ 在
$(V' \to V)^{-1}(Z) = |x'|^{-1}(T') \subset |V'|$ 的所有点处光滑。
由于一个态射的光滑点集合是开集，$x'$ 不光滑的点所成的闭集
$B \subset |V'|$ 与 $(V' \to V)^{-1}(Z)$ 不相交。由于
$V' \to V$ 为真态射，因而是闭态射，故
$(V' \to V)(B) \subset V$ 是与 $Z$ 不相交的闭子集。因此，缩小 $V$ 后
可假设 $B = \emptyset$，即 $x'$ 光滑。由上一段的讨论，这恰好意味着
$\xi$ 在 $V$ 中不属于 $Z$ 的所有有限型点处通用，证明完成。
\end{proof}

\noindent
下面给出最终结果。

\begin{theorem}
\label{artin-theorem-contractions}
\begin{reference}
\cite[Theorem 3.1]{ArtinII}
\end{reference}
设 $S$ 是局部 Noether 概形，且对所有有限型点 $s \in S$，
$\mathcal{O}_{S, s}$ 都是 G-环。设 $X'$ 是 $S$ 上局部有限型的代数空间，
$T' \subset |X'|$ 是闭子集。设 $W$ 是 $S$ 上局部 Noether 的形式代数空间，
且 $W_{red}$ 在 $S$ 上局部有限型。最后，设
$$
g : X'_{/T'} \longrightarrow W
$$
为形式修正；参见《形式空间的代数化》定义
\ref{restricted-definition-formal-modification}。若 $X'$ 和 $W$ 在 $S$
上分离\footnote{参见注 \ref{artin-remark-separated-needed}。}，则存在 $S$ 上
代数空间的真态射 $f : X' \to X$、闭子集 $T \subset |X|$，以及形式代数空间的
同构 $a : X_{/T} \to W$，使得
\begin{enumerate}
\item $T'$ 是 $T$ 在 $|f| : |X'| \to |X|$ 下的逆像；
\item $f : X' \to X$ 将 $X' \setminus T'$ 同构地映到 $X \setminus T$；以及
\item $g = a \circ f_{/T}$，其中
$f_{/T} : X'_{/T'} \to X_{/T}$ 是诱导态射。
\end{enumerate}
换言之，$(f : X' \to X, T, a)$ 是本节前面定义的一个解。
\end{theorem}

\begin{proof}
设 $F$ 是本节中利用 $X'$、$T'$、$W$、$g$ 构造的函子。由引理
\ref{artin-lemma-functor-is-solution}，只需证明 $F$ 对应于 $S$ 上局部有限型的
代数空间 $X$。为此，将应用命题
\ref{artin-proposition-spaces-diagonal-representable-noetherian}。具体而言，
由引理 \ref{artin-lemma-diagonal-contractions}，$F$ 的对角态射可由闭浸入表示；
又由诸引理 \ref{artin-lemma-sheaf}、\ref{artin-lemma-limit-preserving}、
\ref{artin-lemma-rs}, \ref{artin-lemma-finite-dim},
\ref{artin-lemma-formal-object-effective} 和 \ref{artin-lemma-openness-versality}，
我们得到公理 [0]、[1]、[2]、[3]、[4]、[5]。
\end{proof}

\begin{remark}
\label{artin-remark-separated-needed}
定理 \ref{artin-theorem-contractions} 的证明在两处使用了 $X'$ 和 $W$ 在 $S$
上分离这一事实。第一处是用它证明
$\Delta : F \to F \times F$ 可由代数空间表示。若把分离情形下的定理应用于
注 \ref{artin-remark-diagonal} 中的三元组
$E'$、$(E' \to V)^{-1}Z$、$E'_{/Z} \to E_W$，便可证明 $\Delta$ 可表示，
从而完全避免这一次使用该假设（这是对角态射通常的自举过程）。
% P11-E0339: the frozen source says "the triples" although E', the inverse image, and the formal modification are the three components of one triple.
因此，引理 \ref{artin-lemma-formal-object-effective} 的证明，是我们对定理
\ref{artin-theorem-contractions} 的证明中唯一真正需要使用 $X' \to S$ 分离之处。
读者可以检验，我们只用该假设来得到态射
$x' : V' \to X'$。若 $X' \to S$ 拟分离，则利用文献中的结果可以证明
$x'$ 存在；参见《空间的态射续篇》注
\ref{spaces-more-morphisms-remark-weaken-separation-axioms-question}.
由此可知，该定理在把
“分离”替换为“拟分离”后仍按原陈述成立。若以后需要这一版本，我们会在这里
精确陈述并仔细证明。
\end{remark}
